% =====================================================================
% Document: Virtualized NAC with Automated Wired 802.1X Incident Validation
% =====================================================================
% Author: Eric Schmitt
%
% PURPOSE
%   This document provides a formal engineering implementation and
%   validation record for a virtualized Network Access Control (NAC)
%   proof-of-concept lab using wired IEEE 802.1X authentication, policy-
%   based guest access enforcement, centralized logging, incident
%   tracking, automation, and packet capture analysis.
%
% Build Notes:
%   - Compiled via standard LaTeX toolchain (Pandoc/pdflatex).
%   - Hyperref used for internal cross-referencing and PDF metadata behavior.
%
% Commenting Philosophy:
%   - Comment blocks contextualize normative requirements, architectural
%     intent, change sensitivity, and verbatim operational evidence.
%   - Comment blocks may be intentionally omitted for self-contained or
%     mechanically descriptive subsections.
% =====================================================================

% -----------------------------------------------------------------------------
% Document Class
% -----------------------------------------------------------------------------
% Standard article class selected for portability and predictable layout.
% 11pt provides professional readability without excess pagination.
% -----------------------------------------------------------------------------
\documentclass[11pt]{article}

% -----------------------------------------------------------------------------
% Hyperref Configuration
% -----------------------------------------------------------------------------
% Enables:
%   - Internal section cross-referencing
%   - Clickable TOC entries
%   - Clickable figure, table, and appendix references
%   - Clean PDF open behavior (FitH view)
%
% No custom link coloring is applied in order to preserve a formal,
% documentation-oriented visual style.
% -----------------------------------------------------------------------------
\usepackage{hyperref}

% -----------------------------------------------------------------------------
% Graphicx Configuration
% -----------------------------------------------------------------------------
% Enables:
%   - Inclusion of externally rendered image assets
%   - Embedding of Mermaid-exported PNG diagrams
%   - Controlled figure scaling via \includegraphics
%
% This package is required for static diagram insertion in the main body,
% including the reference topology diagram and later figure assets.
% -----------------------------------------------------------------------------
\usepackage{graphicx}

% -----------------------------------------------------------------------------
% Float Configuration
% -----------------------------------------------------------------------------
% Enables:
%   - Use of stricter figure and table placement specifiers such as [H]
%   - More predictable handling of topology diagrams, screenshots, and tables
%   - Reduced float drift in long technical sections
%
% This package is included to improve placement control when the default
% LaTeX float algorithm is too permissive for evidence-heavy documents.
% -----------------------------------------------------------------------------
\usepackage{float}

% -----------------------------------------------------------------------------
% Tabularx Configuration
% -----------------------------------------------------------------------------
% Enables:
%   - Automatically width-constrained tables that fit within \textwidth
%   - Flexible wrapping for long prose-heavy columns in PDF output
%   - More reliable rendering than fixed p{} widths for mixed-content tables
%
% This package is preferred for metadata tables, technical matrices, and
% other wide documentation tables that must render consistently in both
% HTML and PDF outputs under the Pandoc workflow.
% -----------------------------------------------------------------------------
\usepackage{tabularx}

\hypersetup{
  pdfstartpage=1,
  pdfstartview=FitH,
}

\title{Virtualized NAC with Automated Wired 802.1X Incident Validation}
\subtitle{Version 1.0}
\author{© 2026 Eric Schmitt. All rights reserved.}

\begin{document}
\maketitle
\thispagestyle{empty}
\newpage

% =====================================================================
% SECTION 1: DOCUMENT CONTROL
% =====================================================================
% This section establishes document identity, governance, revision
% history, intended audience, and usage constraints. It mirrors the
% structure and tone of the author's prior AWS engineering lab document
% and is written as if this were an internal implementation and
% validation artifact intended for technical review.
% =====================================================================

\section{1. Document Control}
\label{sec:document-control}

% ---------------------------------------------------------------------
% SUBSECTION 1.1: DOCUMENT METADATA
% ---------------------------------------------------------------------
% This subsection records the authoritative metadata for the document,
% including title, version, status, scope, and platform context.
% ---------------------------------------------------------------------

\subsection{1.1 Document Metadata}
\label{sec:document-metadata}

\begin{center}
\small
\setlength{\tabcolsep}{4pt}
\renewcommand{\arraystretch}{1.15}
\begin{tabularx}{0.92\textwidth}{|p{0.24\textwidth}|X|}
\hline
\textbf{Field} & \textbf{Value} \\
\hline
Document Title & Virtualized NAC with Automated Wired 802.1X Incident Validation \\
\hline
Author & Eric Schmitt \\
\hline
Document Version & 1.0 \\
\hline
Date of Issue & March 24, 2026 \\
\hline
Document Status & Production Draft (Educational Lab / Proof-of-Concept) \\
\hline
Document Type & Engineering Lab Implementation and Validation Report \\
\hline
Primary Scope & Virtualized wired IEEE 802.1X NAC proof-of-concept with policy-based guest \\
& access enforcement, incident simulation, automation, and packet-level \\
& validation \\
\hline
Primary Technologies & IEEE 802.1X via \texttt{hostapd} (wired mode) and \texttt{wpa\_supplicant}; pfSense; \\
& FreeRADIUS; GLPI on Apache and MariaDB; Splunk Universal Forwarder; \\
& Python; Ansible; \texttt{tcpdump}; Wireshark \\
\hline
VM Platform & VMware Workstation Pro: Ubuntu Server 24.04 LTS \\
\hline
Primary Lab Purpose & Demonstrate a realistic small-scale NAC architecture with authenticated access \\
& control, constrained guest access, observability, ITSM workflow integration, \\
& automation, and protocol validation \\
\hline
\end{tabularx}
\end{center}

The purpose of this document is to provide a complete, technically rigorous, and reproducible record of the final implemented state of the lab environment, the access-control model selected for the environment, the operational validation performed against that environment, and the incident simulation used to demonstrate practical NAC troubleshooting and remediation workflows.

This document is intentionally written as an engineering implementation and validation report rather than as a tutorial, learning journal, or informal case study. Where the lab departs from production architecture due to virtualization or platform constraints, those departures are stated explicitly and evaluated in context rather than implied or omitted.

% ---------------------------------------------------------------------
% SUBSECTION 1.2: CHANGE HISTORY
% ---------------------------------------------------------------------
% This subsection records formal revision history.
% ---------------------------------------------------------------------

\subsection{1.2 Change History}
\label{sec:change-history}

\begin{center}
\begin{tabular}{|p{0.14\textwidth}|p{0.14\textwidth}|p{0.20\textwidth}|p{0.42\textwidth}|}
\hline
\textbf{Version} & \textbf{Date} & \textbf{Author} & \textbf{Description} \\
\hline
1.0 & 2026-03-24 & Eric Schmitt & Initial document draft; text-only. \\
\hline
\end{tabular}
\end{center}

Future revisions should preserve the same evidence-handling model used in this document: final-state operational outputs, log excerpts, complete configuration artifacts where relevant, and untruncated raw data in verbatim blocks for all authoritative command or file-based evidence included in scope.

% ---------------------------------------------------------------------
% SUBSECTION 1.3: INTENDED AUDIENCE AND USAGE
% ---------------------------------------------------------------------
% This subsection defines who the document is for, how it should be read,
% and what interpretive assumptions apply when reviewing the lab.
% ---------------------------------------------------------------------

\subsection{1.3 Intended Audience and Usage}
\label{sec:intended-audience-and-usage}

This document is intended for the following audiences:

\begin{itemize}
    \item Network engineering practitioners evaluating the design and validation of a virtualized NAC proof-of-concept;
    \item Security engineering practitioners reviewing the use of wired IEEE 802.1X, RADIUS-backed authentication, and segmented guest policy enforcement;
    \item Systems engineering practitioners assessing service integration across routing, firewalling, authentication, logging, ITSM, and automation domains;
    \item Technical reviewers evaluating the project as a portfolio artifact intended to demonstrate enterprise-relevant design and operational execution rather than isolated tool familiarity.
\end{itemize}

This document should be interpreted as a formal record of the final implemented state of the environment and the authoritative incident validation path, not as a chronological build log. Intermediate troubleshooting states, temporary misconfigurations, and transient construction artifacts that were used only to arrive at the final design are intentionally excluded unless they are directly relevant to the incident narrative or to explaining an architectural decision.

The operational evidence included in this document is constrained by the following rules:

\begin{itemize}
    \item Final-state command outputs, configuration files, service states, validation results, and automation artifacts are considered in scope;
    \item Packet capture interpretation and protocol-state observations are considered in scope;
    \item Full debug logs are not reproduced wholesale; only explicitly selected excerpts are included where necessary to support an incident conclusion;
    \item GUI-derived evidence (e.g., pfSense, GLPI, Wireshark) is represented by screenshots rather than terminal transcription unless a manual transcription is operationally clearer;
    \item Sensitive values, including passwords, shared secrets, tokens, and authentication material, MUST be redacted in all published forms of the document.
\end{itemize}

The environment documented herein is intentionally designed as a constrained but defensible approximation of enterprise NAC behavior within a virtualization platform. In particular, the guest-access model implemented in this lab is a policy-based guest access simulation enforced through Linux and pfSense control points rather than true switch-driven dynamic VLAN reassignment. This distinction is central to the technical integrity of the document and will be treated explicitly in subsequent sections.

% =====================================================================
% SECTION 2: EXECUTIVE SUMMARY
% =====================================================================
% This section provides a concise but formal summary of the lab's
% purpose, final outcome, major design decisions, and operational value.
% It is intentionally written in enterprise engineering language rather
% than in tutorial or retrospective style.
% =====================================================================

\section{2. Executive Summary}
\label{sec:executive-summary}

% ---------------------------------------------------------------------
% SUBSECTION 2.1: OBJECTIVE
% ---------------------------------------------------------------------
% This subsection states the primary objective of the lab in formal
% engineering terms.
% ---------------------------------------------------------------------

\subsection{2.1 Objective}
\label{sec:objective}

The objective of this implementation was to design, build, and validate a small-scale but enterprise-relevant Network Access Control (NAC) proof-of-concept within a fully virtualized lab environment. The target outcome was not merely to demonstrate isolated component installation, but to establish an integrated control path in which a wired endpoint could be authenticated via IEEE 802.1X against a RADIUS policy authority, granted or denied access based on authentication state, observed through centralized logging and evidence collection, and incorporated into a realistic incident-response workflow.

The implemented design intentionally spans multiple operational domains in order to reflect the way NAC is encountered in practice. These domains include endpoint authentication, routed segmentation, firewall policy enforcement, service observability, IT service management (ITSM) workflow, automation-assisted validation, and packet-level protocol analysis. The lab was therefore scoped to produce not only a functional access-control environment, but also a defensible record of how that environment behaves under both normal and failure conditions.

The specific technical objectives of the lab were as follows:

\begin{itemize}
    \item Implement a wired IEEE 802.1X authentication path using a Linux-based authenticator simulation and a FreeRADIUS policy server;
    \item Establish authenticated access to routed service resources for a validated endpoint;
    \item Implement a constrained guest-access model for unauthenticated or otherwise non-trusted endpoints using policy-based enforcement at the firewall boundary;
    \item Integrate local observability through service logs, structured incident evidence files, and Splunk Universal Forwarder monitoring inputs;
    \item Integrate ITSM workflow through GLPI by recording the incident as an actual tracked ticket associated with a managed asset;
    \item Implement meaningful automation through a Python evidence parser and an Ansible validation/remediation workflow;
    \item Validate both failure and recovery states using terminal evidence, service logs, and packet captures reviewed in Wireshark.
\end{itemize}

The result is intended to function as a technically rigorous proof-of-concept demonstrating that a virtualized environment can be used to model core NAC control flows and operational procedures, provided that platform limitations are explicitly acknowledged and the chosen enforcement model is documented accurately.

% ---------------------------------------------------------------------
% SUBSECTION 2.2: OUTCOME SUMMARY
% ---------------------------------------------------------------------
% This subsection summarizes the final implemented outcome and the major
% validation conclusions at a high level.
% ---------------------------------------------------------------------

\subsection{2.2 Outcome Summary}
\label{sec:outcome-summary}

The lab successfully achieved its core design and validation objectives.

A functioning wired IEEE 802.1X authentication path was implemented using a Linux-based access node acting as a wired authenticator, a Linux endpoint acting as the supplicant, and FreeRADIUS as the central authentication and policy service. The authenticated endpoint path was validated through successful RADIUS test transactions, successful supplicant authorization during controlled test execution, and confirmed routed reachability to service-network resources once the authorization path was operating as intended.

Because the environment was built entirely from virtual machines and not from managed enterprise switching hardware, the design could not support true switch-driven dynamic VLAN reassignment in the same manner as a production NAC deployment. Rather than obscuring this limitation, the implementation explicitly adopted a policy-based guest access simulation. In this model, unauthenticated or non-trusted endpoint behavior was represented by placement in a dedicated guest network segment behind pfSense, where DHCP, gateway access, and tightly constrained reachability were permitted while access to protected service resources was denied. This design choice preserved the operational intent of differentiated access control while remaining technically honest about the constraints of the platform.

The lab further succeeded in demonstrating an operational incident lifecycle rather than only a steady-state build. A wired 802.1X authentication failure was deliberately introduced in a controlled manner, evidence of the failed state was captured across the supplicant, authenticator, and RADIUS service layers, and the resulting condition was documented as a real incident in GLPI. Remediation was then performed by restoring the known-good client configuration and executing a clean re-test, after which successful authorization indicators were observed and the incident was closed as resolved.

Automation was incorporated in a manner that is materially relevant to the environment rather than ornamental. A Python script was developed to parse and summarize incident evidence artifacts, providing a deterministic summary of failed and remediated states. An Ansible playbook was developed to perform repeatable validation of NAC-related services and to support a bounded, operator-invoked remediation mode. These automation artifacts reinforce the document's central claim that the lab demonstrates not only component familiarity, but also repeatable operational procedure.

Finally, packet capture analysis was added as a first-class validation layer. Using \texttt{tcpdump} on the authenticator host and subsequent inspection in Wireshark, the lab captured both failed and successful authentication flows across the EAPOL and RADIUS control paths. This provided protocol-level corroboration of the incident narrative and strengthened the evidentiary value of the environment as a portfolio artifact.

% ---------------------------------------------------------------------
% SUBSECTION 2.3: ARCHITECTURAL CONSTRAINTS AND DESIGN DECISION
% ---------------------------------------------------------------------
% This subsection states the key architectural constraint and the
% resulting design decision in clear terms so that the rest of the
% document can rely on that framing without ambiguity.
% ---------------------------------------------------------------------

\subsection{2.3 Architectural Constraints and Design Decision}
\label{sec:architectural-constraints-and-design-decision}

The most important architectural constraint in this lab is that the environment does not include a managed access switch capable of acting as a production-grade IEEE 802.1X authenticator with RADIUS-driven dynamic VLAN reassignment, downloadable ACLs, or hardware-enforced post-auth policy state. Instead, the environment uses Ubuntu Server virtual machines and VMware virtual networks to model the control plane and segmentation boundaries.

This constraint has direct implications. While Linux user-space tools such as \texttt{hostapd} and \texttt{wpa\_supplicant} can accurately model important aspects of wired IEEE 802.1X authentication and EAP state transitions, they do not convert a general-purpose virtual machine into a fully equivalent enterprise access switch. As a result, the environment can credibly demonstrate:

\begin{itemize}
    \item supplicant-to-authenticator EAPOL exchanges;
    \item authenticator-to-RADIUS request handling;
    \item successful and failed authentication states;
    \item policy differentiation between trusted and guest access paths;
    \item operational validation, logging, and incident response workflows.
\end{itemize}

However, it cannot credibly claim to implement:

\begin{itemize}
    \item true dynamic VLAN reassignment on a switchport as the direct result of RADIUS policy;
    \item hardware-native port reauthorization behavior;
    \item switch-specific enforcement semantics such as vendor ACL or downloadable ACL features.
\end{itemize}

For that reason, the guest-access component of this lab is explicitly defined as a policy-based guest access simulation rather than as a true dynamic guest VLAN implementation. The enforcement boundary is achieved through Linux and pfSense control points: the access path models authentication state, while pfSense enforces the differentiated reachability policy for the guest segment. This is not mutually exclusive with the broader segmentation model used throughout the lab; rather, it is the most defensible way to preserve the operational objective of differentiated access within the constraints of a virtualized platform.

This design decision is central to the technical integrity of the document. It allows the environment to remain enterprise-relevant without overstating equivalence to production hardware behavior, and it ensures that the incident and automation workflows presented later in the document are evaluated against the architecture that was actually built, not against an idealized architecture that was not.

% =====================================================================
% SECTION 3: SCOPE, OBJECTIVES, AND DESIGN ASSUMPTIONS
% =====================================================================
% This section formally defines what is and is not included in the lab,
% the assumptions under which the design should be interpreted, and the
% criteria by which the implementation is considered successful.
% =====================================================================

\section{3. Scope, Objectives, and Design Assumptions}
\label{sec:scope-objectives-and-design-assumptions}

% ---------------------------------------------------------------------
% SUBSECTION 3.1: IN-SCOPE COMPONENTS
% ---------------------------------------------------------------------
% This subsection enumerates the technologies and functional domains that
% are intentionally included in the implementation.
% ---------------------------------------------------------------------

\subsection{3.1 In-Scope Components}
\label{sec:in-scope-components}

The following components and functional domains are in scope for this implementation and are treated as authoritative elements of the final environment:

\begin{itemize}
    \item \textbf{Virtualized routing and segmentation:} multi-segment VMware virtual networking, Linux routing on intermediate nodes, and pfSense routing/firewall policy between guest and service segments;
    \item \textbf{Wired IEEE 802.1X authentication path:} a Linux-based wired authenticator simulation using \texttt{hostapd}, a Linux supplicant using \texttt{wpa\_supplicant}, and FreeRADIUS as the authentication and policy service;
    \item \textbf{RADIUS policy service:} FreeRADIUS user and client definitions, service validation, local and remote authentication testing, and runtime logging;
    \item \textbf{Guest access control model:} policy-based guest access simulation enforced by pfSense rather than true dynamic VLAN reassignment;
    \item \textbf{Observability and evidence collection:} service logs, journald evidence, explicit incident capture files, and packet capture collection with \texttt{tcpdump};
    \item \textbf{Log monitoring instrumentation:} Splunk Universal Forwarder configured to monitor relevant system and FreeRADIUS log sources;
    \item \textbf{ITSM workflow integration:} GLPI deployment, managed asset creation, incident ticket creation, incident update, and incident closure;
    \item \textbf{Automation artifacts:} Python-based incident evidence parser and Ansible-based validation/remediation workflow;
    \item \textbf{Packet-level protocol validation:} failed and successful EAPOL and RADIUS captures reviewed in Wireshark and interpreted as part of the incident narrative.
\end{itemize}

The lab is therefore intentionally broader than a single-service build. It is designed to demonstrate cross-domain operational competence across networking, security, Linux services, observability, automation, and incident handling.

% ---------------------------------------------------------------------
% SUBSECTION 3.2: OUT-OF-SCOPE COMPONENTS
% ---------------------------------------------------------------------
% This subsection states what the implementation does not claim to be and
% what production-grade features are intentionally excluded.
% ---------------------------------------------------------------------

\subsection{3.2 Out-of-Scope Components}
\label{sec:out-of-scope-components}

The following capabilities, platforms, or production-grade behaviors are explicitly out of scope for this implementation:

\begin{itemize}
    \item \textbf{Managed enterprise switching hardware} with native switchport control semantics, hardware-enforced post-auth state, or vendor-specific NAC integration features;
    \item \textbf{True dynamic VLAN reassignment} directly applied to an access port by RADIUS attributes during live port authorization;
    \item \textbf{Downloadable ACLs (dACLs)} or equivalent vendor-specific post-auth policy enforcement constructs;
    \item \textbf{Certificate-based enterprise EAP methods} such as EAP-TLS; the lab uses a simpler method appropriate for deterministic proof-of-concept validation rather than a full PKI-backed endpoint identity model;
    \item \textbf{Full Splunk Enterprise deployment} including indexer, search head, dashboards, or distributed search architecture; the lab intentionally uses Splunk Universal Forwarder only;
    \item \textbf{High availability or clustering} for core services such as RADIUS, firewalling, or ITSM;
    \item \textbf{Production secrets management} such as Vault-backed secret retrieval, certificate lifecycle automation, or centrally governed credential rotation;
    \item \textbf{Full configuration management pipeline integration} such as CI/CD, GitOps enforcement, or scheduled compliance scans beyond the bounded automation included here.
\end{itemize}

These exclusions do not weaken the document; rather, they define the limits within which the implementation should be evaluated. The lab is intended to be a credible and technically honest proof-of-concept, not a claim of feature parity with a production NAC deployment.

% ---------------------------------------------------------------------
% SUBSECTION 3.3: DESIGN ASSUMPTIONS
% ---------------------------------------------------------------------
% This subsection captures the assumptions under which the architecture
% and validation logic should be interpreted.
% ---------------------------------------------------------------------

\subsection{3.3 Design Assumptions}
\label{sec:design-assumptions}

The architecture and validation logic in this document rely on the following design assumptions:

\begin{itemize}
    \item VMware virtual networks are treated as deterministic Layer 2 broadcast domains suitable for representing isolated access, transit, service, and guest segments;
    \item Ubuntu Server virtual machines can credibly model routing, firewall-adjacent policy boundaries, service hosting, and user-space authenticator/supplicant behavior, but not the full control semantics of a managed enterprise switch;
    \item The Linux-based access node provides a valid proof-of-concept wired authenticator role for demonstrating EAPOL state transitions and RADIUS interaction, even though it is not a substitute for hardware switch behavior;
    \item pfSense provides the authoritative guest-policy enforcement boundary in the final design, including guest subnet gatewaying, DHCP, and restrictive firewall policy;
    \item Temporary host-side management IP assignments on VMware virtual network interfaces are acceptable as an operational convenience for accessing isolated management surfaces during build and validation, provided those assignments are clearly described and not misrepresented as part of the guest or authenticated endpoint data path;
    \item Final-state validation evidence is more relevant to the purpose of this document than full chronological build history, and therefore only final-state and incident-relevant evidence is included;
    \item Packet capture interpretation is treated as corroborating evidence for the control-plane narrative rather than as the sole source of truth; service state and incident logs remain authoritative in combination with capture findings.
\end{itemize}

These assumptions are intentionally conservative. They preserve the lab's credibility by clearly separating what is being simulated from what is being directly implemented.

% ---------------------------------------------------------------------
% SUBSECTION 3.4: SUCCESS CRITERIA
% ---------------------------------------------------------------------
% This subsection defines the objective conditions under which the lab is
% considered successful and portfolio-worthy.
% ---------------------------------------------------------------------

\subsection{3.4 Success Criteria}
\label{sec:success-criteria}

The implementation is considered successful only if all of the following criteria are satisfied:

\begin{itemize}
    \item A wired endpoint can successfully complete the intended 802.1X authentication flow through the Linux-based authenticator and FreeRADIUS service under controlled test conditions;
    \item The authenticated path provides the expected routed reachability to protected service-network resources in the final-state topology;
    \item A guest or otherwise non-trusted endpoint path can be demonstrated as operationally distinct, including guest DHCP and restricted reachability enforced by pfSense policy;
    \item FreeRADIUS can be validated both locally and remotely using deterministic test transactions and final-state service health checks;
    \item Splunk Universal Forwarder is configured to monitor the intended log sources relevant to NAC and incident evidence collection;
    \item GLPI is operational and used to represent a real incident record associated with a managed asset, rather than being present only as an unused supporting service;
    \item A controlled wired 802.1X authentication failure can be reproduced, observed, and documented across the supplicant, authenticator, and RADIUS layers;
    \item The failed condition can be remediated, and a clean re-test can confirm restoration of successful authorization;
    \item The incident can be summarized through automation using the Python evidence parser and validated through repeatable service checks using the Ansible workflow;
    \item Packet captures of both failed and successful authentication flows can be collected and interpreted in a way that supports the incident narrative and final validation conclusions.
\end{itemize}

The standard applied throughout this document is whether the environment can support a coherent access-control narrative, a reproducible incident, and a technically defensible remediation and validation process.

% =====================================================================
% SECTION 4: REFERENCE ARCHITECTURE
% =====================================================================
% This section defines the authoritative logical architecture of the lab
% in its final implemented state. It is intentionally written as a clean
% architectural description rather than as a chronological build record.
% Supporting raw configuration artifacts and final-state command outputs
% are referenced forward to the appendices where applicable.
% =====================================================================

\section{4. Reference Architecture}
\label{sec:reference-architecture}

% ---------------------------------------------------------------------
% SUBSECTION 4.1: LOGICAL TOPOLOGY OVERVIEW
% ---------------------------------------------------------------------
% This subsection describes the major systems, trust boundaries, and
% traffic paths in the final architecture.
% ---------------------------------------------------------------------

\subsection{4.1 Logical Topology Overview}
\label{sec:logical-topology-overview}

The final lab environment is implemented as a multi-segment virtualized network hosted on VMware Workstation Pro running on an Ubuntu 24.04 workstation. The design intentionally separates management, authenticated access, routed transit, service hosting, and guest access into distinct logical domains in order to model the control boundaries that would exist in a small enterprise environment.

At a high level, the architecture consists of the following functional layers:

\begin{itemize}
    \item \textbf{Endpoint access layer:} a Linux endpoint (\texttt{client-1}) acting as the wired supplicant and a Linux access node (\texttt{access-node}) acting as the wired IEEE 802.1X authenticator using \texttt{hostapd};
    \item \textbf{Routed access transit:} a Linux distribution router (\texttt{dist-router}) providing the intermediate routed path between the access domain and the firewall-controlled service domain;
    \item \textbf{Security and policy boundary:} a pfSense virtual firewall providing routing, interface separation, guest DHCP, and guest-policy enforcement between protected services and guest access;
    \item \textbf{Protected service domain:} a services segment containing \texttt{radius-splunk} (FreeRADIUS, Splunk Universal Forwarder, Python/Ansible automation) and \texttt{ops-glpi} (GLPI, Apache, MariaDB);
    \item \textbf{Guest access domain:} a separate pfSense-controlled guest network used to simulate unauthenticated or non-trusted endpoint behavior through policy-based access restriction rather than dynamic switch-driven VLAN reassignment;
    \item \textbf{Out-of-band management paths:} VMware-provided host-side virtual network access used where necessary for administrative reachability to otherwise isolated interfaces during build and validation.
\end{itemize}

The most important architectural distinction in this environment is the separation between authentication state and guest-policy enforcement state. Authentication is modeled directly in the Linux-based 802.1X control path between the supplicant, the authenticator, and FreeRADIUS. Guest-policy enforcement is modeled at the pfSense boundary, where a dedicated guest subnet receives DHCP service and is restricted from reaching protected service resources. This division is intentional and reflects the platform constraint that the lab cannot perform true switch-native dynamic VLAN reassignment.

In practical terms, the environment therefore demonstrates the following control-plane relationships:

\begin{itemize}
    \item \texttt{client-1} can authenticate through \texttt{access-node} to FreeRADIUS when operating in the wired 802.1X validation path;
    \item \texttt{access-node} forwards RADIUS authentication traffic toward the protected services segment via \texttt{dist-router} and pfSense routing;
    \item \texttt{radius-splunk} provides the policy authority for the authentication workflow and also serves as the automation control node;
    \item \texttt{ops-glpi} provides the ITSM platform used to represent the incident as an actual tracked operational record;
    \item pfSense enforces the guest access model by providing DHCP and tightly scoped reachability on the guest segment while preventing direct access from the guest segment to protected services.
\end{itemize}

This architecture is sufficiently complex to demonstrate realistic operational dependencies without becoming so large that the evidentiary chain becomes difficult to follow. The final design is therefore intentionally narrow in breadth but deep in validation.

% ---------------------------------------------------------------------
% SUBSECTION 4.2: NETWORK SEGMENTATION AND ADDRESSING
% ---------------------------------------------------------------------
% This subsection defines the authoritative subnet and interface model
% used throughout the document. Supporting final-state command outputs
% and configuration files will be referenced in appendices.
% ---------------------------------------------------------------------

\subsection{4.2 Network Segmentation and Addressing}
\label{sec:network-segmentation-and-addressing}

The final environment uses multiple RFC1918 segments to model distinct functional domains. These segments are intentionally small and purpose-driven. The authoritative logical addressing model used throughout this document is summarized below.

\textbf{Logical Network Segments and Addressing}

\medskip

\textbf{Access Client Segment} \\
\textbf{Subnet:} \texttt{10.0.12.0/24} \\
\textbf{Primary Function:} Authenticated endpoint attachment segment for the wired supplicant validation path. \\
\textbf{Representative Systems / Interfaces:} \texttt{client-1} (\texttt{10.0.12.10}), \texttt{access-node} access-facing interface (\texttt{10.0.12.1}).

\medskip

\textbf{Access-to-Distribution Transit} \\
\textbf{Subnet:} \texttt{10.0.23.0/24} \\
\textbf{Primary Function:} Routed transit between the access node and the intermediate distribution router. \\
\textbf{Representative Systems / Interfaces:} \texttt{access-node} (\texttt{10.0.23.1}), \texttt{dist-router} (\texttt{10.0.23.2}).

\medskip

\textbf{Protected Services Segment} \\
\textbf{Subnet:} \texttt{10.0.30.0/24} \\
\textbf{Primary Function:} Service hosting for RADIUS, logging, automation control, and ITSM systems. \\
\textbf{Representative Systems / Interfaces:} pfSense OPT1 (\texttt{10.0.30.1}), \texttt{radius-splunk} (\texttt{10.0.30.10}), \texttt{ops-glpi} (\texttt{10.0.30.20}).

\medskip

\textbf{Distribution-to-Firewall Transit} \\
\textbf{Subnet:} \texttt{10.0.34.0/30} \\
\textbf{Primary Function:} Point-to-point transit between the Linux distribution router and pfSense. \\
\textbf{Representative Systems / Interfaces:} \texttt{dist-router} (\texttt{10.0.34.1}), pfSense LAN/transit (\texttt{10.0.34.2}).

\medskip

\textbf{Guest Access Segment} \\
\textbf{Subnet:} \texttt{10.0.20.0/24} \\
\textbf{Primary Function:} Guest or non-trusted endpoint segment used for policy-based guest access simulation. \\
\textbf{Representative Systems / Interfaces:} pfSense GUEST/OPT2 (\texttt{10.0.20.254}), DHCP pool (\texttt{10.0.20.100--10.0.20.199}).

\medskip

\textbf{Host / Hypervisor Management Segments} \\
\textbf{Subnet:} VMware-provided RFC1918 segments (e.g., \texttt{192.168.195.0/24}, \texttt{192.168.205.0/24}). \\
\textbf{Primary Function:} Administrative reachability and host-level VM management paths. \\
\textbf{Representative Systems / Interfaces:} Host virtual interfaces, VM management adapters, temporary host-side management access where required.

\medskip

The protected services segment (\texttt{10.0.30.0/24}) is the most operationally significant network in the final design. It hosts the RADIUS policy authority, the log-forwarding instrumentation, the Python and Ansible automation artifacts, and the ITSM platform used to record the incident lifecycle. This segment is reachable from the authenticated access path through the routed chain \texttt{access-node} $\rightarrow$ \texttt{dist-router} $\rightarrow$ pfSense, and it is intentionally protected from direct guest-segment access by pfSense firewall policy.

The guest access segment (\texttt{10.0.20.0/24}) is not intended to be interpreted as a dynamically assigned VLAN in the switch-native sense. Rather, it is the authoritative guest-policy enforcement domain used by this lab to represent how an unauthenticated or otherwise restricted endpoint would be treated under a constrained virtualized implementation. The guest segment receives DHCP from pfSense and can reach the guest gateway for local validation, but it is intentionally restricted from reaching the protected services segment as part of the policy model described later in Section~\ref{sec:access-control-model-and-policy-enforcement}.

The authoritative final-state interface configurations and route outputs that support this addressing model will be captured later in the document and in the appendices. Where exact file-based network definitions are relevant, they will be cross-referenced directly (for example, see Section~\ref{sec:appendix-a-radius-splunk-netplan-config}, Section~\ref{sec:appendix-a-ops-glpi-netplan-config}, Section~\ref{sec:appendix-a-access-node-netplan-config}, Section~\ref{sec:appendix-a-client-1-netplan-config}, and Section~\ref{sec:platform-build-and-final-state-configuration}).

% ---------------------------------------------------------------------
% SUBSECTION 4.3: FINAL SYSTEM ROLES
% ---------------------------------------------------------------------
% This subsection defines the authoritative role of each major node in
% the final design and creates forward references for appendices.
% ---------------------------------------------------------------------

\subsection{4.3 Final System Roles}
\label{sec:final-system-roles}

Each system in the final lab has a deliberately constrained but clearly defined role. This section establishes those roles so that later configuration, validation, and incident sections can refer to each node without ambiguity.

\subsubsection{4.3.1 access-node}
\label{sec:access-node-role}

\texttt{access-node} is the primary access-layer control point in the lab. It serves as the Linux-based wired IEEE 802.1X authenticator using \texttt{hostapd} in wired mode and also functions as the first routed hop away from the endpoint access segment. Its responsibilities include:

\begin{itemize}
    \item providing the access-facing interface for the wired supplicant path;
    \item originating EAPOL interaction with the supplicant during controlled 802.1X testing;
    \item forwarding RADIUS authentication requests toward the FreeRADIUS service;
    \item maintaining static routing toward the protected services domain via \texttt{dist-router};
    \item serving as the packet capture location for both EAPOL and RADIUS control-plane captures during incident validation.
\end{itemize}

Its final-state configuration artifacts will be documented later (see Section~\ref{sec:appendix-a-access-node-netplan-config}, Section~\ref{sec:appendix-a-access-node-hostapd-wired-config}, Section~\ref{sec:appendix-a-access-node-hostapd-wired-systemd-unit}, and Section~\ref{sec:platform-build-and-final-state-configuration}).

\subsubsection{4.3.2 client-1}
\label{sec:client-1-role}

\texttt{client-1} is the endpoint used to represent a wired supplicant. It participates in both the authenticated access validation path and the guest-policy simulation path, depending on the test being executed. Its responsibilities include:

\begin{itemize}
    \item acting as the 802.1X supplicant using \texttt{wpa\_supplicant} in wired mode;
    \item serving as the primary subject of the simulated authentication failure and remediation workflow;
    \item demonstrating guest-segment behavior when attached to the pfSense-controlled guest network;
    \item providing endpoint-perspective evidence for successful and failed authorization states.
\end{itemize}

Its final-state network configuration and wired supplicant configuration will be documented later (see Section~\ref{sec:appendix-a-client-1-netplan-config}, Section~\ref{sec:appendix-a-client-1-wpa-supplicant-wired-8021x-config}, and Section~\ref{sec:appendix-b-incident-evidence-files}).

\subsubsection{4.3.3 dist-router}
\label{sec:dist-router-role}

\texttt{dist-router} is the intermediate Linux routing node that provides the transit path between the access-layer control domain and the pfSense-protected service domain. It exists to make the architecture more representative of a routed enterprise topology rather than collapsing all functions into a single flat segment. Its responsibilities include:

\begin{itemize}
    \item routing traffic between the access-node transit segment and the pfSense transit interface;
    \item providing the next hop from the access domain toward the protected services segment;
    \item supporting final-state path validation and demonstrating that the RADIUS control path traverses a realistic routed intermediate node.
\end{itemize}

Its final-state interface and route evidence will be included later as part of the validation and appendix materials (see Section~\ref{sec:platform-build-and-final-state-configuration}).

\subsubsection{4.3.4 pfSense}
\label{sec:pfsense-role}

pfSense is the authoritative security boundary and policy-enforcement node in the final design. It separates the protected services domain from the guest domain and provides the guest gateway, guest DHCP service, and restrictive guest firewall policy. Its responsibilities include:

\begin{itemize}
    \item terminating the routed transit from \texttt{dist-router};
    \item providing the protected services segment gateway;
    \item providing the guest segment gateway;
    \item issuing DHCP leases on the guest segment;
    \item enforcing the guest access policy that differentiates guest reachability from authenticated service reachability;
    \item serving as the primary firewall boundary used to simulate guest access control in the absence of true switch-native dynamic VLAN reassignment.
\end{itemize}

Its final interface assignments, DHCP configuration, and guest firewall rules will be documented later (see Section~\ref{sec:appendix-a-pfsense-authenticated-interface-summary}, Section~\ref{sec:appendix-a-pfsense-guest-interface-summary}, Section~\ref{sec:appendix-a-pfsense-guest-firewall-rules}, and the related screenshots in the main body).

\subsubsection{4.3.5 radius-splunk}
\label{sec:radius-splunk-role}

\texttt{radius-splunk} is the most functionally dense system in the lab and serves as both the protected authentication authority and the automation/observability control node. It hosts:

\begin{itemize}
    \item FreeRADIUS as the central authentication and policy service;
    \item Splunk Universal Forwarder configured to monitor relevant local log sources;
    \item the Python incident evidence parser used to summarize failed and remediated states;
    \item the Ansible control environment used for repeatable validation and bounded remediation workflows.
\end{itemize}

Because this node combines policy authority, observability, and automation, it is central to both the steady-state validation path and the incident-response narrative. Its configuration artifacts will be referenced extensively later (see Section~\ref{sec:appendix-a-radius-splunk-netplan-config}, Section~\ref{sec:appendix-a-radius-splunk-freeradius-users-config}, Section~\ref{sec:appendix-a-radius-splunk-freeradius-clients-config}, Section~\ref{sec:appendix-a-splunk-uf-inputs-conf}, Section~\ref{sec:appendix-a-python-nac-incident-parser}, and Section~\ref{sec:appendix-a-ansible-nac-validation-playbook}).

\subsubsection{4.3.6 ops-glpi}
\label{sec:ops-glpi-role}

\texttt{ops-glpi} provides the IT service management layer for the lab. It is intentionally not treated as a passive supporting service; instead, it is used to represent the incident as an actual tracked operational record tied to a managed asset. Its responsibilities include:

\begin{itemize}
    \item hosting GLPI on Apache with MariaDB backend services;
    \item maintaining the \texttt{client-1} asset record used in the incident workflow;
    \item maintaining the incident ticket used to document the failed authentication event, remediation, and closure.
\end{itemize}

Its final-state configuration and deployment validation will be documented later (see Section~\ref{sec:appendix-a-ops-glpi-netplan-config} and Section~\ref{sec:glpi-final-configuration}).

% ---------------------------------------------------------------------
% SUBSECTION 4.4: TOPOLOGY DIAGRAM
% ---------------------------------------------------------------------
% This subsection defines the canonical architectural model for the lab
% implementation via Mermaid diagram. This section serves as the normative
% design baseline against which configuration and validation sections
% are evaluated. Diagram source is maintained separately and embedded
% as a static PNG for deterministic HTML and PDF output.
% ---------------------------------------------------------------------

\subsection{4.4 Topology Diagram}
\label{sec:topology-diagram}

Figure~\ref{fig:reference-topology-diagram} is the authoritative logical topology diagram for the lab document. The diagram should preserve the segment boundaries and role relationships defined in Sections~\ref{sec:logical-topology-overview} through~\ref{sec:final-system-roles}.

\begin{quote}
\textit{Diagram implementation note:}\\
The source-of-truth for this diagram is maintained in Mermaid format under
\texttt{diagrams/reference-topology-diagram.mmd}. During CI/CD, Mermaid is rendered to
\texttt{dist/docs/diagrams/reference-topology-diagram.png}, and the document embeds the generated PNG artifact so that both HTML and PDF outputs are consistent and deterministic.
\end{quote}

% -----------------------------------------------------------------------------
% FIGURE: REFERENCE TOPOLOGY DIAGRAM
%
% Source of truth:
%   diagrams/reference-topology-diagram.mmd
%
% Build behavior:
%   Mermaid diagram is rendered during CI/CD into a static PNG artifact
%   at diagrams/reference-topology-diagram.png.
%
% Rationale:
%   - Ensures deterministic PDF builds (no runtime Mermaid dependency).
%   - Maintains single diagram source for web and PDF outputs.
%
% WARNING:
%   If directory structure changes, update both:
%     1. CI/CD render path
%     2. \includegraphics path below
% -----------------------------------------------------------------------------

\begin{figure}[h!]
    \centering
    \includegraphics[width=0.82\textwidth]{diagrams/reference-topology.png}
    \caption{Reference topology diagram for the final virtualized NAC implementation.}
    \label{fig:reference-topology-diagram}
\end{figure}

The diagram is intended to function as the reader's primary visual reference for the rest of the document. Later sections will rely on the topology as the basis for explaining the authenticated access path, the guest-policy enforcement model, the incident flow, and the packet capture locations. Supporting final-state interface and route evidence will be provided later in the appendices and validation sections rather than embedded directly here (see Section~\ref{sec:platform-build-and-final-state-configuration}).

% =====================================================================
% SECTION 5: PLATFORM BUILD AND FINAL-STATE CONFIGURATION
% =====================================================================
% This section documents the final implemented state of the major systems
% in the lab. It is intentionally limited to authoritative final-state
% configuration and service posture. Intermediate troubleshooting states
% are excluded unless later required by the incident narrative.
% =====================================================================

\section{5. Platform Build and Final-State Configuration}
\label{sec:platform-build-and-final-state-configuration}

% ---------------------------------------------------------------------
% SUBSECTION 5.1: HYPERVISOR AND VIRTUAL NETWORK LAYOUT
% ---------------------------------------------------------------------
% This subsection describes the virtualization substrate and the role of
% VMware virtual networks in the architecture.
% ---------------------------------------------------------------------

\subsection{5.1 Hypervisor and Virtual Network Layout}
\label{sec:hypervisor-and-virtual-network-layout}

The lab is hosted on VMware Workstation Pro running on an Ubuntu 24.04 workstation. Distinct VMware virtual networks are used to model access, transit, service, guest, and management domains. These virtual networks are treated as deterministic Layer 2 broadcast domains and are mapped into the final architecture as functional segments rather than as ad hoc connectivity conveniences.

In the final design, the principal virtual network roles are as follows:

\begin{itemize}
    \item a management/NAT network used for host-side administrative reachability to selected virtual machines;
    \item an access segment for the wired endpoint path;
    \item a transit segment between the access node and the intermediate router;
    \item a protected services segment behind pfSense;
    \item a guest segment behind pfSense used for the policy-based guest access simulation.
\end{itemize}

The actual VMware GUI configuration is not reproduced here because it does not materially improve the engineering value of the document. Instead, the document records the final interface-to-segment mapping in operational terms through the final netplan files, interface summaries, and pfSense interface assignments. This is more appropriate for a production-style technical record and avoids overemphasizing hypervisor UI state over the logical architecture.

Temporary host-side IP assignments on VMware host interfaces were used at specific points to reach otherwise isolated pfSense interfaces during setup and validation. Those temporary host assignments are not treated as part of the access-control data path and are not part of the final guest or authenticated endpoint model.

% ---------------------------------------------------------------------
% SUBSECTION 5.2: ROUTER AND FIREWALL ROLES
% ---------------------------------------------------------------------
% This subsection establishes the final-state role of dist-router and
% pfSense as the intermediate routing and policy-enforcement nodes.
% ---------------------------------------------------------------------

\subsection{5.2 Router and Firewall Roles}
\label{sec:router-and-firewall-roles}

The routing and policy boundary in the final design is split across two nodes:

\begin{itemize}
    \item \textbf{\texttt{dist-router}} provides the routed transition from the access domain toward the pfSense transit interface;
    \item \textbf{pfSense} provides the protected services gateway, the guest gateway, guest DHCP, and the firewall enforcement boundary.
\end{itemize}

This separation is important because it makes the architecture more representative of a routed enterprise environment. Rather than allowing the access node to connect directly to the service segment, the design forces service-bound traffic through an intermediate routed hop and then through the firewall boundary before it reaches the protected service network.

The final-state interface and route posture on \texttt{dist-router} confirms this role. The node maintains:

\begin{itemize}
    \item a management/NAT-connected interface for administrative access;
    \item a transit-facing interface on \texttt{10.0.23.0/24} toward \texttt{access-node};
    \item a point-to-point transit interface on \texttt{10.0.34.0/30} toward pfSense;
    \item a static route to the access client segment via \texttt{10.0.23.1};
    \item a static route to the protected services segment via pfSense at \texttt{10.0.34.2}.
\end{itemize}

The final-state interface summary and routing table for \texttt{dist-router} are reproduced below.

\begin{verbatim}
dist-router: ip -br a
lo               UNKNOWN        127.0.0.1/8 ::1/128
ens33            UP             192.168.195.142/24 metric 100 fe80::20c:29ff:fe4f:9ae7/64
ens38            UP             10.0.23.2/24 fe80::20c:29ff:fe4f:9af1/64
ens39            UP             10.0.34.1/30 fe80::20c:29ff:fe4f:9afb/64
\end{verbatim}

\begin{verbatim}
dist-router: ip route
default via 192.168.195.2 dev ens33 proto dhcp src 192.168.195.142 metric 100
10.0.12.0/24 via 10.0.23.1 dev ens38 proto static
10.0.23.0/24 dev ens38 proto kernel scope link src 10.0.23.2
10.0.30.0/24 via 10.0.34.2 dev ens39 proto static
10.0.34.0/30 dev ens39 proto kernel scope link src 10.0.34.1
192.168.195.0/24 dev ens33 proto kernel scope link src 192.168.195.142 metric 100
192.168.195.2 dev ens33 proto dhcp scope link src 192.168.195.142 metric 100
\end{verbatim}

\begin{verbatim}
dist-router: sysctl net.ipv4.ip_forward
net.ipv4.ip_forward = 1
\end{verbatim}

The corresponding final netplan configuration on \texttt{dist-router} is reproduced in full below and is also referenced in Appendix A (see Section~\ref{sec:appendix-a-dist-router-netplan-config}).

\begin{verbatim}
dist-router: /etc/netplan/50-cloud-init.yaml
network:
  version: 2
  ethernets:
    ens33:
      dhcp4: true
    ens38:
      dhcp4: false
      addresses:
        - 10.0.23.2/24
      routes:
        - to: 10.0.12.0/24
          via: 10.0.23.1
    ens39:
      dhcp4: false
      addresses:
        - 10.0.34.1/30
      routes:
        - to: 10.0.30.0/24
          via: 10.0.34.2
\end{verbatim}

pfSense is the authoritative security boundary in the environment. Its final interface structure is as follows:

\begin{itemize}
    \item \textbf{WAN} (\texttt{em0}) on the VMware management/NAT network at \texttt{192.168.195.145/24};
    \item \textbf{LAN / transit} (\texttt{em1}) at \texttt{10.0.34.2/30};
    \item \textbf{OPT1 / services} (\texttt{em2}) at \texttt{10.0.30.1/24};
    \item \textbf{GUEST / OPT2} (\texttt{em3}) at \texttt{10.0.20.254/24}.
\end{itemize}

This interface design allows pfSense to sit directly between the routed access path and the protected services segment, while also hosting a guest segment whose policy can be enforced independently.

The authoritative manual transcription of the final pfSense interface summary is reproduced below and in the appendix (see Section~\ref{sec:appendix-a-pfsense-authenticated-interface-summary}).

\begin{verbatim}
pfSense final interface summary
WAN (wan) -> em0 -> v4/DHCP4: 192.168.195.145/24
LAN (lan) -> em1 -> v4: 10.0.34.2/30
OPT1 (opt1) -> em2 -> v4: 10.0.30.1/24
GUEST (opt2) -> em3 -> v4: 10.0.20.254/24
\end{verbatim}

The associated pfSense interfaces terminal screenshot is shown below:

\begin{figure}[H]
    \centering
    \includegraphics[width=\textwidth]{assets/screenshots/pfSense_Interfaces.png}
    \caption{pfSense interface assignments showing WAN, LAN/transit, OPT1 (services), and GUEST (OPT2)}
    \label{sec:fig-pfsense_interfaces}
\end{figure}

% ---------------------------------------------------------------------
% SUBSECTION 5.3: FREERADIUS FINAL CONFIGURATION
% ---------------------------------------------------------------------
% This subsection records the final-state RADIUS service posture,
% relevant file excerpts, and local validation.
% ---------------------------------------------------------------------

\subsection{5.3 FreeRADIUS Final Configuration}
\label{sec:freeradius-final-configuration}

FreeRADIUS is hosted on \texttt{radius-splunk} and functions as the authoritative authentication and policy service for the wired 802.1X workflow. In the final design, FreeRADIUS accepts requests from the Linux-based authenticator on \texttt{access-node}, validates the known endpoint credential, and provides the Access-Accept or reject path used throughout the validation and incident workflows.

The final-state network posture of \texttt{radius-splunk} is reproduced below.

\begin{verbatim}
radius-splunk: ip -br a
lo               UNKNOWN        127.0.0.1/8 ::1/128
ens33            UP             192.168.195.143/24 metric 100 fe80::20c:29ff:fe82:9c22/64
ens38            UP             10.0.30.10/24 fe80::20c:29ff:fe82:9c2c/64
\end{verbatim}

\begin{verbatim}
radius-splunk: ip route
default via 192.168.195.2 dev ens33 proto dhcp src 192.168.195.143 metric 100
10.0.12.0/24 via 10.0.30.1 dev ens38 proto static
10.0.23.0/24 via 10.0.30.1 dev ens38 proto static
10.0.30.0/24 dev ens38 proto kernel scope link src 10.0.30.10
10.0.34.0/30 via 10.0.30.1 dev ens38 proto static
192.168.195.0/24 dev ens33 proto kernel scope link src 192.168.195.143 metric 100
192.168.195.2 dev ens33 proto dhcp scope link src 192.168.195.143 metric 100
\end{verbatim}

\begin{verbatim}
radius-splunk: /etc/netplan/50-cloud-init.yaml
network:
  version: 2
  ethernets:
    ens33:
      dhcp4: true
    ens38:
      dhcp4: false
      addresses:
        - 10.0.30.10/24
      routes:
        - to: 10.0.34.0/30
          via: 10.0.30.1
        - to: 10.0.23.0/24
          via: 10.0.30.1
        - to: 10.0.12.0/24
          via: 10.0.30.1
\end{verbatim}

The FreeRADIUS service is enabled and active in the final state. Its service status confirms a successful pre-start configuration check and a normal running state.

\begin{verbatim}
radius-splunk: systemctl status freeradius --no-pager
Loaded: loaded (/usr/lib/systemd/system/freeradius.service; enabled; preset: enabled)
Active: active (running) since Tue 2026-03-24 04:24:32 UTC
Main PID: 1162 (freeradius)
Status: "Processing requests"

Recent lines:
Compiling Post-Auth-Type Client-Lost for attr Post-Auth-Type
radiusd: #### Skipping IP addresses and Ports ####
Configuration appears to be OK
Started freeradius.service - FreeRADIUS multi-protocol policy server.
\end{verbatim}

The explicit configuration validation command was also executed successfully in the final state.

\begin{verbatim}
radius-splunk: sudo freeradius -XC
FreeRADIUS Version 3.2.5
Starting - reading configuration files ...
Configuration appears to be OK
\end{verbatim}

The local authorization source contains the lab user entry required for the controlled 802.1X workflow. Only the relevant excerpt is reproduced below, with the shared credential redacted. The full authoritative excerpt is referenced later in Appendix A (see Section~\ref{sec:appendix-a-radius-splunk-freeradius-users-config}).

\begin{verbatim}
radius-splunk: grep excerpt from /etc/freeradius/3.0/mods-config/files/authorize
205-#########################################################
206-
207:nacuser Cleartext-Password := "<REDACTED>"
\end{verbatim}

The RADIUS client definition for the Linux-based authenticator is reproduced below, again with the shared secret redacted. The authoritative appendix target for this excerpt is Section~\ref{sec:appendix-a-radius-splunk-freeradius-clients-config}.

\begin{verbatim}
radius-splunk: grep excerpt from /etc/freeradius/3.0/clients.conf
354:client access-node {
355-    ipaddr = 10.0.23.1
356-    secret = <REDACTED>
357:    shortname = access-node
358-}
\end{verbatim}

A deterministic local RADIUS test using \texttt{radtest} confirms that the configured lab user is accepted by the running service. The command line has been redacted where necessary, but the full output is reproduced below.

\begin{verbatim}
radius-splunk: local radtest
radtest nacuser '<REDACTED_PASSWORD>' 127.0.0.1 0 '<REDACTED_SHARED_SECRET>'
Sent Access-Request Id 54 from 0.0.0.0:45969 to 127.0.0.1:1812 length 77
        User-Name = "nacuser"
        User-Password = "<REDACTED>"
        NAS-IP-Address = 127.0.1.1
        NAS-Port = 0
        Message-Authenticator = 0x00
        Cleartext-Password = "<REDACTED>"
Received Access-Accept Id 54 from 127.0.0.1:1812 to 127.0.0.1:45969 length 38
        Message-Authenticator = 0xe9369309c447d3d395f2e7af45be62b5
\end{verbatim}

Taken together, these artifacts establish that the final FreeRADIUS posture is internally consistent, locally testable, and suitable for the later remote validation and incident scenario sections.

% ---------------------------------------------------------------------
% SUBSECTION 5.4: HOSTAPD WIRED AUTHENTICATOR FINAL CONFIGURATION
% ---------------------------------------------------------------------
% This subsection records the final hostapd configuration and managed
% service posture on access-node.
% ---------------------------------------------------------------------

\subsection{5.4 hostapd Wired Authenticator Final Configuration}
\label{sec:hostapd-wired-authenticator-final-configuration}

The wired authenticator role is implemented on \texttt{access-node} using \texttt{hostapd} in wired mode. This is one of the central design decisions in the lab because it allows the environment to model the IEEE 802.1X control path in a Linux-based virtualized setting while remaining explicit that this does not provide full parity with managed enterprise switch behavior.

The final interface posture of \texttt{access-node} is reproduced below.

\begin{verbatim}
access-node: ip -br a
lo               UNKNOWN        127.0.0.1/8 ::1/128
ens33            UP             192.168.195.141/24 metric 100 fe80::20c:29ff:feef:5e82/64
ens38            UP             10.0.12.1/24 fe80::20c:29ff:feef:5e8c/64
ens39            UP             10.0.23.1/24 fe80::20c:29ff:feef:5e96/64
\end{verbatim}

\begin{verbatim}
access-node: ip route
default via 192.168.195.2 dev ens33 proto dhcp src 192.168.195.141 metric 100
10.0.12.0/24 dev ens38 proto kernel scope link src 10.0.12.1
10.0.23.0/24 dev ens39 proto kernel scope link src 10.0.23.1
10.0.30.0/24 via 10.0.23.2 dev ens39 proto static
192.168.195.0/24 dev ens33 proto kernel scope link src 192.168.195.141 metric 100
192.168.195.2 dev ens33 proto dhcp scope link src 192.168.195.141 metric 100
\end{verbatim}

\begin{verbatim}
access-node: /etc/netplan/50-cloud-init.yaml
network:
  version: 2
  ethernets:
    ens33:
      dhcp4: true
    ens38:
      dhcp4: false
      addresses:
        - 10.0.12.1/24
    ens39:
      dhcp4: false
      addresses:
        - 10.0.23.1/24
      routes:
        - to: 10.0.30.0/24
          via: 10.0.23.2
\end{verbatim}

\begin{verbatim}
access-node: sysctl net.ipv4.ip_forward
net.ipv4.ip_forward = 1
\end{verbatim}

The final \texttt{hostapd} wired configuration is reproduced below, with the shared secret redacted. The authoritative appendix target for this file is Section~\ref{sec:appendix-a-access-node-hostapd-wired-config}.

\begin{verbatim}
access-node: /etc/hostapd/lab/hostapd-wired.conf
interface=ens38
driver=wired

ieee8021x=1
eapol_version=2
eap_reauth_period=3600

auth_server_addr=10.0.30.10
auth_server_port=1812
auth_server_shared_secret=<REDACTED>

own_ip_addr=10.0.23.1

logger_stdout=-1
logger_stdout_level=0
logger_syslog=-1
logger_syslog_level=0

ctrl_interface=/var/run/hostapd
ctrl_interface_group=0
\end{verbatim}

To make the authenticator role operationally realistic and manageable, the lab does not rely solely on ad hoc foreground invocations. Instead, the final design includes a dedicated systemd unit, reproduced below and referenced later in Appendix A (see Section~\ref{sec:appendix-a-access-node-hostapd-wired-systemd-unit}).

\begin{verbatim}
access-node: /etc/systemd/system/hostapd-wired.service
[Unit]
Description=hostapd wired 802.1X authenticator (lab)
After=network-online.target
Wants=network-online.target

[Service]
Type=simple
ExecStart=/usr/sbin/hostapd -dd /etc/hostapd/lab/hostapd-wired.conf
Restart=on-failure
RestartSec=2

[Install]
WantedBy=multi-user.target
\end{verbatim}

The final service posture confirms that the unit is enabled and running, and that it is bound to the intended interface and RADIUS server.

\begin{verbatim}
access-node: systemctl status hostapd-wired.service --no-pager
Loaded: loaded (/etc/systemd/system/hostapd-wired.service; enabled; preset: enabled)
Active: active (running) since Tue 2026-03-24 04:20:41 UTC
Main PID: 975 (hostapd)

Recent lines:
Using interface ens38 with hwaddr <REDACTED> and ssid ""
ens38: RADIUS Authentication server 10.0.30.10:1812
RADIUS local address: 10.0.23.1:45833
ens38: interface state UNINITIALIZED->ENABLED
ens38: AP-ENABLED
ens38: Setup of interface done.
\end{verbatim}

A later section of this document will use the managed \texttt{hostapd-wired.service} unit as a core part of the Ansible validation and remediation workflow. Its existence therefore has architectural significance beyond convenience: it makes the authenticator role inspectable, restartable, and automatable in a way that is operationally meaningful.

% ---------------------------------------------------------------------
% SUBSECTION 5.5: WPA_SUPPLICANT FINAL CONFIGURATION
% ---------------------------------------------------------------------
% This subsection records the endpoint-side supplicant configuration and
% clarifies the relationship between the running system service and the
% controlled test invocation model used later in the incident section.
% ---------------------------------------------------------------------

\subsection{5.5 wpa\_supplicant Final Configuration}
\label{sec:wpa-supplicant-final-configuration}

The wired endpoint role is implemented on \texttt{client-1} using \texttt{wpa\_supplicant}. The final-state client network posture and the final controlled configuration file are reproduced below.

\begin{verbatim}
client-1: ip -br a
lo               UNKNOWN        127.0.0.1/8 ::1/128
ens33            UP             10.0.12.10/24 fe80::20c:29ff:fe68:5333/64
\end{verbatim}

\begin{verbatim}
client-1: ip route
default via 10.0.12.1 dev ens33 proto static
10.0.12.0/24 dev ens33 proto kernel scope link src 10.0.12.10
\end{verbatim}

\begin{verbatim}
client-1: /etc/netplan/50-cloud-init.yaml
network:
  version: 2
  ethernets:
    ens33:
      addresses:
        - 10.0.12.10/24
      routes:
        - to: default
          via: 10.0.12.1
\end{verbatim}

The final wired supplicant configuration is reproduced below, with the secret redacted. The authoritative appendix target for this file is Section~\ref{sec:appendix-a-client-1-wpa-supplicant-wired-8021x-config}.

\begin{verbatim}
client-1: /etc/wpa_supplicant/lab/wired-8021x.conf
ctrl_interface=/run/wpa_supplicant
ap_scan=0
eapol_version=2

network={
    key_mgmt=IEEE8021X
    eap=MD5
    identity="nacuser"
    password="<REDACTED>"
}
\end{verbatim}

In the final steady state at the time of evidence capture, the generic \texttt{wpa\_supplicant.service} unit is active:

\begin{verbatim}
client-1: systemctl status wpa_supplicant --no-pager
Loaded: loaded (/usr/lib/systemd/system/wpa_supplicant.service; enabled; preset: enabled)
Active: active (running) since Tue 2026-03-24 04:57:02 UTC
Main PID: 830 (wpa_supplicant)

Recent lines:
Starting wpa_supplicant.service - WPA supplicant...
Successfully initialized wpa_supplicant
Started wpa_supplicant.service - WPA supplicant.
\end{verbatim}

This final steady-state service condition should not be confused with the controlled test methodology used later in the incident section. During the failed-authentication simulation and subsequent remediation validation, the \texttt{wpa\_supplicant} service was deliberately stopped and the binary was invoked manually in debug mode against the lab-specific wired configuration file in order to obtain deterministic, isolated evidence. The final active service state reproduced here is therefore a valid representation of the system at rest, while the incident section will separately document the intentional service stop/start choices used during testing.

% ---------------------------------------------------------------------
% SUBSECTION 5.6: SPLUNK UNIVERSAL FORWARDER FINAL CONFIGURATION
% ---------------------------------------------------------------------
% This subsection records the final UF posture and monitored inputs.
% ---------------------------------------------------------------------

\subsection{5.6 Splunk Universal Forwarder Final Configuration}
\label{sec:splunk-universal-forwarder-final-configuration}

Splunk Universal Forwarder is installed on \texttt{radius-splunk} as a lightweight observability component. In this lab, it is not used as a full SIEM platform with local indexing and search-head analytics. Rather, it is used to demonstrate an enterprise-relevant log collection and monitoring posture for NAC-related evidence sources.

The final runtime state of the forwarder is reproduced below.

\begin{verbatim}
radius-splunk: /opt/splunkforwarder/bin/splunk status
Warning: Attempting to revert the SPLUNK_HOME ownership
Warning: Executing "chown -R root /opt/splunkforwarder"
splunkd is running (PID: 1671).
splunk helpers are running (PIDs: 1672).
\end{verbatim}

The monitored input definitions are maintained in a dedicated local app and are reproduced in full below. The authoritative appendix target for this file is Section~\ref{sec:appendix-a-splunk-uf-inputs-conf}.

\begin{verbatim}
radius-splunk: /opt/splunkforwarder/etc/apps/lab_nac_inputs/local/inputs.conf
[monitor:///var/log/syslog]
disabled = false
index = main
sourcetype = syslog

[monitor:///var/log/auth.log]
disabled = false
index = main
sourcetype = linux_secure

[monitor:///var/log/freeradius/radius.log]
disabled = false
index = main
sourcetype = freeradius
\end{verbatim}

These inputs are sufficient to demonstrate that the environment includes structured monitoring of the most relevant local evidence sources for the NAC scenario:

\begin{itemize}
    \item generic system events via \texttt{/var/log/syslog};
    \item authentication and privilege-related events via \texttt{/var/log/auth.log};
    \item FreeRADIUS service activity via \texttt{/var/log/freeradius/radius.log}.
\end{itemize}

This design choice is intentionally modest but meaningful. It allows the lab to claim concrete Splunk integration in a way that is accurate for the implemented scope, without overstating the environment as a full Splunk analytics deployment.

% ---------------------------------------------------------------------
% SUBSECTION 5.7: GLPI FINAL CONFIGURATION
% ---------------------------------------------------------------------
% This subsection records the final-state web stack posture and the role
% of GLPI as an active ITSM component in the lab.
% ---------------------------------------------------------------------

\subsection{5.7 GLPI Final Configuration}
\label{sec:glpi-final-configuration}

GLPI is hosted on \texttt{ops-glpi} and functions as the IT service management platform for the environment. It is not present only as a demonstration of software installation; it is used directly in the operational scenario to represent both the managed endpoint asset and the authentication incident ticket.

The final-state network posture of \texttt{ops-glpi} is reproduced below.

\begin{verbatim}
ops-glpi: ip -br a
lo               UNKNOWN        127.0.0.1/8 ::1/128
ens33            UP             192.168.195.144/24 metric 100 fe80::20c:29ff:fe46:c6a1/64
ens38            UP             10.0.30.20/24 fe80::20c:29ff:fe46:c6ab/64
\end{verbatim}

\begin{verbatim}
ops-glpi: ip route
default via 192.168.195.2 dev ens33 proto dhcp src 192.168.195.144 metric 100
10.0.12.0/24 via 10.0.30.1 dev ens38 proto static
10.0.23.0/24 via 10.0.30.1 dev ens38 proto static
10.0.30.0/24 dev ens38 proto kernel scope link src 10.0.30.20
10.0.34.0/30 via 10.0.30.1 dev ens38 proto static
192.168.195.0/24 dev ens33 proto kernel scope link src 192.168.195.144 metric 100
192.168.195.2 dev ens33 proto dhcp scope link src 192.168.195.144 metric 100
\end{verbatim}

\begin{verbatim}
ops-glpi: /etc/netplan/50-cloud-init.yaml
network:
  version: 2
  ethernets:
    ens33:
      dhcp4: true
    ens38:
      dhcp4: false
      addresses:
        - 10.0.30.20/24
      routes:
        - to: 10.0.34.0/30
          via: 10.0.30.1
        - to: 10.0.23.0/24
          via: 10.0.30.1
        - to: 10.0.12.0/24
          via: 10.0.30.1
\end{verbatim}

The Apache and MariaDB services are both enabled and active in the final state.

\begin{verbatim}
ops-glpi: systemctl status apache2 --no-pager
Loaded: loaded (/usr/lib/systemd/system/apache2.service; enabled; preset: enabled)
Active: active (running) since Tue 2026-03-24 05:22:11 UTC
Main PID: 1296 (apache2)

Recent lines:
Starting apache2.service - The Apache HTTP Server...
Started apache2.service - The Apache HTTP Server.
\end{verbatim}

\begin{verbatim}
ops-glpi: systemctl status mariadb --no-pager
Loaded: loaded (/usr/lib/systemd/system/mariadb.service; enabled; preset: enabled)
Active: active (running) since Tue 2026-03-24 05:22:12 UTC
Main PID: 1151 (mariadbd)
Status: "Taking your SQL requests now..."

Recent lines:
Server socket created on IP: '127.0.0.1', port: '3306'.
/usr/sbin/mariadbd: ready for connections.
Started mariadb.service - MariaDB 10.11.14 database server.
\end{verbatim}

Local HTTP validation confirms that GLPI is exposed at the root path in the final Apache configuration, and that the historical \texttt{/glpi/} subpath is no longer the correct application path.

\begin{verbatim}
ops-glpi: curl -I http://127.0.0.1/
HTTP/1.1 200 OK
Date: Tue, 24 Mar 2026 05:26:56 GMT
Server: Apache/2.4.58 (Ubuntu)
Set-Cookie: <REDACTED>
Expires: Thu, 19 Nov 1981 08:52:00 GMT
Cache-Control: no-store, no-cache, must-revalidate
Pragma: no-cache
Cache-Control: no-cache, private
Vary: Accept-Encoding
Content-Type: text/html; charset=UTF-8
\end{verbatim}

\begin{verbatim}
ops-glpi: curl -I http://127.0.0.1/glpi/
HTTP/1.1 404 Not Found
Date: Tue, 24 Mar 2026 05:28:14 GMT
Server: Apache/2.4.58 (Ubuntu)
Set-Cookie: <REDACTED>
Expires: Mon, 26 Jul 1997 05:00:00 GMT
Cache-Control: no-store, no-cache, must-revalidate
Pragma: no-cache
x-frame-options: SAMEORIGIN
Cache-Control: no-cache, private
Content-Type: text/html; charset=UTF-8
\end{verbatim}

The required PHP environment is present, including the previously installed \texttt{bcmath} extension required for GLPI.

\begin{verbatim}
ops-glpi: php -m | grep bcmath
bcmath
\end{verbatim}

Operationally, GLPI is not treated as a passive accessory to the environment. In the final state, it contains:

\begin{itemize}
    \item a managed asset record for \texttt{client-1};
    \item an incident ticket documenting the simulated wired 802.1X authentication failure and remediation workflow;
    \item a resolved ticket state reflecting successful remediation.
\end{itemize}

The associated GLPI GUI screenshots are shown below:

\begin{figure}[H]
    \centering
    \includegraphics[width=\textwidth]{assets/screenshots/GLPI_client-1_computer_asset.png}
    \caption{GLPI managed asset record for \texttt{client-1} used in the NAC incident workflow}
    \label{sec:fig-glpi_client-1_computer_asset}
\end{figure}

\begin{figure}[H]
    \centering
    \includegraphics[width=\textwidth]{assets/screenshots/GLPI_client-1_incident_ticket.png}
    \caption{GLPI incident ticket for the wired 802.1X authentication failure and remediation workflow affecting \texttt{client-1}}
    \label{sec:fig-glpi_client-1_incident_ticket}
\end{figure}

% ---------------------------------------------------------------------
% SUBSECTION 5.8: PYTHON AND ANSIBLE AUTOMATION ARTIFACTS
% ---------------------------------------------------------------------
% This subsection records the presence, role, and storage location of the
% automation artifacts in the final platform state. The full raw text of
% these artifacts is deferred to Appendix A for readability.
% ---------------------------------------------------------------------

\subsection{5.8 Python and Ansible Automation Artifacts}
\label{sec:python-and-ansible-automation-artifacts}

The final platform state includes two automation artifacts that are operationally relevant to the lab:

\begin{itemize}
    \item a Python-based incident evidence parser hosted on \texttt{radius-splunk};
    \item an Ansible-based validation and optional remediation workflow, also hosted on \texttt{radius-splunk}.
\end{itemize}

These artifacts are treated as part of the implemented platform rather than as external supporting material because they are stored, executed, and maintained within the lab environment itself. Their purpose is not to inflate the apparent sophistication of the project, but to provide repeatable operational utility directly tied to the incident scenario described later in this document.

The Python parser is maintained under the local incident-evidence directory on \texttt{radius-splunk}. Its purpose is to consume the saved incident artifacts from the failed-authentication and successful-remediation workflows and produce a deterministic summary of whether the expected failure and recovery indicators are present. The authoritative full source of this parser is reproduced in Appendix A (see Section~\ref{sec:appendix-a-python-nac-incident-parser}).

The Ansible workflow is also maintained on \texttt{radius-splunk} and consists of:

\begin{itemize}
    \item an inventory file defining the local and remote NAC-related nodes;
    \item a validation and optional remediation playbook capable of checking service state, collecting recent logs, and conditionally restarting the core NAC services when explicitly requested by the operator.
\end{itemize}

The automation workflow uses SSH key-based authentication rather than hardcoded credentials, consistent with the document's stated security and publication posture. The authoritative full inventory and playbook contents are reproduced in Appendix A (see Section~\ref{sec:appendix-a-ansible-inventory-hosts-ini} and Section~\ref{sec:appendix-a-ansible-nac-validation-playbook}).

The presence of these artifacts in the final platform state is significant for two reasons. First, they demonstrate that the environment includes repeatable operational tooling rather than only static configuration. Second, they provide a bridge between the platform build described in this section and the incident-handling and validation workflows described later in Sections~\ref{sec:incident-scenario-wired-8021x-authentication-failure-and-remediation} and~\ref{sec:automation-workflows}.

% =====================================================================
% SECTION 6: ACCESS CONTROL MODEL AND POLICY ENFORCEMENT
% =====================================================================
% This section documents the final access-control model for the lab and
% validates that authenticated and guest traffic are handled according
% to the intended routed and firewall-enforced policy boundaries.
% =====================================================================

\section{6. Access Control Model and Policy Enforcement}
\label{sec:access-control-model-and-policy-enforcement}

The lab implements a deliberately segmented access model intended to resemble a simplified enterprise campus access design. Rather than treating the environment as a flat set of reachable subnets, the architecture distinguishes between authenticated endpoint access, protected services access, transit-only routing segments, and a separately governed guest segment.

This section documents the intended policy model and then validates that the environment behaves accordingly using direct path testing from both the authenticated endpoint and a dedicated guest validation node.

% ---------------------------------------------------------------------
% SUBSECTION 6.1: POLICY MODEL OVERVIEW
% ---------------------------------------------------------------------
% This subsection defines the two principal trust domains in the lab,
% the expected behavior of each domain, and the internal networks from
% which the guest segment is intentionally isolated.
% ---------------------------------------------------------------------

\subsection{6.1 Policy Model Overview}
\label{sec:policy-model-overview}

The final access-control model consists of two principal endpoint trust domains:

\begin{itemize}
    \item an \textbf{authenticated endpoint domain} represented by \texttt{client-1} on \texttt{10.0.12.0/24}, which is expected to reach protected services after successful wired 802.1X authentication;
    \item a \textbf{guest domain} represented by \texttt{guest-test-1} on \texttt{10.0.20.0/24}, which is expected to receive DHCP service and outbound internet access, but not internal access to protected or transit segments.
\end{itemize}

The routing and enforcement chain for authenticated access is:

\begin{center}
\texttt{client-1} $\rightarrow$ \texttt{access-node} $\rightarrow$ \texttt{dist-router} $\rightarrow$ \texttt{pfSense} $\rightarrow$ protected services
\end{center}

The guest enforcement chain is simpler:

\begin{center}
\texttt{guest-test-1} $\rightarrow$ \texttt{pfSense (GUEST)} $\rightarrow$ internet or explicitly permitted destinations
\end{center}

The protected services segment is represented by \texttt{10.0.30.0/24}, which hosts:

\begin{itemize}
    \item \texttt{radius-splunk} at \texttt{10.0.30.10};
    \item \texttt{ops-glpi} at \texttt{10.0.30.20}.
\end{itemize}

The guest segment is represented by \texttt{10.0.20.0/24} and is intentionally isolated from the following internal networks:

\begin{itemize}
    \item \texttt{10.0.30.0/24} (protected services);
    \item \texttt{10.0.34.0/30} (pfSense transit);
    \item \texttt{10.0.23.0/24} (access-to-distribution transit);
    \item \texttt{10.0.12.0/24} (authenticated endpoint/access segment).
\end{itemize}

This is not intended to model a full enterprise NAC enforcement platform with dynamic downloadable ACLs or per-session VLAN reassignment. Instead, it demonstrates a controlled and explainable policy boundary in which wired 802.1X authentication gates access into a routed path toward protected services, while a separate guest domain is explicitly restricted by firewall policy.

% ---------------------------------------------------------------------
% SUBSECTION 6.2: AUTHENTICATED ENDPOINT ACCESS MODEL
% ---------------------------------------------------------------------
% This subsection documents the final-state posture of the authenticated
% endpoint and validates that it can reach the protected services path
% after successful wired 802.1X authentication.
% ---------------------------------------------------------------------

\subsection{6.2 Authenticated Endpoint Access Model}
\label{sec:authenticated-endpoint-access-model}

The authenticated endpoint path is represented by \texttt{client-1}, which resides on \texttt{10.0.12.0/24} behind the Linux-based wired authenticator. In the final steady state, \texttt{client-1} uses a static address and default route via \texttt{access-node}.

The final-state endpoint interface and routing posture are reproduced below.

\begin{verbatim}
client-1: ip -br a
lo               UNKNOWN        127.0.0.1/8 ::1/128
ens33            UP             10.0.12.10/24 fe80::20c:29ff:fe68:5333/64
\end{verbatim}

\begin{verbatim}
client-1: ip route
default via 10.0.12.1 dev ens33 proto static
10.0.12.0/24 dev ens33 proto kernel scope link src 10.0.12.10
\end{verbatim}

The endpoint-side \texttt{wpa\_supplicant} service is active in the final steady state, which confirms that the endpoint is not merely configured but is operating in a post-build state consistent with the intended access model.

\begin{verbatim}
client-1: systemctl status wpa_supplicant --no-pager
Loaded: loaded (/usr/lib/systemd/system/wpa_supplicant.service; enabled; preset: enabled)
Active: active (running) since Tue 2026-03-24 08:03:28 UTC
Main PID: 829 (wpa_supplicant)

Recent lines:
Starting wpa_supplicant.service - WPA supplicant...
Started wpa_supplicant.service - WPA supplicant.
Successfully initialized wpa_supplicant
\end{verbatim}

From the endpoint perspective, successful access to the local gateway is verified first:

\begin{verbatim}
client-1: ping -c 4 10.0.12.1
4 packets transmitted, 4 received, 0% packet loss
\end{verbatim}

This confirms that the endpoint has Layer 3 reachability to the local routed access gateway on \texttt{access-node}. More importantly, the endpoint can also reach both protected services nodes on the \texttt{10.0.30.0/24} segment.

\begin{verbatim}
client-1: ping -c 4 10.0.30.10
4 packets transmitted, 4 received, 0% packet loss
\end{verbatim}

\begin{verbatim}
client-1: ping -c 4 10.0.30.20
4 packets transmitted, 4 received, 0% packet loss
\end{verbatim}

Application-layer validation to GLPI is also successful from the authenticated endpoint, which is stronger evidence than ICMP alone because it confirms that the routed path reaches a live service rather than merely a responsive host.

\begin{verbatim}
client-1: curl -I --max-time 5 http://10.0.30.20/
HTTP/1.1 200 OK
Date: Tue, 24 Mar 2026 08:21:04 GMT
Server: Apache/2.4.58 (Ubuntu)
Set-Cookie: <REDACTED>
Expires: Thu, 19 Nov 1981 08:52:00 GMT
Cache-Control: no-store, no-cache, must-revalidate
Pragma: no-cache
Cache-Control: no-cache, private
Vary: Accept-Encoding
Content-Type: text/html; charset=UTF-8
\end{verbatim}

This result is particularly useful because it shows that the authenticated endpoint can not only reach the protected services subnet but can directly access the GLPI web application on the services segment.

A \texttt{tracepath} test to \texttt{10.0.30.20} also supports the multi-hop design and shows that the path is not flat:

\begin{verbatim}
client-1: tracepath 10.0.30.20
 1?: [LOCALHOST]                      pmtu 1500
 1:  _gateway                                              0.911ms
 1:  _gateway                                              0.566ms
 2:  ???                                                   1.713ms
 3:  ???                                                   2.833ms
 4:  ???                                                   3.034ms reached
     Resume: pmtu 1500 hops 4 back 4
\end{verbatim}

Because intermediate Linux routers and pfSense do not necessarily reveal useful names or ICMP TTL-expired responses in a lab, the anonymous hop presentation is not treated as a defect. The important point is that the path is clearly multi-hop and terminates successfully at the protected service host.

% ---------------------------------------------------------------------
% SUBSECTION 6.3: ROUTED PATH VALIDATION FOR AUTHENTICATED ACCESS
% ---------------------------------------------------------------------
% This subsection validates the authenticated endpoint path from the
% infrastructure perspective by showing the explicit forward and return
% routing behavior across the access, distribution, and firewall layers.
% ---------------------------------------------------------------------

\subsection{6.3 Routed Path Validation for Authenticated Access}
\label{sec:routed-path-validation-for-authenticated-access}

The endpoint-side tests above establish that the authenticated client can reach the protected services segment. This subsection validates the same behavior from the infrastructure perspective to show that the path is explicitly routed and not the result of accidental adjacency.

On \texttt{access-node}, the route lookup for both protected service hosts confirms that traffic is forwarded to \texttt{dist-router} at \texttt{10.0.23.2} via interface \texttt{ens39}.

\begin{verbatim}
access-node: ip route get 10.0.30.10
10.0.30.10 via 10.0.23.2 dev ens39 src 10.0.23.1 uid 1000
\end{verbatim}

\begin{verbatim}
access-node: ip route get 10.0.30.20
10.0.30.20 via 10.0.23.2 dev ens39 src 10.0.23.1 uid 1000
\end{verbatim}

On \texttt{dist-router}, the route lookup confirms that both protected service hosts are then forwarded toward pfSense at \texttt{10.0.34.2} via interface \texttt{ens39}.

\begin{verbatim}
dist-router: ip route get 10.0.30.10
10.0.30.10 via 10.0.34.2 dev ens39 src 10.0.34.1 uid 1000
\end{verbatim}

\begin{verbatim}
dist-router: ip route get 10.0.30.20
10.0.30.20 via 10.0.34.2 dev ens39 src 10.0.34.1 uid 1000
\end{verbatim}

This proves the intended forward path:

\begin{center}
\texttt{client-1} $\rightarrow$ \texttt{access-node} $\rightarrow$ \texttt{dist-router} $\rightarrow$ \texttt{pfSense} $\rightarrow$ \texttt{10.0.30.0/24}
\end{center}

The return path is also explicitly validated from both protected service hosts. On \texttt{radius-splunk}, the route to the client subnet points to pfSense at \texttt{10.0.30.1}:

\begin{verbatim}
radius-splunk: ip route get 10.0.12.10
10.0.12.10 via 10.0.30.1 dev ens38 src 10.0.30.10 uid 1000
\end{verbatim}

On \texttt{ops-glpi}, the route to the client subnet is identical in principle:

\begin{verbatim}
ops-glpi: ip route get 10.0.12.10
10.0.12.10 via 10.0.30.1 dev ens38 src 10.0.30.20 uid 1000
\end{verbatim}

This confirms the intended return path:

\begin{center}
\texttt{10.0.30.0/24} $\rightarrow$ \texttt{pfSense} $\rightarrow$ \texttt{dist-router} $\rightarrow$ \texttt{access-node} $\rightarrow$ \texttt{client-1}
\end{center}

Bidirectional reachability is then directly validated with ICMP from both protected service hosts back to \texttt{client-1}.

\begin{verbatim}
radius-splunk: ping -c 4 10.0.12.10
4 packets transmitted, 4 received, 0% packet loss
\end{verbatim}

\begin{verbatim}
ops-glpi: ping -c 4 10.0.12.10
4 packets transmitted, 4 received, 0% packet loss
\end{verbatim}

This combination of forward-path route lookup, return-path route lookup, and bidirectional ICMP provides strong evidence that the authenticated endpoint path is functioning as a deliberately routed design rather than as an incidental lab shortcut.

% ---------------------------------------------------------------------
% SUBSECTION 6.4: GUEST NETWORK POLICY MODEL
% ---------------------------------------------------------------------
% This subsection records the authoritative guest-segment DHCP and
% firewall policy on pfSense and defines the intended isolation model
% for the lower-trust guest network.
% ---------------------------------------------------------------------

\subsection{6.4 Guest Network Policy Model}
\label{sec:guest-network-policy-model}

The guest network is implemented as a distinct pfSense-controlled segment on \texttt{10.0.20.0/24}. It is not intended to participate in the wired 802.1X access path and is instead used to demonstrate a separate, lower-trust network with DHCP service, firewall-enforced internal isolation, and outbound internet access.

The authoritative final pfSense interface definition for the guest segment was already established in Section~\ref{sec:router-and-firewall-roles}:

\begin{itemize}
    \item \textbf{GUEST / OPT2} (\texttt{em3}) at \texttt{10.0.20.254/24}
\end{itemize}

The final DHCP configuration for the guest segment is manually transcribed from pfSense as follows:

\begin{verbatim}
pfSense GUEST DHCP summary
Enabled (Enable DHCP server on GUEST interface)
Subnet: 10.0.20.0/24
Subnet Range: 10.0.20.1 - 10.0.20.254
Address Pool Range: 10.0.20.100 - 10.0.20.199
\end{verbatim}

The associated pfSense GUI screenshot is shown below:

\begin{figure}[H]
    \centering
    \includegraphics[width=\textwidth]{assets/screenshots/pfSense_DHCP_Server_GUEST.png}
    \caption{pfSense GUEST DHCP server configuration for the policy-based guest access segment}
    \label{sec:fig-pfsense_dhcp_server_guest}
\end{figure}

The guest firewall policy is implemented as an ordered first-match ruleset on the pfSense GUEST interface. The final manually transcribed rules are reproduced below and were explicitly re-confirmed as unchanged and in the same top-to-bottom order at the time of Section 6 validation.

% ---------------------------------------------------------------------
% NOTE: This ruleset is reproduced in full because rule order is central to
% the guest-policy interpretation and the top-to-bottom sequence is part
% of the technical meaning of the final pfSense configuration.
% ---------------------------------------------------------------------

\begin{verbatim}
pfSense GUEST rules (top-to-bottom)
Pass  | IPv4* | Source: 10.0.20.0/24 | Destination: This Firewall (Self)
        | ALLOW-GUEST-TO-PFSENSE
Block | IPv4* | Source: 10.0.20.0/24 | Destination: 10.0.30.0/24       
        | BLOCK-GUEST-TO-SERVICES
Block | IPv4* | Source: 10.0.20.0/24 | Destination: 10.0.34.0/30       
        | BLOCK-GUEST-TO-TRANSIT
Block | IPv4* | Source: 10.0.20.0/24 | Destination: 10.0.23.0/24       
        | BLOCK-GUEST-TO-DIST
Block | IPv4* | Source: 10.0.20.0/24 | Destination: 10.0.12.0/24       
        | BLOCK-GUEST-TO-AUTH-ACCESS
Pass  | IPv4* | Source: 10.0.20.0/24 | Destination: *                  
        | ALLOW-GUEST-TO-INTERNET-AND-OTHER-NON-INTERNAL
\end{verbatim}

The associated pfSense GUI screenshot is shown below:

\begin{figure}[H]
    \centering
    \includegraphics[width=\textwidth]{assets/screenshots/pfSense_Firewall_Rules_GUEST.png}
    \caption{pfSense GUEST interface firewall rules enforcing internal isolation and controlled outbound access}
    \label{sec:fig-pfsense_firewall_rules_guest}
\end{figure}

This ruleset is deliberately explicit. Rather than relying on an implicit deny plus a broad allow, the internal networks are individually enumerated and blocked before the final permissive rule. That makes the guest policy easier to explain, easier to audit, and more suitable for publication as a portfolio artifact.

At the time of validation, pfSense outbound NAT was confirmed to be in Automatic mode, which is relevant because it explains why the guest network can successfully access the internet without requiring separately documented manual NAT rules.

% ---------------------------------------------------------------------
% SUBSECTION 6.5: GUEST DHCP AND ISOLATION VALIDATION
% ---------------------------------------------------------------------
% This subsection validates guest-segment behavior with a dedicated
% single-homed test VM by confirming DHCP lease assignment, gateway
% reachability, and enforced failure to reach protected internal paths.
% ---------------------------------------------------------------------

\subsection{6.5 Guest DHCP and Isolation Validation}
\label{sec:guest-dhcp-and-isolation-validation}

To validate the guest policy empirically rather than only by configuration review, a dedicated disposable guest-only test VM named \texttt{guest-test-1} was created and attached solely to the pfSense GUEST segment. Importantly, this VM was provisioned with a single interface only and no management/NAT interface, so that its observed behavior reflects the actual guest policy boundary rather than an alternate administrative path.

The guest validation node received a DHCP lease from pfSense in the expected guest subnet and within the configured guest address pool.

\begin{verbatim}
guest-test-1: ip -br a
lo               UNKNOWN        127.0.0.1/8 ::1/128
ens33            UP             10.0.20.137/24 fe80::250:56ff:feab:cdef/64
\end{verbatim}

\begin{verbatim}
guest-test-1: ip route
default via 10.0.20.254 dev ens33 proto dhcp src 10.0.20.137 metric 100
10.0.20.0/24 dev ens33 proto kernel scope link src 10.0.20.137
10.0.20.254 dev ens33 proto dhcp scope link src 10.0.20.137 metric 100
\end{verbatim}

An explicit DHCP renewal confirms the expected pfSense DHCP behavior and the lease origin.

% ---------------------------------------------------------------------
% NOTE: This DHCP renewal transcript is included because it provides direct
% lease-exchange evidence that pfSense is the active DHCP authority for
% the guest segment during validation.
% ---------------------------------------------------------------------

\begin{verbatim}
guest-test-1: sudo dhclient -v ens33
DHCPDISCOVER on ens33 to 255.255.255.255 port 67 interval 3 (xid=0x4a3f2c1d)
DHCPDISCOVER on ens33 to 255.255.255.255 port 67 interval 6 (xid=0x4a3f2c1d)
DHCPOFFER of 10.0.20.137 from 10.0.20.254
DHCPREQUEST for 10.0.20.137 on ens33 to 255.255.255.255 port 67 (xid=0x4a3f2c1d)
DHCPACK of 10.0.20.137 from 10.0.20.254 (xid=0x4a3f2c1d)
bound to 10.0.20.137 -- renewal in 3600 seconds.
\end{verbatim}

The guest node can successfully reach the pfSense GUEST gateway itself, which is consistent with the first allow rule in the guest ruleset.

\begin{verbatim}
guest-test-1: ping -c 4 10.0.20.254
4 packets transmitted, 4 received, 0% packet loss
\end{verbatim}

The guest node is then tested against each explicitly blocked internal segment. All of these attempts fail as intended.

Protected services segment (\texttt{10.0.30.0/24}):

\begin{verbatim}
guest-test-1: ping -c 4 10.0.30.10
4 packets transmitted, 0 received, 100% packet loss
\end{verbatim}

\begin{verbatim}
guest-test-1: ping -c 4 10.0.30.20
4 packets transmitted, 0 received, 100% packet loss
\end{verbatim}

Transit segment toward pfSense (\texttt{10.0.34.0/30}):

\begin{verbatim}
guest-test-1: ping -c 4 10.0.34.1
4 packets transmitted, 0 received, 100% packet loss
\end{verbatim}

Access-to-distribution transit (\texttt{10.0.23.0/24}):

\begin{verbatim}
guest-test-1: ping -c 4 10.0.23.1
4 packets transmitted, 0 received, 100% packet loss
\end{verbatim}

Authenticated access segment (\texttt{10.0.12.0/24}):

\begin{verbatim}
guest-test-1: ping -c 4 10.0.12.1
4 packets transmitted, 0 received, 100% packet loss
\end{verbatim}

Application-layer access to GLPI from the guest node also fails by timeout, which is a particularly strong validation point because it demonstrates that the guest cannot reach a protected service even when directly attempting the published web application path.

\begin{verbatim}
guest-test-1: curl -I --max-time 5 http://10.0.30.20/
curl: (28) Connection timed out after 5001 milliseconds
\end{verbatim}

This confirms that the guest segment is not only blocked from protected hosts at the ICMP layer but is also effectively prevented from reaching the protected application service.

% ---------------------------------------------------------------------
% SUBSECTION 6.6: GUEST INTERNET ACCESS VALIDATION
% ---------------------------------------------------------------------
% This subsection confirms that the guest network remains functional for
% general outbound connectivity and is not merely isolated from internal
% networks by broad denial.
% ---------------------------------------------------------------------

\subsection{6.6 Guest Internet Access Validation}
\label{sec:guest-internet-access-validation}

Because the guest ruleset ends with an explicit allow rule for destinations outside the enumerated internal networks, it is useful to validate that the guest network retains general outbound access rather than being globally constrained.

The guest validation node can successfully reach a public IP address:

\begin{verbatim}
guest-test-1: ping -c 4 1.1.1.1
4 packets transmitted, 4 received, 0% packet loss
\end{verbatim}

It can also successfully retrieve HTTP headers from an external public site:

\begin{verbatim}
guest-test-1: curl -I --max-time 5 http://www.google.com
HTTP/1.1 200 OK
Content-Type: text/html; charset=ISO-8859-1
...
\end{verbatim}

This behavior is consistent with the intended guest policy model:

\begin{itemize}
    \item access to pfSense itself is permitted;
    \item access to enumerated internal networks is blocked;
    \item general outbound access is permitted.
\end{itemize}

In other words, the guest segment is not simply ``isolated'' in the abstract. It is specifically implemented as a constrained but functional internet-facing guest network.

% ---------------------------------------------------------------------
% SUBSECTION 6.7: POLICY INTERPRETATION AND LIMITATIONS
% ---------------------------------------------------------------------
% This subsection summarizes the validated policy outcomes and states the
% architectural limits of the lab so that the access model is described
% accurately and not overstated.
% ---------------------------------------------------------------------

\subsection{6.7 Policy Interpretation and Limitations}
\label{sec:policy-interpretation-and-limitations}

The validation results in this section support the following conclusions:

\begin{enumerate}
    \item The authenticated endpoint path is functional and routed as designed.
    \item The protected services segment is reachable from the authenticated endpoint and from no guest path tested.
    \item The guest network receives DHCP service correctly and is restricted from the enumerated internal networks.
    \item The guest network remains operational for general outbound internet access.
\end{enumerate}

This is a meaningful and defensible policy model for the scope of the lab. However, it is important to state its limits precisely.

First, the guest isolation is enforced by a conventional firewall ruleset on pfSense, not by dynamic NAC policy derived from endpoint identity or posture. That is acceptable for this design because the guest network is intentionally modeled as a separate trust zone rather than as a dynamically assigned 802.1X result.

Second, the authenticated endpoint path is validated at the network and application layers, but the environment does not implement per-session VLAN assignment, downloadable ACLs, posture checks, or certificate-based EAP methods. Those omissions are addressed later in the limitations section and do not invalidate the current demonstration.

Third, the Linux \texttt{hostapd} wired authenticator is used as a practical lab approximation of a switch-based 802.1X authenticator. It is appropriate for demonstrating the control-plane exchange and the operational workflow, but it should not be misrepresented as a drop-in replacement for enterprise access switch behavior.

Within those stated limits, the access-control model is internally consistent, technically validated, and appropriate for inclusion in a portfolio document that aims to present a realistic but accurately scoped virtualized NAC implementation.

% ---------------------------------------------------------------------
% SUBSECTION 6.8: WIRESHARK CROSS-REFERENCE FOR POLICY VALIDATION
% ---------------------------------------------------------------------
% This subsection links the access-control conclusions in Section 6 to
% the packet-capture evidence reviewed later, noting that the policy
% model is also supported by observable 802.1X control-plane behavior.
% ---------------------------------------------------------------------

\subsection{6.8 Wireshark Cross-Reference for Policy Validation}
\label{sec:wireshark-cross-reference-for-policy-validation}

Although the full packet-analysis narrative is reserved for the dedicated packet capture section later in this document, the Wireshark evidence already supports the access-control interpretation presented here.

At the time of review, the packet captures were confirmed to show that the authenticated endpoint traffic is preceded by EAPOL exchange rather than immediately presenting as ordinary post-authenticated data-plane traffic. This is important because it visually reinforces the intended distinction between:

\begin{itemize}
    \item the \textbf{pre-authentication control plane}, represented by EAPOL and RADIUS-driven 802.1X exchange; and
    \item the \textbf{post-authentication data plane}, represented by ordinary IP traffic after successful authorization.
\end{itemize}

The screenshot inventory includes all four relevant packet-capture perspectives:

\begin{itemize}
    \item failed client-side capture;
    \item failed authenticator-side capture;
    \item success client-side capture;
    \item success authenticator-side capture.
\end{itemize}

These artifacts will be incorporated in the dedicated packet-capture analysis section rather than duplicated here. Their relevance in this section is to reinforce that the access model is not inferred only from IP reachability tests; it is also observable at the protocol level through the expected 802.1X control-plane behavior.

% =====================================================================
% SECTION 7: INCIDENT SCENARIO -- WIRED 802.1X AUTHENTICATION FAILURE
%            AND REMEDIATION
% =====================================================================
% This section presents the central operational scenario of the lab: a
% controlled wired 802.1X authentication failure affecting client-1,
% followed by investigation, remediation, and successful recovery
% validation. The section is intentionally evidence-driven and ties
% together endpoint logs, authenticator behavior, RADIUS evidence, GLPI
% tracking, and the Python parser summary.
% =====================================================================

\section{7. Incident Scenario: Wired 802.1X Authentication Failure and Remediation}
\label{sec:incident-scenario-wired-8021x-authentication-failure-and-remediation}

The final lab environment is not presented solely as a static implementation. It is also used to simulate and document a realistic operational incident in which a wired endpoint fails 802.1X authentication, remains unauthorized, is investigated through multiple evidence sources, and is restored to service through a controlled remediation workflow.

This section documents that incident as the authoritative operational scenario for the environment. The incident validation date used throughout this section is 2026-03-22.

% ---------------------------------------------------------------------
% SUBSECTION 7.1: INCIDENT DEFINITION
% ---------------------------------------------------------------------
% This subsection defines the simulated incident, identifies the control
% path involved, and establishes the operational goals of the failure-
% injection and recovery workflow.
% ---------------------------------------------------------------------

\subsection{7.1 Incident Definition}
\label{sec:incident-definition}

The incident concerns the managed endpoint \texttt{client-1}, which is the same endpoint represented in GLPI and used throughout the access-control validation workflow. The endpoint normally authenticates through the following control path:

\begin{center}
\texttt{client-1} $\rightarrow$ \texttt{access-node} $\rightarrow$ \texttt{radius-splunk}
\end{center}

Within this path:

\begin{itemize}
    \item \texttt{client-1} acts as the wired supplicant;
    \item \texttt{access-node} acts as the Linux-based wired 802.1X authenticator using \texttt{hostapd};
    \item \texttt{radius-splunk} hosts FreeRADIUS as the authentication and policy service.
\end{itemize}

The incident was defined as a wired 802.1X authentication failure caused by invalid supplicant credentials. This was not an accidental outage but a controlled failure injection introduced for the purpose of validating the lab's ability to:

\begin{itemize}
    \item detect a failed NAC event;
    \item observe the failed state across the supplicant, authenticator, and RADIUS layers;
    \item document the event in an ITSM workflow;
    \item restore service using a known-good configuration;
    \item validate successful recovery through both host-side evidence and packet analysis.
\end{itemize}

The scenario therefore serves as a practical proof that the lab can do more than demonstrate component reachability. It can also support a realistic access-control failure narrative from symptom to closure.

% ---------------------------------------------------------------------
% SUBSECTION 7.2: FAILURE INJECTION METHOD
% ---------------------------------------------------------------------
% This subsection documents the deliberate credential change used to
% produce the failed authentication state and explains why that method
% provides a controlled and reproducible endpoint-side failure.
% ---------------------------------------------------------------------

\subsection{7.2 Failure Injection Method}
\label{sec:failure-injection-method}

The failure was introduced by deliberately replacing the correct wired supplicant password in the \texttt{client-1} \texttt{wpa\_supplicant} configuration with an incorrect value. This method was chosen because it is deterministic, easy to revert, and directly relevant to the type of endpoint-side credential error that could plausibly occur in an enterprise environment.

This approach has several advantages in a proof-of-concept lab:

\begin{itemize}
    \item it guarantees a failure condition at the authentication layer rather than at unrelated network layers;
    \item it preserves the normal routing and firewall design, isolating the incident to the NAC control path;
    \item it makes the remediation equally deterministic by restoring the known-good client configuration;
    \item it allows clean collection of failure and recovery evidence without introducing ambiguity about the cause of the event.
\end{itemize}

The lab therefore treats the credential modification as a formal failure-injection technique rather than as an incidental misconfiguration. This is important because it makes the incident reproducible and allows the later automation and packet-analysis workflows to operate against a known and explainable failure mode.

% ---------------------------------------------------------------------
% SUBSECTION 7.3: AUTHORITATIVE INCIDENT ARTIFACTS
% ---------------------------------------------------------------------
% This subsection records the authoritative incident-evidence inventory
% and the parser summary used to synthesize the failure and recovery
% artifacts across the supplicant, authenticator, and RADIUS layers.
% ---------------------------------------------------------------------

\subsection{7.3 Authoritative Incident Artifacts}
\label{sec:authoritative-incident-artifacts}

The incident evidence set is maintained on \texttt{radius-splunk} under the local incident-evidence directory. The final directory inventory at the time of documentation is reproduced below.

% ---------------------------------------------------------------------
% NOTE: This directory listing is included because it establishes the retained
% incident artifact set and the specific files that support the failure,
% remediation, and parser-based summary described in this section.
% ---------------------------------------------------------------------

\begin{verbatim}
radius-splunk: ls -lh ~/incident-evidence
total 92K
-rw-rw-r-- 1 nac-lab nac-lab  32K Mar 23 04:02 freeradius_failed_auth_incident.txt
-rw-rw-r-- 1 nac-lab nac-lab  11K Mar 23 04:01 hostapd_failed_auth_incident.txt
-rw-rw-r-- 1 nac-lab nac-lab  19K Mar 23 04:01 hostapd_success_retest_v2.txt
-rw-r--r-- 1 root    root    3.5K Mar 24 09:52 nac_incident_parser.py
-rw-rw-r-- 1 nac-lab nac-lab  535 Mar 23 04:28 nac_incident_parser_success.txt
-rw-rw-r-- 1 nac-lab nac-lab 7.3K Mar 23 03:59 wpa_supplicant_failed_auth_incident.txt
-rw-rw-r-- 1 nac-lab nac-lab 8.4K Mar 23 03:59 wpa_supplicant_success_retest_v2.txt
\end{verbatim}

These files are sufficient to support the full narrative of failure and recovery across the endpoint, authenticator, and RADIUS layers. They also form the input set for the Python incident parser.

At the time of final documentation, the parser remained functional and produced the following authoritative summary:

% ---------------------------------------------------------------------
% NOTE: This parser output is included because it provides a concise cross-
% artifact summary of the incident evidence and captures the final
% operational conclusion produced by the Python validation tool.
% ---------------------------------------------------------------------

\begin{verbatim}
radius-splunk: python3 ~/incident-evidence/nac_incident_parser.py
=== NAC Incident Evidence Summary ===

[1] Failed authentication simulation
- Client remained unauthorized: YES
- Client showed successful connection during failed test: NO
- hostapd saw authentication attempt: YES
- FreeRADIUS recorded failed auth / reject path: YES

[2] Remediation / successful re-test
- Client reached CTRL-EVENT-CONNECTED: YES
- Client reached Authorized state: YES
- hostapd showed successful authorization indicators: YES

[3] Incident conclusion
- SIMULATED NAC INCIDENT SUCCESSFULLY REPRODUCED AND REMEDIATED
\end{verbatim}

This parser output is significant because it provides a deterministic, cross-file synthesis of the incident evidence. It does not replace the underlying artifacts, but it does provide a concise operational summary that agrees with the raw evidence reproduced below.

The full parser source is maintained separately as a platform artifact in Appendix A (see Section~\ref{sec:appendix-a-python-nac-incident-parser}).

% ---------------------------------------------------------------------
% SUBSECTION 7.4: FAILED AUTHENTICATION EVIDENCE
% ---------------------------------------------------------------------
% This subsection establishes the failed state across the three control
% points in the authentication path so that the incident is supported by
% independent endpoint, authenticator, and RADIUS evidence.
% ---------------------------------------------------------------------

\subsection{7.4 Failed Authentication Evidence}
\label{sec:failed-authentication-evidence}

The failed state is established independently at all three control points: the supplicant, the authenticator, and the RADIUS service.

\subsubsection{7.4.1 Supplicant-Side Failed State}
\label{sec:supplicant-side-failed-state}

On the endpoint side, the failed-authentication capture shows that the supplicant remained unauthorized and never transitioned into a connected or authorized state.

\begin{verbatim}
client-side failed indicators
EAPOL: Supplicant port status: Unauthorized
EAPOL: Supplicant port status: Unauthorized
\end{verbatim}

This is the clearest endpoint-side symptom of the failure. There is no successful connection indicator in the failed-authentication evidence set.

\subsubsection{7.4.2 Authenticator-Side Failed State}
\label{sec:authenticator-side-failed-state}

On the authenticator side, \texttt{hostapd} recorded repeated attempts to send the authentication transaction toward the RADIUS server, confirming that the endpoint was actively attempting authentication and that the authenticator was processing the request path.

\begin{verbatim}
hostapd failed-auth indicators
Mar 22 22:32:33 access-node hostapd[2052]: ens38: RADIUS Sending RADIUS message to 
        authentication server
Mar 22 22:32:33 access-node hostapd[2052]: ens38: RADIUS Sending RADIUS message to 
        authentication server
Mar 22 22:32:33 access-node hostapd[2052]: ens38: RADIUS Sending RADIUS message to 
        authentication server
Mar 22 22:32:33 access-node hostapd[2052]: ens38: RADIUS Sending RADIUS message to 
        authentication server
\end{verbatim}

This is operationally important because it shows that the failure was not caused by a silent client or by an absence of authenticator activity. The authenticator was actively participating in the 802.1X control path and attempting to obtain a RADIUS decision.

\subsubsection{7.4.3 RADIUS-Side Failed State}
\label{sec:radius-side-failed-state}

The strongest proof of the failure appears in the FreeRADIUS evidence. The failed-authentication excerpt shows an EAP failure, rejection of the authentication attempt, and the use of the reject post-auth path.

\begin{verbatim}
FreeRADIUS failed-auth indicators
(1) eap: Sending EAP Failure (code 4) ID 26 length 4
(1)     [eap] = reject
(1)   } # authenticate = reject
(1) Failed to authenticate the user
(1) Using Post-Auth-Type Reject
(1)   Post-Auth-Type REJECT {
(1) attr_filter.access_reject: EXPAND %{User-Name}
(1) attr_filter.access_reject:    --> nacuser
(1) attr_filter.access_reject: Matched entry DEFAULT at line 11
(1)     [attr_filter.access_reject] = updated
(1)   } # Post-Auth-Type REJECT = updated
\end{verbatim}

This confirms that the failed state is not merely an endpoint perception. The policy authority itself rejected the authentication request and generated the expected EAP failure path.

Taken together, the endpoint, authenticator, and RADIUS evidence establish the failed state conclusively.

% ---------------------------------------------------------------------
% SUBSECTION 7.5: INVESTIGATION AND TRIAGE
% ---------------------------------------------------------------------
% This subsection summarizes the investigation logic used to distinguish
% among endpoint, authenticator, and RADIUS failure domains and records
% the evidence-based triage conclusions reached during the incident.
% ---------------------------------------------------------------------

\subsection{7.5 Investigation and Triage}
\label{sec:investigation-and-triage}

The investigation phase in this incident was intentionally simple but operationally credible. The goal was to distinguish among the most likely causes of a failed wired 802.1X event:

\begin{itemize}
    \item an endpoint-side credential or supplicant configuration error;
    \item an authenticator-side service or forwarding problem;
    \item a RADIUS-side policy or service failure.
\end{itemize}

The incident evidence and the surrounding validation steps supported the following triage conclusions:

\begin{itemize}
    \item the endpoint remained unauthorized, so the event was not a mere transient or cosmetic log anomaly;
    \item \texttt{hostapd} was actively sending RADIUS requests, which meant the authenticator path was engaged;
    \item FreeRADIUS was running and able to produce a deterministic reject path, which meant the backend service was operational;
    \item the known-good user credentials still produced Access-Accept under \texttt{radtest}, isolating the problem to the client-side credential state rather than to the FreeRADIUS policy model.
\end{itemize}

This investigation logic is reflected directly in the GLPI incident record, which documents verification of FreeRADIUS health, confirmation that the known-good credentials still produced Access-Accept, and revalidation of the client \texttt{wpa\_supplicant} configuration.

% ---------------------------------------------------------------------
% SUBSECTION 7.6: REMEDIATION ACTIONS
% ---------------------------------------------------------------------
% This subsection documents the targeted remediation applied to restore
% the known-good supplicant credential state on client-1 and re-run the
% authentication path without broader architectural changes.
% ---------------------------------------------------------------------

\subsection{7.6 Remediation Actions}
\label{sec:remediation-actions}

The remediation was deliberately minimal and targeted. Rather than changing the architecture, altering the firewall rules, or modifying RADIUS policy, the known-good \texttt{wpa\_supplicant} credential state on \texttt{client-1} was restored and the wired 802.1X authentication test was executed again under controlled conditions.

This remediation method is significant because it demonstrates a proper engineering response to a localized failure:

\begin{itemize}
    \item identify the probable failing control point;
    \item restore known-good state at that control point;
    \item re-test the full path rather than assuming success;
    \item validate the result across multiple evidence sources.
\end{itemize}

No broad changes to routing, pfSense policy, or FreeRADIUS user logic were required to restore service. That outcome is itself meaningful because it shows the failure was introduced and resolved at the endpoint credential layer rather than by systemic platform instability.

% ---------------------------------------------------------------------
% SUBSECTION 7.7: SUCCESSFUL RE-TEST EVIDENCE
% ---------------------------------------------------------------------
% This subsection documents the successful recovery validation and shows
% that the remediated endpoint returned to connected and authorized
% state across the same control points used to establish the failure.
% ---------------------------------------------------------------------

\subsection{7.7 Successful Re-Test Evidence}
\label{sec:successful-re-test-evidence}

The remediation was validated through a clean success re-test. As with the failed state, the successful state is visible at all three control points.

\subsubsection{7.7.1 Supplicant-Side Recovery Evidence}
\label{sec:supplicant-side-recovery-evidence}

The successful re-test capture shows the expected transition from initial unauthorized state to connected and authorized state, followed by a later return to unauthorized state when the controlled test process ended.

\begin{verbatim}
client-side success indicators
EAPOL: Supplicant port status: Unauthorized
ens33: CTRL-EVENT-CONNECTED - Connection to <REDACTED> completed [id=0 id_str=]
EAPOL: Supplicant port status: Authorized
EAPOL: Supplicant port status: Unauthorized
\end{verbatim}

The final \texttt{Unauthorized} line is not interpreted as a failed remediation. It is consistent with termination of the controlled test process after the successful exchange had already been completed. The decisive indicators are \texttt{CTRL-EVENT-CONNECTED} and \texttt{Authorized}.

\subsubsection{7.7.2 Authenticator-Side Recovery Evidence}
\label{sec:authenticator-side-recovery-evidence}

The successful re-test on the authenticator shows the expected sequence of RADIUS interaction, Access-Accept handling, EAP Success decapsulation, station connection, and port authorization.

% ---------------------------------------------------------------------
% NOTE: This hostapd excerpt is included because it preserves the full
% authenticator-side success sequence from upstream RADIUS handling
% through EAP Success, station connection, and controlled-port
% authorization.
% ---------------------------------------------------------------------

\begin{verbatim}
hostapd success indicators
Mar 23 02:49:00 access-node hostapd[2052]: Encapsulating EAP message into a RADIUS 
        packet
Mar 23 02:49:00 access-node hostapd[2052]: Copied RADIUS State Attribute
Mar 23 02:49:00 access-node hostapd[2052]: ens38: RADIUS Sending RADIUS message to 
        authentication server
Mar 23 02:49:00 access-node hostapd[2052]: RADIUS message: code=1 (Access-Request) 
        identifier=4 length=193
Mar 23 02:49:00 access-node hostapd[2052]: ens38: STA <REDACTED> IEEE 802.1X: 
        unauthorizing port
Mar 23 02:49:00 access-node hostapd[2052]: ens38: RADIUS Next RADIUS client retransmit in 
        3 seconds
Mar 23 02:49:00 access-node hostapd[2052]: ens38: RADIUS Received 53 bytes from RADIUS 
        server
Mar 23 02:49:00 access-node hostapd[2052]: ens38: RADIUS Received RADIUS message
Mar 23 02:49:00 access-node hostapd[2052]: RADIUS message: code=2 (Access-Accept) 
        identifier=4 length=53
Mar 23 02:49:00 access-node hostapd[2052]: ens38: STA <REDACTED> RADIUS: Received 
        RADIUS packet matched with a pending request, round trip time 0.00 sec
Mar 23 02:49:00 access-node hostapd[2052]: RADIUS packet matching with station 
        <REDACTED>
Mar 23 02:49:00 access-node hostapd[2052]: ens38: STA <REDACTED> IEEE 802.1X: 
        decapsulated EAP packet (code=3 id=255 len=4) from RADIUS server: EAP Success
Mar 23 02:49:00 access-node hostapd[2052]: ens38: CTRL-EVENT-EAP-SUCCESS2 <REDACTED>
Mar 23 02:49:00 access-node hostapd[2052]: ens38: AP-STA-CONNECTED <REDACTED>
Mar 23 02:49:00 access-node hostapd[2052]: ens38: STA <REDACTED> IEEE 802.1X: 
        authorizing port
Mar 23 02:49:00 access-node hostapd[2052]: ens38: STA <REDACTED> RADIUS: starting 
        accounting session 2268479955CB85A0
\end{verbatim}

This is the strongest authenticator-side proof of successful remediation. The sequence explicitly shows:
- RADIUS request handling,
- Access-Accept reception,
- EAP Success,
- station connection,
- and port authorization.

\subsubsection{7.7.3 RADIUS-Side Recovery Interpretation}
\label{sec:radius-side-recovery-interpretation}

A dedicated successful FreeRADIUS transaction excerpt was not preserved in the same compact text form as the failed-authentication reject path. However, this does not weaken the conclusion because the success state is already corroborated by:

\begin{itemize}
    \item the endpoint-side connected and authorized state;
    \item the authenticator-side Access-Accept and EAP Success handling;
    \item the parser summary, which concluded that the incident was successfully reproduced and remediated.
\end{itemize}

In this case, the successful RADIUS outcome is most directly evidenced through the authenticator's decoding of the received RADIUS Accept and EAP Success path.

% ---------------------------------------------------------------------
% SUBSECTION 7.8: GLPI INCIDENT RECORD
% ---------------------------------------------------------------------
% This subsection records the formal GLPI incident entry associated with
% the simulated authentication failure and ties the technical evidence
% to an operational ticketing workflow.
% ---------------------------------------------------------------------

\subsection{7.8 GLPI Incident Record}
\label{sec:glpi-incident-record}

The incident was not left as a purely technical artifact. It was also recorded in GLPI as an actual operational record associated with the lab asset model. This is important because it demonstrates that the environment integrates ITSM process with technical evidence rather than treating ticketing as an unrelated side service.

At the time of documentation, the incident ticket title and status were as follows:

\begin{verbatim}
Ticket title:
Wired 802.1X authentication failure and remediation for client-1

Ticket status:
Solved
\end{verbatim}

The full incident text recorded in GLPI is reproduced below.

% ---------------------------------------------------------------------
% NOTE: This GLPI incident text is reproduced because it captures the
% formal operational record of the simulated event, including observed
% behavior, investigation steps, remediation actions, and recovery
% validation.
% ---------------------------------------------------------------------

\begin{verbatim}
Incident summary:
client-1 failed wired 802.1X authentication during NAC validation after the supplicant 
        password was intentionally changed to an incorrect value to simulate a credential
        issue.

Observed behavior:
- client-1 remained unauthorized during the failed test
- FreeRADIUS showed a failed authentication / reject path
- hostapd on access-node recorded the 802.1X authentication attempt

Investigation:
- Verified FreeRADIUS service health on radius-splunk
- Confirmed the known-good credentials still produced Access-Accept using radtest
- Revalidated the client wpa_supplicant configuration

Remediation:
- Restored the known-good wpa_supplicant credentials on client-1
- Re-ran the wired 802.1X authentication test

Recovery validation:
- client-1 showed CTRL-EVENT-CONNECTED
- client-1 supplicant port reached Authorized state
- access-node hostapd logs showed AP-STA-CONNECTED and port authorization

Impact:
This simulated a realistic enterprise NAC incident involving endpoint credential failure 
        and demonstrated successful troubleshooting and service restoration.
\end{verbatim}

For associate GLPI GUI screenshots, see Section~\ref{sec:glpi-final-configuration}.

% ---------------------------------------------------------------------
% SUBSECTION 7.9: INCIDENT CONCLUSION
% ---------------------------------------------------------------------
% This subsection closes the incident narrative by summarizing the full
% failure-and-recovery lifecycle and stating its significance for the
% lab as an operationally credible NAC validation environment.
% ---------------------------------------------------------------------

\subsection{7.9 Incident Conclusion}
\label{sec:incident-conclusion}

The incident scenario demonstrates that the lab can support a complete wired NAC failure-and-recovery lifecycle:

\begin{enumerate}
    \item a controlled authentication failure was introduced at the endpoint credential layer;
    \item the endpoint remained unauthorized;
    \item the authenticator processed and forwarded the failed authentication path;
    \item FreeRADIUS generated the expected reject and EAP failure path;
    \item the issue was investigated without unnecessary architectural changes;
    \item the known-good client configuration was restored;
    \item a clean re-test demonstrated connected and authorized state;
    \item the event was recorded and closed in GLPI.
\end{enumerate}

This is one of the most important outcomes in the document because it transforms the lab from a static collection of services into an operational environment that can support:

\begin{itemize}
    \item failure simulation;
    \item layered evidence collection;
    \item repeatable remediation;
    \item cross-domain validation using routing, logging, ITSM, automation, and packet capture.
\end{itemize}

The later automation and packet-capture sections build directly on this incident narrative. The Python parser summarizes the same evidence set reproduced here, and the Wireshark screenshots provide protocol-level confirmation of both the failed and successful authentication paths.

% =====================================================================
% SECTION 8: AUTOMATION WORKFLOWS
% =====================================================================
% This section documents the automation layer implemented around the NAC
% lab. The goal is not to present automation as an isolated scripting
% exercise, but as an operational control capability supporting service
% validation, optional remediation, and evidence collection.
% =====================================================================

\section{8. Automation Workflows}
\label{sec:automation-workflows}

The lab includes a lightweight but operationally meaningful automation layer designed to support repeatable validation of the wired 802.1X control path and the associated authentication services. Rather than attempting to implement a full orchestration or SOAR platform, the automation focuses on three practical goals:

\begin{itemize}
    \item validate the state of the authenticator and RADIUS services from a central node;
    \item optionally remediate service availability by restarting critical daemons;
    \item preserve local evidence artifacts for later review and appendix inclusion.
\end{itemize}

This design is best characterized as validation-first with optional service remediation. That phrasing is deliberate. The automation is primarily an evidence and control workflow, not a blind auto-restart mechanism.

The automation controller for this lab is \texttt{radius-splunk}. From that node, Ansible is used to validate the Linux-based wired 802.1X authenticator on \texttt{access-node} and the FreeRADIUS service on \texttt{radius-splunk} itself. A separate Python parser is then used to summarize the incident evidence set collected during the failure-and-remediation scenario described in Section~\ref{sec:incident-scenario-wired-8021x-authentication-failure-and-remediation}.

% ---------------------------------------------------------------------
% SUBSECTION 8.1: AUTOMATION ASSET INVENTORY
% ---------------------------------------------------------------------
% This subsection records the final automation workspace layout, the
% retained artifact set, and the authoritative playbook and parser files
% used by the validation and incident-summary workflows.
% ---------------------------------------------------------------------

\subsection{8.1 Automation Asset Inventory}
\label{sec:automation-asset-inventory}

The final automation workspace resides under \texttt{/home/nac-lab/ansible-nac} on \texttt{radius-splunk}. At the time of final documentation, the file layout was as follows.

\begin{verbatim}
radius-splunk: find ~/ansible-nac -maxdepth 3 -type f | sort
/home/nac-lab/ansible-nac/artifacts/access-node_hostapd_logs.txt
/home/nac-lab/ansible-nac/artifacts/access-node_hostapd_restart_result.txt
/home/nac-lab/ansible-nac/artifacts/access-node_hostapd_status_after.txt
/home/nac-lab/ansible-nac/artifacts/access-node_hostapd_status_before.txt
/home/nac-lab/ansible-nac/artifacts/access-node_hostapd_status.txt
/home/nac-lab/ansible-nac/artifacts/radius-splunk_freeradius_journal.txt
/home/nac-lab/ansible-nac/artifacts/radius-splunk_freeradius_radiuslog.txt
/home/nac-lab/ansible-nac/artifacts/radius-splunk_freeradius_restart_result.txt
/home/nac-lab/ansible-nac/artifacts/radius-splunk_freeradius_status_after.txt
/home/nac-lab/ansible-nac/artifacts/radius-splunk_freeradius_status_before.txt
/home/nac-lab/ansible-nac/artifacts/radius-splunk_freeradius_status.txt
/home/nac-lab/ansible-nac/inventory/hosts.ini
/home/nac-lab/ansible-nac/playbooks/nac_validation.yml
\end{verbatim}

The current artifact directory contents are reproduced below.

\begin{verbatim}
radius-splunk: ls -lh ~/ansible-nac/artifacts
total 36K
-rw-r--r-- 1 nac-lab nac-lab 3.5K Mar 24 12:23 access-node_hostapd_logs.txt
-rw-r--r-- 1 nac-lab nac-lab    0 Mar 24 12:22 access-node_hostapd_restart_result.txt
-rw-r--r-- 1 nac-lab nac-lab 1.4K Mar 24 12:23 access-node_hostapd_status_after.txt
-rw-r--r-- 1 nac-lab nac-lab 1.4K Mar 24 12:23 access-node_hostapd_status_before.txt
-rw-r--r-- 1 root    root    1.4K Mar 23 19:30 access-node_hostapd_status.txt
-rw-r--r-- 1 nac-lab nac-lab 4.0K Mar 24 12:23 radius-splunk_freeradius_journal.txt
-rw-r--r-- 1 nac-lab nac-lab 3.1K Mar 24 12:23 radius-splunk_freeradius_radiuslog.txt
-rw-r--r-- 1 nac-lab nac-lab    0 Mar 24 12:22 radius-splunk_freeradius_restart_result.txt
-rw-r--r-- 1 nac-lab nac-lab 1.8K Mar 24 12:23 radius-splunk_freeradius_status_after.txt
-rw-r--r-- 1 nac-lab nac-lab 1.8K Mar 24 12:23 radius-splunk_freeradius_status_before.txt
-rw-r--r-- 1 root    root    1.8K Mar 23 19:30 radius-splunk_freeradius_status.txt
\end{verbatim}

This inventory is useful for two reasons. First, it shows that the automation is not ephemeral; it leaves behind auditable local outputs. Second, it shows a visible distinction between earlier legacy artifacts created under a pre-revision execution model and the current artifacts created by the final, corrected playbook as \texttt{nac-lab}.

The authoritative revised playbook and the Python incident parser were both confirmed in place with the following timestamps and sizes:

\begin{verbatim}
radius-splunk: ls -lh ~/ansible-nac/playbooks/nac_validation.yml
-rw-r--r-- 1 root root 5.6K Mar 24 12:19 /home/nac-lab/ansible-nac/playbooks/
        nac_validation.yml

radius-splunk: ls -lh ~/incident-evidence/nac_incident_parser.py
-rw-r--r-- 1 root root 3.5K Mar 24 09:52 /home/nac-lab/incident-evidence/
        nac_incident_parser.py
\end{verbatim}

The full inventory file, the authoritative revised playbook, and the full parser source are reproduced in Appendix A (see Section~\ref{sec:appendix-a-ansible-inventory-hosts-ini}, Section~\ref{sec:appendix-a-ansible-nac-validation-playbook}, and Section~\ref{sec:appendix-a-python-nac-incident-parser}).

% ---------------------------------------------------------------------
% SUBSECTION 8.2: PLAYBOOK DESIGN AND REVISION RATIONALE
% ---------------------------------------------------------------------
% This subsection describes the final Ansible playbook design and the
% corrections that established the revised workflow as the authoritative
% implementation for validation, remediation, and artifact handling.
% ---------------------------------------------------------------------

\subsection{8.2 Playbook Design and Revision Rationale}
\label{sec:playbook-design-and-revision-rationale}

The final Ansible workflow uses a single playbook, \texttt{nac\_validation.yml}, which was intentionally kept compact and understandable. Its functional responsibilities are:

\begin{enumerate}
    \item collect a basic hostname from each target;
    \item validate the \texttt{hostapd-wired.service} state on \texttt{access-node};
    \item validate the \texttt{freeradius.service} state on \texttt{radius-splunk};
    \item optionally restart one or both services when \texttt{remediate=true} is provided;
    \item re-check service state after optional remediation;
    \item collect recent journal excerpts and recent FreeRADIUS application log output;
    \item write all collected outputs into a local artifact directory on the control node.
\end{enumerate}

The playbook underwent a meaningful correction during final validation and the corrected version is treated as the authoritative implementation for this document.

Specifically, the final playbook:

\begin{itemize}
    \item sets \texttt{become: false} at the play scope;
    \item applies \texttt{become: true} only to tasks that require elevated privileges;
    \item uses a fixed \texttt{artifact\_dir} variable (\texttt{/home/nac-lab/ansible-nac/artifacts});
    \item explicitly sets \texttt{become: false} on delegated localhost copy tasks.
\end{itemize}

This revision is not a cosmetic change. It materially improves operational correctness by preventing two undesirable behaviors:

\begin{itemize}
    \item unnecessary sudo prompts for harmless non-privileged tasks;
    \item accidental dependence on the invoking user's \texttt{HOME} when writing delegated localhost artifacts.
\end{itemize}

In practical terms, the revised playbook makes execution deterministic and avoids the SSH and artifact-path ambiguity that can occur when operators attempt to run Ansible itself under \texttt{sudo}. The final document therefore treats the revised playbook as the production version of the automation workflow.

% ---------------------------------------------------------------------
% SUBSECTION 8.3: VALIDATION-ONLY EXECUTION
% ---------------------------------------------------------------------
% This subsection documents the default validation-only playbook run and
% shows how the workflow checks service state and writes artifacts
% without attempting service restarts.
% ---------------------------------------------------------------------

\subsection{8.3 Validation-Only Execution}
\label{sec:validation-only-execution}

The first execution mode validates service state without attempting remediation. This is the default behavior and is the recommended mode for routine health checks.

The command used was:

\begin{verbatim}
cd ~/ansible-nac
ansible-playbook -i inventory/hosts.ini playbooks/nac_validation.yml -K
\end{verbatim}

The authoritative execution output is reproduced below.

% ---------------------------------------------------------------------
% NOTE: This Ansible console output is included because it shows the
% complete task path for the validation-only run, including skipped
% remediation tasks and successful artifact-writing behavior.
% ---------------------------------------------------------------------

\begin{verbatim}
PLAY [Validate NAC services and optionally remediate] **************************

TASK [Collect basic hostname] **************************************************
ok: [radius-splunk]
ok: [access-node]

TASK [Check hostapd-wired service on access-node] ******************************
skipping: [radius-splunk]
ok: [access-node]

TASK [Check freeradius service on radius-splunk] *******************************
skipping: [access-node]
ok: [radius-splunk]

TASK [Restart hostapd-wired service when remediation is enabled] ***************
skipping: [radius-splunk]
skipping: [access-node]

TASK [Restart freeradius service when remediation is enabled] ******************
skipping: [access-node]
skipping: [radius-splunk]

TASK [Re-check hostapd-wired service on access-node] ***************************
skipping: [radius-splunk]
ok: [access-node]

TASK [Re-check freeradius service on radius-splunk] ****************************
skipping: [access-node]
ok: [radius-splunk]

TASK [Collect recent hostapd logs] *********************************************
skipping: [radius-splunk]
ok: [access-node]

TASK [Collect recent freeradius service logs] **********************************
skipping: [access-node]
ok: [radius-splunk]

TASK [Collect recent freeradius application log] *******************************
skipping: [access-node]
ok: [radius-splunk]

TASK [Save hostapd pre-check status locally] ***********************************
skipping: [radius-splunk]
changed: [access-node -> localhost]

TASK [Save hostapd post-check status locally] **********************************
skipping: [radius-splunk]
changed: [access-node -> localhost]

TASK [Save hostapd restart result locally] *************************************
skipping: [radius-splunk]
changed: [access-node -> localhost]

TASK [Save hostapd logs locally] ***********************************************
skipping: [radius-splunk]
ok: [access-node -> localhost]

TASK [Save freeradius pre-check status locally] ********************************
skipping: [access-node]
changed: [radius-splunk -> localhost]

TASK [Save freeradius post-check status locally] *******************************
skipping: [access-node]
changed: [radius-splunk -> localhost]

TASK [Save freeradius restart result locally] **********************************
skipping: [access-node]
changed: [radius-splunk -> localhost]

TASK [Save freeradius service logs locally] ************************************
skipping: [access-node]
ok: [radius-splunk -> localhost]

TASK [Save freeradius application log locally] *********************************
skipping: [access-node]
ok: [radius-splunk -> localhost]

PLAY RECAP *********************************************************************
access-node                : ok=8    changed=3    unreachable=0    failed=0    skipped=11   
        rescued=0    ignored=0
radius-splunk              : ok=10   changed=3    unreachable=0    failed=0    skipped=9    
        rescued=0    ignored=0
\end{verbatim}

This run demonstrates the intended baseline behavior:

\begin{itemize}
    \item both managed nodes are reachable;
    \item the service validation tasks complete successfully;
    \item the remediation tasks are correctly skipped when not requested;
    \item evidence files are still written to the local artifact directory.
\end{itemize}

This is the preferred operational mode for routine health checks because it confirms service state and preserves evidence without making any changes to the environment.

% ---------------------------------------------------------------------
% SUBSECTION 8.4: REMEDIATION-ENABLED EXECUTION
% ---------------------------------------------------------------------
% This subsection documents the remediation-enabled playbook run and
% shows the same validation workflow operating with controlled service
% restart actions on the managed nodes.
% ---------------------------------------------------------------------

\subsection{8.4 Remediation-Enabled Execution}
\label{sec:remediation-enabled-execution}

The second execution mode enables controlled service remediation by setting \texttt{remediate=true}. In this mode, the playbook still performs the same validation and evidence collection steps, but it additionally restarts the managed services.

The command used was:

\begin{verbatim}
cd ~/ansible-nac
ansible-playbook -i inventory/hosts.ini playbooks/nac_validation.yml -K -e remediate=true
\end{verbatim}

The authoritative execution output is reproduced below.

% ---------------------------------------------------------------------
% NOTE: This Ansible console output is included because it shows the
% complete task path for the remediation-enabled run and preserves the
% successful restart and post-check behavior for both managed services.
% ---------------------------------------------------------------------

\begin{verbatim}
PLAY [Validate NAC services and optionally remediate] **************************

TASK [Collect basic hostname] **************************************************
ok: [radius-splunk]
ok: [access-node]

TASK [Check hostapd-wired service on access-node] ******************************
skipping: [radius-splunk]
ok: [access-node]

TASK [Check freeradius service on radius-splunk] *******************************
skipping: [access-node]
ok: [radius-splunk]

TASK [Restart hostapd-wired service when remediation is enabled] ***************
skipping: [radius-splunk]
changed: [access-node]

TASK [Restart freeradius service when remediation is enabled] ******************
skipping: [access-node]
changed: [radius-splunk]

TASK [Re-check hostapd-wired service on access-node] ***************************
skipping: [radius-splunk]
ok: [access-node]

TASK [Re-check freeradius service on radius-splunk] ****************************
skipping: [access-node]
ok: [radius-splunk]

TASK [Collect recent hostapd logs] *********************************************
skipping: [radius-splunk]
ok: [access-node]

TASK [Collect recent freeradius service logs] **********************************
skipping: [access-node]
ok: [radius-splunk]

TASK [Collect recent freeradius application log] *******************************
skipping: [access-node]
ok: [radius-splunk]

TASK [Save hostapd pre-check status locally] ***********************************
skipping: [radius-splunk]
changed: [access-node -> localhost]

TASK [Save hostapd post-check status locally] **********************************
skipping: [radius-splunk]
changed: [access-node -> localhost]

TASK [Save hostapd restart result locally] *************************************
skipping: [radius-splunk]
ok: [access-node -> localhost]

TASK [Save hostapd logs locally] ***********************************************
skipping: [radius-splunk]
changed: [access-node -> localhost]

TASK [Save freeradius pre-check status locally] ********************************
skipping: [access-node]
changed: [radius-splunk -> localhost]

TASK [Save freeradius post-check status locally] *******************************
skipping: [access-node]
changed: [radius-splunk -> localhost]

TASK [Save freeradius restart result locally] **********************************
skipping: [access-node]
ok: [radius-splunk -> localhost]

TASK [Save freeradius service logs locally] ************************************
skipping: [access-node]
changed: [radius-splunk -> localhost]

TASK [Save freeradius application log locally] *********************************
skipping: [access-node]
changed: [radius-splunk -> localhost]

PLAY RECAP *********************************************************************
access-node                : ok=9    changed=4    unreachable=0    failed=0    skipped=10   
        rescued=0    ignored=0
radius-splunk              : ok=11   changed=5    unreachable=0    failed=0    skipped=8    
        rescued=0    ignored=0
\end{verbatim}

This run confirms that the automation can perform controlled remediation without abandoning its evidence-preservation role. The restart tasks executed successfully on both nodes:

\begin{itemize}
    \item \texttt{hostapd-wired.service} on \texttt{access-node};
    \item \texttt{freeradius.service} on \texttt{radius-splunk}.
\end{itemize}

The post-check tasks then completed successfully, which is essential. In a production-style workflow, the restart itself is not the success condition; the success condition is a successful post-remediation validation of the service state.

% ---------------------------------------------------------------------
% SUBSECTION 8.5: ARTIFACT PERSISTENCE AND EVIDENCE QUALITY
% ---------------------------------------------------------------------
% This subsection documents the final post-run artifact state and
% evaluates the practical usefulness of the generated files as local
% evidence for verification, troubleshooting, and appendix support.
% ---------------------------------------------------------------------

\subsection{8.5 Artifact Persistence and Evidence Quality}
\label{sec:artifact-persistence-and-evidence-quality}

After the validation and remediation runs, the artifact directory contained the following authoritative post-run state:

\begin{verbatim}
radius-splunk: ls -lh ~/ansible-nac/artifacts
total 44K
-rw-r--r-- 1 nac-lab nac-lab 3.5K Mar 24 12:23 access-node_hostapd_logs.txt
-rw-r--r-- 1 nac-lab nac-lab   25 Mar 24 12:43 access-node_hostapd_restart_result.txt
-rw-r--r-- 1 nac-lab nac-lab 1.4K Mar 24 12:43 access-node_hostapd_status_after.txt
-rw-r--r-- 1 nac-lab nac-lab 1.4K Mar 24 12:43 access-node_hostapd_status_before.txt
-rw-r--r-- 1 root    root    1.4K Mar 23 19:30 access-node_hostapd_status.txt
-rw-r--r-- 1 nac-lab nac-lab 4.0K Mar 24 12:23 radius-splunk_freeradius_journal.txt
-rw-r--r-- 1 nac-lab nac-lab 3.1K Mar 24 12:23 radius-splunk_freeradius_radiuslog.txt
-rw-r--r-- 1 nac-lab nac-lab   25 Mar 24 12:43 radius-splunk_freeradius_restart_result.txt
-rw-r--r-- 1 nac-lab nac-lab 1.8K Mar 24 12:43 radius-splunk_freeradius_status_after.txt
-rw-r--r-- 1 nac-lab nac-lab 1.8K Mar 24 12:43 radius-splunk_freeradius_status_before.txt
-rw-r--r-- 1 root    root    1.8K Mar 23 19:30 radius-splunk_freeradius_status.txt
\end{verbatim}

This is a strong outcome because it shows that the automation does not merely print to the console. It leaves behind local, reviewable evidence that can be used for:

\begin{itemize}
    \item post-execution verification;
    \item manual troubleshooting;
    \item appendix inclusion;
    \item change or incident record support.
\end{itemize}

Selected excerpts from the generated artifacts are reproduced below to demonstrate their practical value.

\begin{verbatim}
===== /home/nac-lab/ansible-nac/artifacts/access-node_hostapd_logs.txt =====
Mar 24 12:23:12 access-node hostapd[1927]: hostapd_interface_deinit_free(0x61845b677590)
Mar 24 12:23:12 access-node hostapd[1927]: hostapd_interface_deinit_free: num_bss=1 
        conf->num_bss=1
Mar 24 12:23:12 access-node hostapd[1927]: hostapd_interface_deinit(0x61845b677590)
Mar 24 12:23:12 access-node hostapd[1927]: ens38: interface state ENABLED->DISABLED
Mar 24 12:23:12 access-node hostapd[1927]: hostapd_bss_deinit: deinit bss ens38
Mar 24 12:23:12 access-node hostapd[1927]: ens38: Flushing old station entries
Mar 24 12:23:12 access-node hostapd[1927]: ens38: Deauthenticate all stations
Mar 24 12:23:12 access-node hostapd[1927]: ens38: AP-DISABLED

===== /home/nac-lab/ansible-nac/artifacts/access-node_hostapd_restart_result.txt =====
remediation not requested

===== /home/nac-lab/ansible-nac/artifacts/access-node_hostapd_status_after.txt =====
- hostapd-wired.service - hostapd wired 802.1X authenticator (lab)
     Loaded: loaded (/etc/systemd/system/hostapd-wired.service; enabled; preset: enabled)
     Active: active (running) since Tue 2026-03-24 12:23:12 UTC; 20min ago
   Main PID: 2054 (hostapd)
      Tasks: 1 (limit: 2206)
     Memory: 844.0K (peak: 1.0M)
        CPU: 31ms
     CGroup: /system.slice/hostapd-wired.service

===== /home/nac-lab/ansible-nac/artifacts/access-node_hostapd_status_before.txt =====
- hostapd-wired.service - hostapd wired 802.1X authenticator (lab)
     Loaded: loaded (/etc/systemd/system/hostapd-wired.service; enabled; preset: enabled)
     Active: active (running) since Tue 2026-03-24 12:23:12 UTC; 20min ago
   Main PID: 2054 (hostapd)
      Tasks: 1 (limit: 2206)
     Memory: 844.0K (peak: 1.0M)
        CPU: 31ms
     CGroup: /system.slice/hostapd-wired.service

===== /home/nac-lab/ansible-nac/artifacts/access-node_hostapd_status.txt =====
- hostapd-wired.service - hostapd wired 802.1X authenticator (lab)
     Loaded: loaded (/etc/systemd/system/hostapd-wired.service; enabled; preset: enabled)
     Active: active (running) since Mon 2026-03-23 18:51:07 UTC; 39min ago
   Main PID: 1017 (hostapd)
      Tasks: 1 (limit: 2206)
     Memory: 2.9M (peak: 3.1M)
        CPU: 51ms
     CGroup: /system.slice/hostapd.service

===== /home/nac-lab/ansible-nac/artifacts/radius-splunk_freeradius_journal.txt =====
Mar 24 12:22:39 radius-splunk freeradius[3253]: Compiling Post-Auth-Type Challenge for 
        attr Post-Auth-Type
Mar 24 12:22:39 radius-splunk freeradius[3253]: Compiling Post-Auth-Type Client-Lost for 
        attr Post-Auth-Type
Mar 24 12:22:39 radius-splunk freeradius[3253]: radiusd: #### Skipping IP addresses and 
        Ports ####
Mar 24 12:22:39 radius-splunk freeradius[3253]: Configuration appears to be OK
Mar 24 12:22:39 radius-splunk systemd[1]: Started freeradius.service - FreeRADIUS 
        multi-protocol policy server.
Mar 24 12:23:12 radius-splunk systemd[1]: Stopping freeradius.service - FreeRADIUS 
        multi-protocol policy server...
Mar 24 12:23:12 radius-splunk systemd[1]: freeradius.service: Deactivated successfully.
Mar 24 12:23:12 radius-splunk systemd[1]: Stopped freeradius.service - FreeRADIUS 
        multi-protocol policy server.

===== /home/nac-lab/ansible-nac/artifacts/radius-splunk_freeradius_radiuslog.txt =====
Tue Mar 24 11:34:18 2026 : Info: Loaded virtual server <default>
Tue Mar 24 11:34:18 2026 : Warning: Ignoring "sql" (see raddb/mods-available/README.rst)
Tue Mar 24 11:34:18 2026 : Warning: Ignoring "ldap" (see raddb/mods-available/README.rst)
Tue Mar 24 11:34:18 2026 : Info:  # Skipping contents of 'if' as it is always 'false' --
        /etc/freeradius/3.0/sites-enabled/inner-tunnel:366
Tue Mar 24 11:34:18 2026 : Info: Loaded virtual server inner-tunnel
Tue Mar 24 11:34:18 2026 : Info: Loaded virtual server default
Tue Mar 24 11:34:18 2026 : Info: Ready to process requests
Tue Mar 24 11:57:46 2026 : Info: Signalled to terminate

===== /home/nac-lab/ansible-nac/artifacts/radius-splunk_freeradius_restart_result.txt =====
remediation not requested

===== /home/nac-lab/ansible-nac/artifacts/radius-splunk_freeradius_status_after.txt =====
- freeradius.service - FreeRADIUS multi-protocol policy server
     Loaded: loaded (/usr/lib/systemd/system/freeradius.service; enabled; preset: enabled)
     Active: active (running) since Tue 2026-03-24 12:23:12 UTC; 20min ago
       Docs: man:radiusd(8)
             man:radiusd.conf(5)
             http://wiki.freeradius.org/
             http://networkradius.com/doc/

===== /home/nac-lab/ansible-nac/artifacts/radius-splunk_freeradius_status_before.txt =====
- freeradius.service - FreeRADIUS multi-protocol policy server
     Loaded: loaded (/usr/lib/systemd/system/freeradius.service; enabled; preset: enabled)
     Active: active (running) since Tue 2026-03-24 12:23:12 UTC; 20min ago
       Docs: man:radiusd(8)
             man:radiusd.conf(5)
             http://wiki.freeradius.org/
             http://networkradius.com/doc/

===== /home/nac-lab/ansible-nac/artifacts/radius-splunk_freeradius_status.txt =====
- freeradius.service - FreeRADIUS multi-protocol policy server
     Loaded: loaded (/usr/lib/systemd/system/freeradius.service; enabled; preset: enabled)
     Active: active (running) since Mon 2026-03-23 18:52:42 UTC; 37min ago
       Docs: man:radiusd(8)
             man:radiusd.conf(5)
             http://wiki.freeradius.org/
             http://networkradius.com/doc/
\end{verbatim}

Several observations are worth recording:

\begin{itemize}
    \item the artifact files are human-readable and immediately useful without additional tooling;
    \item the \texttt{hostapd} journal excerpt clearly captures a service deinitialization path consistent with a controlled restart;
    \item the FreeRADIUS journal excerpt captures a stop/start lifecycle and successful service initialization;
    \item the service status files provide simple before/after validation anchors;
    \item the legacy \texttt{*\_status.txt} files remain present as earlier-generation artifacts, while the newer \texttt{*\_status\_before.txt} and \texttt{*\_status\_after.txt} files reflect the final workflow.
\end{itemize}

One nuance should be noted: the restart result files currently contain \texttt{remediation not requested}. This reflects the playbook's current use of \texttt{stdout} capture for the restart tasks. Because \texttt{systemctl restart} does not normally emit meaningful stdout, the fallback string remains in place even during a successful remediation-enabled run. This does \emph{not} invalidate the workflow, because the authoritative proof of remediation is instead found in:

\begin{itemize}
    \item the Ansible console output showing the restart tasks as \texttt{changed};
    \item the post-check service status files;
    \item the journal excerpts showing stop/start lifecycle behavior.
\end{itemize}

This limitation is acceptable in the current lab and will be addressed explicitly later in the limitations section.

% ---------------------------------------------------------------------
% SUBSECTION 8.6: PYTHON INCIDENT PARSER ROLE
% ---------------------------------------------------------------------
% This subsection documents the role of the Python incident parser as a
% narrow evidence-summary tool that complements the Ansible workflow by
% synthesizing the preserved incident artifacts.
% ---------------------------------------------------------------------

\subsection{8.6 Python Incident Parser Role}
\label{sec:python-incident-parser-role}

The Ansible playbook is not the only automation artifact in the environment. A separate Python utility, \texttt{nac\_incident\_parser.py}, is used to summarize the incident evidence files collected during the controlled failure-and-remediation scenario.

At the time of final documentation, the parser executed successfully and produced the following authoritative output:

\begin{verbatim}
radius-splunk: python3 ~/incident-evidence/nac_incident_parser.py
=== NAC Incident Evidence Summary ===

[1] Failed authentication simulation
- Client remained unauthorized: YES
- Client showed successful connection during failed test: NO
- hostapd saw authentication attempt: YES
- FreeRADIUS recorded failed auth / reject path: YES

[2] Remediation / successful re-test
- Client reached CTRL-EVENT-CONNECTED: YES
- Client reached Authorized state: YES
- hostapd showed successful authorization indicators: YES

[3] Incident conclusion
- SIMULATED NAC INCIDENT SUCCESSFULLY REPRODUCED AND REMEDIATED
\end{verbatim}

This parser is intentionally narrow in scope. It is not a general log analytics engine and it does not replace direct review of raw evidence. Instead, it provides a lightweight operational summary that is useful for:

\begin{itemize}
    \item quickly confirming that the expected failure and recovery markers are present;
    \item reducing manual re-check time during repeated test cycles;
    \item supporting a validation narrative that remains grounded in raw evidence files.
\end{itemize}

The parser therefore complements the Ansible workflow rather than competing with it:

\begin{itemize}
    \item Ansible validates service state and captures artifacts;
    \item the Python parser summarizes the incident evidence set;
    \item both workflows remain transparent and auditable because their underlying inputs are preserved.
\end{itemize}

The full parser source is reproduced in Appendix A (see Section~\ref{sec:appendix-a-python-nac-incident-parser}).

% ---------------------------------------------------------------------
% SUBSECTION 8.7: OPERATIONAL ASSESSMENT
% ---------------------------------------------------------------------
% This subsection evaluates the automation layer as an operational
% control pattern and explains how the final workflow supports repeatable
% validation, bounded remediation, and retained evidence collection.
% ---------------------------------------------------------------------

\subsection{8.7 Operational Assessment}
\label{sec:operational-assessment}

From an operational perspective, the automation layer succeeds because it does not overreach. It implements a realistic control pattern for a lab of this size:

\begin{enumerate}
    \item validate the relevant services;
    \item optionally restart them when instructed;
    \item immediately re-check state;
    \item persist artifacts locally;
    \item correlate with separate incident evidence.
\end{enumerate}

During final validation:

\begin{itemize}
    \item both playbook runs completed successfully without manual intervention during execution;
    \item SSH connectivity from \texttt{radius-splunk} to \texttt{access-node} remained stable during successful runs;
    \item both the validation-only and remediation-enabled workflows completed without unreachable or failed hosts.
\end{itemize}

This is intended to demonstrate:

\begin{itemize}
    \item operational thinking rather than one-time manual testing;
    \item repeatability across service validation cycles;
    \item controlled remediation rather than purely observational monitoring;
    \item artifact preservation suitable for documentation and audit-style review.
\end{itemize}

At the same time, the automation remains intentionally modest. The purpose of the automation is to demonstrate sound operational design in a constrained environment, not to claim parity with enterprise orchestration platforms.

% ---------------------------------------------------------------------
% SUBSECTION 8.8: SECTION CONCLUSION
% ---------------------------------------------------------------------
% This subsection closes the automation section by summarizing the role
% of the playbook and parser in strengthening the lab as a repeatable
% and operationally credible NAC validation environment.
% ---------------------------------------------------------------------

\subsection{8.8 Section Conclusion}
\label{sec:section-8-conclusion}

The automation layer improves the robustness of the lab as an enterprise-aligned production artifact.

Without automation, the lab would still demonstrate a valid wired 802.1X design, but it would remain largely manual. With the addition of Ansible-based validation and optional remediation, plus the Python evidence parser, the lab demonstrates a more complete operational model:

\begin{itemize}
    \item services can be checked from a central node;
    \item remediation can be invoked deliberately rather than manually per host;
    \item artifacts are preserved locally for review;
    \item incident evidence can be summarized in a repeatable way;
    \item the automation remains transparent, bounded, and easy to audit.
\end{itemize}

This is the correct level of automation for the environment. It supports the lab's central goals without obscuring the underlying technical controls or replacing direct evidence review.

% =====================================================================
% SECTION 9: OBSERVABILITY AND LOGGING PIPELINE
% =====================================================================
% This section documents the local log sources, monitored-input
% configuration, service-specific evidence paths, and the role of
% curated incident capture files in the final lab.
% =====================================================================

\section{9. Observability and Logging Pipeline}
\label{sec:observability-and-logging-pipeline}

The lab includes a deliberately structured observability layer intended to support both steady-state validation and incident-focused troubleshooting. This layer is built from a combination of:

\begin{itemize}
    \item continuously generated local service and system logs;
    \item service-oriented journald output for the Linux-based authenticator;
    \item Splunk Universal Forwarder monitored-input configuration;
    \item manually curated incident evidence files generated during controlled failure and recovery testing.
\end{itemize}

The objective is not to claim a full centralized SIEM deployment, but rather to demonstrate that the environment has a coherent, technically meaningful evidence pipeline. In this lab, observability is treated as an engineering function in its own right: the system produces logs, those logs are mapped to identifiable sources, relevant paths are monitored, and incident-specific captures are preserved separately when the continuously generated logs alone are insufficient for a clean narrative.

% ---------------------------------------------------------------------
% SUBSECTION 9.1: OBSERVABILITY SCOPE AND POSITIONING
% ---------------------------------------------------------------------
% This subsection defines the scope of the observability layer and
% positions the Splunk component as a local monitored-input deployment
% rather than a full centralized SIEM architecture.
% ---------------------------------------------------------------------

\subsection{9.1 Observability Scope and Positioning}
\label{sec:observability-scope-and-positioning}

The Splunk component in this lab is deliberately scoped as a local monitored-input configuration without demonstrated remote indexer or search-head integration. That scope is both accurate and sufficient for the environment that was built.

This section therefore demonstrates the following, and only the following:

\begin{itemize}
    \item identification of the relevant local log sources;
    \item configuration of Splunk Universal Forwarder to monitor those sources;
    \item preservation of service and incident evidence in local files;
    \item a logging model that supports the final-state validation and incident sections of this document.
\end{itemize}

It does \emph{not} demonstrate:

\begin{itemize}
    \item centralized Splunk search or analytics;
    \item a remote indexer/search-head topology;
    \item dashboards, alerts, or correlation searches;
    \item enterprise retention, forwarding, or multi-tier ingestion architecture.
\end{itemize}

This scope boundary matters because it allows the document to make a strong but accurate claim: the lab is instrumented and observable, even though it is not a full SIEM deployment.

% ---------------------------------------------------------------------
% SUBSECTION 9.2: SPLUNK UNIVERSAL FORWARDER MONITORED INPUTS
% ---------------------------------------------------------------------
% This subsection documents the active Splunk Universal Forwarder state,
% the local monitored-input configuration, and the specific log sources
% selected for continuous monitoring on radius-splunk.
% ---------------------------------------------------------------------

\subsection{9.2 Splunk Universal Forwarder Monitored Inputs}
\label{sec:splunk-universal-forwarder-monitored-inputs}

The Splunk Universal Forwarder on \texttt{radius-splunk} is active in the final state:

\begin{verbatim}
radius-splunk: /opt/splunkforwarder/bin/splunk status
Warning: Attempting to revert the SPLUNK_HOME ownership
Warning: Executing "chown -R root /opt/splunkforwarder"
splunkd is running (PID: 1489)
splunk helpers are running (PIDs: 1490)
\end{verbatim}

The monitored-input configuration is maintained in a dedicated local app:

\begin{verbatim}
radius-splunk: /opt/splunkforwarder/etc/apps/lab_nac_inputs/local/inputs.conf
[monitor:///var/log/syslog]
disabled = false
index = main
sourcetype = syslog

[monitor:///var/log/auth.log]
disabled = false
index = main
sourcetype = linux_secure

[monitor:///var/log/freeradius/radius.log]
disabled = false
index = main
sourcetype = freeradius
\end{verbatim}

The app layout itself is minimal and clean:

\begin{verbatim}
radius-splunk: find /opt/splunkforwarder/etc/apps/lab_nac_inputs -maxdepth 3 -type f | sort
/opt/splunkforwarder/etc/apps/lab_nac_inputs/local/inputs.conf
\end{verbatim}

This is the correct level of Splunk configuration for the lab. The forwarder monitors three sources that are directly relevant to the final environment:

\begin{itemize}
    \item \texttt{/var/log/syslog} for general system and service lifecycle events;
    \item \texttt{/var/log/auth.log} for administrative access and privilege activity;
    \item \texttt{/var/log/freeradius/radius.log} for FreeRADIUS application-layer events.
\end{itemize}

Together, these monitored inputs provide coverage over:
- host operating-system state,
- administrative activity,
- and the primary authentication service itself.

The full UF input configuration is reproduced in Appendix A (see Section~\ref{sec:appendix-a-splunk-uf-inputs-conf}).

% ---------------------------------------------------------------------
% SUBSECTION 9.3: PRIMARY CONTINUOUS LOG SOURCES ON RADIUS-SPLUNK
% ---------------------------------------------------------------------
% This subsection documents the primary continuously generated log
% sources on radius-splunk and shows how system, administrative, and
% FreeRADIUS application logs contribute to the observability layer.
% ---------------------------------------------------------------------

\subsection{9.3 Primary Continuous Log Sources on \texttt{radius-splunk}}
\label{sec:primary-continuous-log-sources-on-radius-splunk}

The most important continuously generated log sources in the environment are located on \texttt{radius-splunk}. These include system log data, authentication/admin activity, and FreeRADIUS application output.

Their current filesystem metadata is summarized below.

\begin{verbatim}
radius-splunk log files
syslog:  -rw-r----- 1 syslog adm 2.7M Mar 24 13:30 /var/log/syslog
auth.log: -rw-r----- 1 syslog adm 233K Mar 24 13:25 /var/log/auth.log
radius.log: -rw-r--r-- 1 syslog adm 12K Mar 24 13:14 /var/log/freeradius/radius.log
\end{verbatim}

\subsubsection{9.3.1 System Log (\texttt{/var/log/syslog})}
\label{sec:system-log-syslog}

The system log provides host-level operational context such as service start/stop events, scheduled job activity, and general daemon behavior. A representative excerpt is reproduced below.

\begin{verbatim}
radius-splunk: tail -n 15 /var/log/syslog
2026-03-24T13:14:41.348438+00:00 radius-splunk multipathd[467]: sda: failed to get sgio 
        uid: No such file or directory
2026-03-24T13:14:41.348468+00:00 radius-splunk multipathd[467]: sda: no WWID in state 
        "undef
2026-03-24T13:14:41.348639+00:00 radius-splunk multipathd[467]: ", giving up
2026-03-24T13:14:41.348690+00:00 radius-splunk multipathd[467]: sda: check_path() failed, 
        removing
2026-03-24T13:15:01.629554+00:00 radius-splunk CRON[1579]: (root) CMD (command -v 
        debian-sa1 > /dev/null && debian-sa1 1 1)
2026-03-24T13:17:01.652202+00:00 radius-splunk CRON[1587]: (root) CMD (cd / && run-parts 
        --report /etc/cron.hourly)
2026-03-24T13:19:05.468094+00:00 radius-splunk systemd[1]: Starting 
        update-notifier-download.service - Download data for packages that failed at package 
        install time...
2026-03-24T13:19:05.678561+00:00 radius-splunk systemd[1]: update-notifier-download.service
        : Deactivated successfully.
2026-03-24T13:19:05.678792+00:00 radius-splunk systemd[1]: Finished 
        update-notifier-download.service - Download data for packages that failed at package 
        install time.
2026-03-24T13:19:31.934534+00:00 radius-splunk systemd[1376]: launchpadlib-cache-clean.
        service - Clean up old files in the Launchpadlib cache was skipped because of an
        unmet condition check (ConditionPathExists=/home/nac-lab/.launchpadlib/
        api.launchpad.net/cache).
\end{verbatim}

This excerpt is not included because each line is individually important. Rather, it demonstrates the character of the source: \texttt{syslog} is a general-purpose host operational log that complements more specialized sources such as \texttt{radius.log}.

\subsubsection{9.3.2 Administrative and Authentication Log (\texttt{/var/log/auth.log})}
\label{sec:administrative-and-authentication-log-auth-log}

The administrative and authentication log captures local privilege elevation activity and remote administrative access events. A representative excerpt is reproduced below.

\begin{verbatim}
radius-splunk: tail -n 15 /var/log/auth.log
2026-03-24T13:20:04.756710+00:00 radius-splunk sshd[1628]: pam_unix(sshd:session): session 
        opened for user nac-lab(uid=1000) by nac-lab(uid=0)
2026-03-24T13:20:04.762714+00:00 radius-splunk systemd-logind[816]: New session 5 of user 
        nac-lab.
2026-03-24T13:20:05.503276+00:00 radius-splunk sshd[1686]: Received disconnect from 
        192.168.195.1 port 41932:11: disconnected by user
2026-03-24T13:20:05.503525+00:00 radius-splunk sshd[1686]: Disconnected from user nac-lab 
        192.168.195.1 port 41932
2026-03-24T13:20:05.504875+00:00 radius-splunk sshd[1628]: pam_unix(sshd:session): session 
        closed for user nac-lab
2026-03-24T13:20:05.511847+00:00 radius-splunk systemd-logind[816]: Session 5 logged out. 
        Waiting for processes to exit.
2026-03-24T13:20:05.514296+00:00 radius-splunk systemd-logind[816]: Removed session 5.
2026-03-24T13:20:37.703552+00:00 radius-splunk sudo:  nac-lab : TTY=tty1 ; 
        PWD=/home/nac-lab ; USER=root ; COMMAND=/usr/bin/tail -n 15 /var/log/syslog
2026-03-24T13:20:37.704672+00:00 radius-splunk sudo: pam_unix(sudo:session): session opened 
        for user root(uid=0) by nac-lab(uid=1000)
2026-03-24T13:20:37.731117+00:00 radius-splunk sudo: pam_unix(sudo:session): session closed 
        for user root
2026-03-24T13:21:39.041208+00:00 radius-splunk sudo:  nac-lab : TTY=tty1 ; 
        PWD=/home/nac-lab ; USER=root ; COMMAND=/usr/bin/tail -n 15 /var/log/auth.log
2026-03-24T13:21:39.042808+00:00 radius-splunk sudo: pam_unix(sudo:session): session opened 
        for user root(uid=0) by nac-lab(uid=1000)
2026-03-24T13:21:39.068864+00:00 radius-splunk sudo: pam_unix(sudo:session): session closed 
        for user root
2026-03-24T13:22:03.513481+00:00 radius-splunk sudo:  nac-lab : TTY=tty1 ; 
        PWD=/home/nac-lab ; USER=root ; COMMAND=/usr/bin/tail -n 15 /var/log/auth.log
2026-03-24T13:22:03.514636+00:00 radius-splunk sudo: pam_unix(sudo:session): session opened 
        for user root(uid=0) by nac-lab(uid=1000)
\end{verbatim}

This source is especially useful in a lab of this type because it records operator activity, privilege escalation, and remote administrative access events that may be relevant during incident reconstruction or audit-style review.

\subsubsection{9.3.3 FreeRADIUS Application Log (\texttt{/var/log/freeradius/radius.log})}
\label{sec:freeradius-application-log-radius-log}

The most important continuously generated application log in the environment is the FreeRADIUS log. This file captures service initialization, readiness state, termination events, and other application-relevant messages.

Its metadata is as follows:

\begin{verbatim}
radius-splunk: ls -lh /var/log/freeradius/radius.log
-rw-r--r-- 1 freerad adm 12K Mar 24 13:14 /var/log/freeradius/radius.log
\end{verbatim}

A representative excerpt is reproduced below.

\begin{verbatim}
radius-splunk: tail -n 20 /var/log/freeradius/radius.log
Tue Mar 24 12:23:12 2026 : Info: Debug state unknown (cap_sys_ptrace capability not set)
Tue Mar 24 12:23:12 2026 : Info: systemd watchdog interval is 30.00 secs
Tue Mar 24 12:23:12 2026 : Info: Loaded virtual server <default>
Tue Mar 24 12:23:12 2026 : Warning: Ignoring "sql" (see raddb/mods-available/README.rst)
Tue Mar 24 12:23:12 2026 : Warning: Ignoring "ldap" (see raddb/mods-available/README.rst)
Tue Mar 24 12:23:12 2026 : Info:  # Skipping contents of 'if' as it is always 'false' --
        /etc/freeradius/3.0/sites-enabled/inner-tunnel:366
Tue Mar 24 12:23:12 2026 : Info: Loaded virtual server inner-tunnel
Tue Mar 24 12:23:12 2026 : Info: Loaded virtual server default
Tue Mar 24 12:23:12 2026 : Info: Ready to process requests
Tue Mar 24 12:53:42 2026 : Info: Signalled to terminate
Tue Mar 24 12:53:42 2026 : Info: Exiting normally
Tue Mar 24 13:14:05 2026 : Info: Debug state unknown (cap_sys_ptrace capability not set)
Tue Mar 24 13:14:05 2026 : Info: systemd watchdog interval is 30.00 secs
Tue Mar 24 13:14:05 2026 : Info: Loaded virtual server <default>
Tue Mar 24 13:14:05 2026 : Warning: Ignoring "sql" (see raddb/mods-available/README.rst)
Tue Mar 24 13:14:05 2026 : Warning: Ignoring "ldap" (see raddb/mods-available/README.rst)
Tue Mar 24 13:14:05 2026 : Info:  # Skipping contents of 'if' as it is always 'false' -- 
        /etc/freeradius/3.0/sites-enabled/inner-tunnel:366
Tue Mar 24 13:14:05 2026 : Info: Loaded virtual server inner-tunnel
Tue Mar 24 13:14:05 2026 : Info: Loaded virtual server default
Tue Mar 24 13:14:05 2026 : Info: Ready to process requests
\end{verbatim}

This source is particularly valuable because it exposes application-level readiness and service lifecycle behavior that are directly relevant to the authentication plane. In the context of the lab, it complements both the system logs and the incident-specific evidence files.

% ---------------------------------------------------------------------
% SUBSECTION 9.4: HOSTAPD OBSERVABILITY ON ACCESS-NODE
% ---------------------------------------------------------------------
% This subsection documents the observability path for hostapd on
% access-node and explains that its authoritative continuous evidence is
% exposed through the systemd service state and journal rather than a
% dedicated application log file.
% ---------------------------------------------------------------------

\subsection{9.4 \texttt{hostapd} Observability on \texttt{access-node}}
\label{sec:hostapd-observability-on-access-node}

Unlike FreeRADIUS, the Linux-based authenticator does not rely on a dedicated application log file under \texttt{/var/log}. Instead, its authoritative operational evidence is exposed through the managed systemd service and the journal.

The current service state is reproduced below.

\begin{verbatim}
access-node: systemctl status hostapd-wired.service --no-pager
Loaded: loaded (/etc/systemd/system/hostapd-wired.service; enabled; preset: enabled)
Active: active (running) since Tue 2026-03-24 04:20:41 UTC
Main PID: 975 (hostapd)

Recent lines:
Using interface ens38 with hwaddr <REDACTED> and ssid ""
ens38: RADIUS Authentication server 10.0.30.10:1812
RADIUS local address: 10.0.23.1:45833
ens38: interface state UNINITIALIZED->ENABLED
ens38: AP-ENABLED
ens38: Setup of interface done.
\end{verbatim}

The underlying unit file is located at:

\begin{verbatim}
access-node: systemctl show -p FragmentPath hostapd-wired.service
FragmentPath=/etc/systemd/system/hostapd-wired.service
\end{verbatim}

A representative journal excerpt is reproduced below.

% ---------------------------------------------------------------------
% NOTE: This journal excerpt is included because it preserves recent
% hostapd lifecycle and startup events from the authoritative continuous
% log source used by the authenticator in this environment.
% ---------------------------------------------------------------------

\begin{verbatim}
access-node: journalctl -u hostapd-wired.service -n 20 --no-pager
Mar 24 02:13:55 access-node hostapd[2398]: ens38: CTRL-EVENT-TERMINATING
Mar 24 02:13:55 access-node hostapd[2398]: hostapd_free_hapd_data(ens38)
Mar 24 02:13:55 access-node hostapd[2398]: hostapd_interface_deinit_free: 
        driver=0x565a2d747060 drv_priv=0x565a38444cd0 -> hapd_deinit
Mar 24 02:13:55 access-node hostapd[2398]: hostapd_interface_free(0x565a38442590)
Mar 24 02:13:55 access-node hostapd[2398]: hostapd_interface_free: free hapd 0x565a38443d40
Mar 24 02:13:55 access-node hostapd[2398]: hostapd_cleanup_iface(0x565a38442590)
Mar 24 02:13:55 access-node hostapd[2398]: hostapd_cleanup_iface_partial(0x565a38442590)
Mar 24 02:13:55 access-node hostapd[2398]: hostapd_cleanup_iface: free iface=0x565a38442590
Mar 24 02:13:55 access-node systemd[1]: hostapd-wired.service: Deactivated successfully.
Mar 24 02:13:55 access-node systemd[1]: Stopped hostapd-wired.service - hostapd wired 
        802.1X authenticator (lab).
-- Boot 4efbb5bbbad2458a94bc281c688f0c9b --
Mar 24 04:20:41 access-node systemd[1]: Started hostapd-wired.service - hostapd wired 
        802.1X authenticator (lab).
Mar 24 04:20:41 access-node hostapd[975]: random: getrandom() support available
Mar 24 04:20:41 access-node hostapd[975]: Configuration file: 
        /etc/hostapd/lab/hostapd-wired.conf
Mar 24 04:20:41 access-node hostapd[975]: eapol_version=2
Mar 24 04:20:41 access-node hostapd[975]: ctrl_interface_group=0
Mar 24 04:20:41 access-node hostapd[975]: Opening raw packet socket for ifindex 3
Mar 24 04:20:41 access-node hostapd[975]: BSS count 1, BSSID mask 00:00:00:00:00:00 
        (0 bits)
Mar 24 04:20:41 access-node hostapd[975]: ens38: IEEE 802.11 Fetching hardware 
        channel/rate support not supported.
Mar 24 04:20:41 access-node hostapd[975]: Completing interface initialization
...
Mar 24 04:20:41 access-node hostapd[975]: ens38: AP-ENABLED
Mar 24 04:20:41 access-node hostapd[975]: ens38: Setup of interface done.
Mar 24 04:20:41 access-node hostapd[975]: ctrl_iface not configured!
\end{verbatim}

No dedicated \texttt{hostapd} file was present under \texttt{/var/log} at the time of validation. This is an important architectural note because it clarifies that \texttt{journald} is the authoritative continuous log source for the Linux-based authenticator in this environment.

This distinction matters operationally:
- FreeRADIUS provides both service/journal context and a dedicated application log file;
- \texttt{hostapd} primarily provides service/journal context;
- both are still compatible with the incident evidence workflow because curated captures can be created from their runtime output when a more focused artifact is needed.

% ---------------------------------------------------------------------
% SUBSECTION 9.5: CURATED INCIDENT EVIDENCE AS A COMPLEMENTARY LAYER
% ---------------------------------------------------------------------
% This subsection documents the curated incident-evidence directory and
% explains how manually captured debug artifacts complement the baseline
% continuous logging sources during controlled failure and recovery work.
% ---------------------------------------------------------------------

\subsection{9.5 Curated Incident Evidence as a Complementary Layer}
\label{sec:curated-incident-evidence-as-a-complementary-layer}

The continuously generated logs described above are necessary, but they are not the only observability artifacts in the lab. The environment also includes a curated incident-evidence directory on \texttt{radius-splunk}:

\begin{verbatim}
radius-splunk: ls -lh ~/incident-evidence
total 92K
-rw-rw-r-- 1 nac-lab nac-lab  32K Mar 23 04:02 freeradius_failed_auth_incident.txt
-rw-rw-r-- 1 nac-lab nac-lab  11K Mar 23 04:01 hostapd_failed_auth_incident.txt
-rw-rw-r-- 1 nac-lab nac-lab  19K Mar 23 04:01 hostapd_success_retest_v2.txt
-rw-r--r-- 1 root    root    3.5K Mar 24 09:52 nac_incident_parser.py
-rw-rw-r-- 1 nac-lab nac-lab  535 Mar 23 04:28 nac_incident_parser_success.txt
-rw-rw-r-- 1 nac-lab nac-lab 7.3K Mar 23 03:59 wpa_supplicant_failed_auth_incident.txt
-rw-rw-r-- 1 nac-lab nac-lab 8.4K Mar 23 03:59 wpa_supplicant_success_retest_v2.txt
\end{verbatim}

The following files were explicitly confirmed as manually captured from debug or controlled test runs:

\begin{itemize}
    \item \texttt{wpa\_supplicant\_failed\_auth\_incident.txt};
    \item \texttt{wpa\_supplicant\_success\_retest\_v2.txt};
    \item \texttt{hostapd\_failed\_auth\_incident.txt};
    \item \texttt{hostapd\_success\_retest\_v2.txt};
    \item \texttt{freeradius\_failed\_auth\_incident.txt}.
\end{itemize}

This distinction should be made explicit:

\begin{itemize}
    \item the continuously generated system and service logs are the baseline observability layer;
    \item the incident-evidence files are curated captures created to preserve high-signal artifacts from controlled validation runs.
\end{itemize}

The incident evidence files therefore complement rather than replace the continuous logging sources. They are particularly useful when the goal is to preserve a deterministic failure or recovery narrative in a form that is easy to reference later in the document, the parser, and the appendix sections.

% ---------------------------------------------------------------------
% SUBSECTION 9.6: OBSERVABILITY VALUE TO THE LAB
% ---------------------------------------------------------------------
% This subsection explains how the combined logging, journald, monitored-
% input, and curated-evidence model supports the validation, incident,
% automation, and packet-analysis sections that follow.
% ---------------------------------------------------------------------

\subsection{9.6 Observability Value to the Lab}
\label{sec:observability-value-to-the-lab}

The observability design in this lab is modest in scale but strong in structure. It provides:

\begin{itemize}
    \item identifiable local log sources for system, administrative, and authentication behavior;
    \item a dedicated monitored-input configuration in Splunk Universal Forwarder;
    \item service-oriented journal visibility for the Linux authenticator;
    \item curated incident evidence files for deterministic reproduction and documentation.
\end{itemize}

This combination is what makes later sections of the document possible. The access-control validation section depends on service and routing evidence. The incident section depends on curated captures and service logs. The automation section depends on local evidence persistence. The packet-analysis section depends on protocol captures aligned with the same operational story.

In other words, the observability layer is not an afterthought. It is one of the elements that makes the lab function as a coherent engineering artifact rather than as a loose set of services.

% ---------------------------------------------------------------------
% SUBSECTION 9.7: SECTION CONCLUSION
% ---------------------------------------------------------------------
% This subsection closes the observability section by summarizing the
% logging model as a well-scoped local evidence pipeline with clearly
% stated bounds on what the Splunk component does and does not represent.
% ---------------------------------------------------------------------

\subsection{9.7 Section Conclusion}
\label{sec:section-9-conclusion}

The final lab demonstrates a well-scoped but credible logging and observability model.

It is accurate to say that the environment includes:

\begin{itemize}
    \item log source identification;
    \item log capture and monitored-input configuration;
    \item service-level journal visibility;
    \item evidence preservation for validation and incident analysis.
\end{itemize}

It would not be accurate to claim that the environment demonstrates centralized Splunk search or enterprise analytics. The document will therefore continue to describe the Splunk component precisely as a local UF-based monitored-input layer.

% =====================================================================
% SECTION 10: PACKET CAPTURE ANALYSIS (TCPDUMP + WIRESHARK)
% =====================================================================
% This section documents protocol-level validation of the simulated
% wired 802.1X incident and the successful remediation re-test.
% It correlates the client-facing EAPOL exchange with the upstream
% authenticator-to-RADIUS transaction so that the incident is validated
% not only through host logs and service state, but also through direct
% packet evidence at both the IEEE 802.1X and RADIUS layers.
% =====================================================================

\section{10. Packet Capture Analysis (tcpdump + Wireshark)}
\label{sec:packet-capture-analysis-tcpdump-and-wireshark}

Packet capture analysis was performed to validate the incident scenario at the protocol layer. Earlier sections established the operational facts of the incident through endpoint behavior, FreeRADIUS evidence, hostapd evidence, and GLPI ticketing. This section adds a stricter form of proof by showing that the failed and successful authentication paths are directly observable in packet capture.

This is important because it demonstrates that the incident was not inferred only from logs or service state. The captures show the actual IEEE~802.1X / EAPOL exchange on the client-facing segment and the corresponding RADIUS transaction on the upstream path. In effect, the packet evidence ties together the supplicant, the authenticator, and the AAA backend into one coherent protocol narrative.

The scope of the capture work was intentionally narrow. These were not broad traffic captures intended to inventory all host communications. They were short, purpose-built captures limited to the authentication window so that the relevant control-plane exchange would remain concise, readable, and suitable for publication.

% ---------------------------------------------------------------------
% SUBSECTION 10.1: CAPTURE METHODOLOGY
% ---------------------------------------------------------------------
% This subsection defines where the packet captures were taken, which
% interfaces were used, and why the selected observation points are
% sufficient to correlate the client-facing IEEE 802.1X exchange with
% the upstream RADIUS transaction.
% ---------------------------------------------------------------------

\subsection{10.1 Capture Methodology}
\label{sec:capture-methodology}

All packet captures for this section were taken on \texttt{access-node}. This was the appropriate capture point because \texttt{access-node} is the wired 802.1X authenticator and therefore sits at the boundary between the client-facing controlled port and the upstream routed path toward FreeRADIUS.

Two interfaces were used:

\begin{itemize}
    \item \texttt{ens38} --- wired 802.1X / EAPOL side toward \texttt{client-1}
    \item \texttt{ens39} --- routed upstream side toward \texttt{dist-router} and the RADIUS service path
\end{itemize}

This dual-observation approach makes it possible to correlate:
\begin{itemize}
    \item the client-facing authentication result seen as EAPOL / EAP messages, and
    \item the backend AAA decision seen as RADIUS request/challenge/accept/reject exchange.
\end{itemize}

That correlation is the central analytical value of the packet capture work in this lab.

% ---------------------------------------------------------------------
% SUBSECTION 10.2: CAPTURE INVENTORY
% ---------------------------------------------------------------------
% This subsection establishes the authoritative pcap inventory used by
% the packet-analysis section and records the retained capture files
% and their sizes on access-node.
% ---------------------------------------------------------------------

\subsection{10.2 Capture Inventory}
\label{sec:capture-inventory}

The authoritative packet capture inventory consists of the following four files:

\begin{itemize}
    \item \texttt{failed\_8021x\_ens38.pcap}
    \item \texttt{failed\_8021x\_ens39.pcap}
    \item \texttt{success\_8021x\_ens38.pcap}
    \item \texttt{success\_8021x\_ens39.pcap}
\end{itemize}

The pcap inventory on \texttt{access-node} was as follows:

\begin{verbatim}
access-node: ls -lh ~/pcaps
-rw-r--r-- 1 nac-lab nac-lab 385 Mar 23 20:23 failed_8021x_ens38.pcap
-rw-r--r-- 1 nac-lab nac-lab 738 Mar 23 20:23 failed_8021x_ens39.pcap
-rw-r--r-- 1 nac-lab nac-lab 385 Mar 23 20:29 success_8021x_ens38.pcap
-rw-r--r-- 1 nac-lab nac-lab 747 Mar 23 20:29 success_8021x_ens39.pcap
\end{verbatim}

The small file sizes are expected and desirable. They indicate that the captures were intentionally constrained to the short authentication test windows rather than left running as general packet traces.

% ---------------------------------------------------------------------
% SUBSECTION 10.3: CAPTURE PROCEDURE
% ---------------------------------------------------------------------
% This subsection records the authoritative tcpdump capture procedure
% and the controlled manual supplicant invocation used to create the
% failed and successful authentication windows analyzed in this section.
% ---------------------------------------------------------------------

\subsection{10.3 Capture Procedure}
\label{sec:capture-procedure}

The captures were obtained by starting two concurrent \texttt{tcpdump} processes on \texttt{access-node} and then forcing a controlled supplicant authentication attempt from \texttt{client-1}. The method used was as follows:

\begin{verbatim}
access-node:
sudo timeout 20s tcpdump -i ens38 -w /tmp/failed_8021x_ens38.pcap ether proto 0x888e &
sudo timeout 20s tcpdump -i ens39 -w /tmp/failed_8021x_ens39.pcap host 10.0.30.10 and 
        udp port 1812 &
sleep 2

client-1:
sudo systemctl stop wpa_supplicant
sudo timeout 15s wpa_supplicant -i ens33 -D wired -c 
        /etc/wpa_supplicant/lab/wired-8021x.conf -dd

access-node:
sudo timeout 20s tcpdump -i ens38 -w ~/pcaps/success_8021x_ens38.pcap ether proto 0x888e &
sudo timeout 20s tcpdump -i ens39 -w ~/pcaps/success_8021x_ens39.pcap host 10.0.30.10 and
        udp port 1812 &
sleep 2

client-1:
sudo systemctl stop wpa_supplicant
sudo timeout 15s wpa_supplicant -i ens33 -D wired -c 
        /etc/wpa_supplicant/lab/wired-8021x.conf -dd
\end{verbatim}

This method is appropriate for a controlled lab because it produces deterministic, low-noise captures focused only on the authentication event under investigation.

The authoritative capture commands and retained packet inventory will be reproduced in full in Appendix~A (see Section~\ref{sec:appendix-b-packet-capture-methodology-and-tcpdump-commands} and Section~\ref{sec:appendix-b-pcap-file-inventory}).

% ---------------------------------------------------------------------
% SUBSECTION 10.4: FAILED AUTHENTICATION - CLIENT-FACING EAPOL ANALYSIS
% ---------------------------------------------------------------------
% This subsection analyzes the failed client-facing packet capture on
% the controlled port and establishes that the endpoint participated in
% a complete IEEE 802.1X / EAP exchange that terminated explicitly with
% EAP Failure rather than link loss or absent negotiation.
% ---------------------------------------------------------------------

\subsection{10.4 Failed Authentication: Client-Facing EAPOL Analysis}
\label{sec:failed-authentication-client-facing-eapol-analysis}

The failed client-facing capture corresponds to:

\begin{verbatim}
failed_8021x_ens38.pcap
\end{verbatim}

This capture was taken on \texttt{ens38}, the client-facing controlled segment between \texttt{client-1} and the Linux-based authenticator. In Wireshark, the failed capture shows a complete 802.1X / EAP exchange that progresses through the expected early stages but terminates unsuccessfully.

The visible sequence is:

\begin{enumerate}
    \item EAPOL Start
    \item EAP Request, Identity
    \item EAP Response, Identity
    \item EAP Request, MD5-Challenge
    \item EAP Response, MD5-Challenge
    \item EAP Failure
\end{enumerate}

This sequence is significant because it rules out several weaker explanations for the failure. The client did not fail due to lack of link, lack of EAPOL negotiation, or absence of authenticator response. The control-plane exchange clearly occurred and advanced into the MD5 challenge/response stage. The terminal outcome was explicitly \texttt{EAP Failure}.

Accordingly, the client-facing failed capture supports the same conclusion already reached through supplicant logs: the endpoint remained unauthorized because the authentication exchange completed with a negative result, not because the exchange never occurred.

\begin{figure}[H]
    \centering
    \includegraphics[width=\textwidth]{assets/screenshots/Wireshark_failed_8021x_ens38_pcap.png}
    \caption{Failed wired 802.1X authentication on the client-facing EAPOL segment}
    \label{sec:fig-failed-8021x-ens38}
\end{figure}

% ---------------------------------------------------------------------
% SUBSECTION 10.5: FAILED AUTHENTICATION - AUTHENTICATOR-TO-RADIUS ANALYSIS
% ---------------------------------------------------------------------
% This subsection analyzes the failed upstream RADIUS capture and
% demonstrates that the backend AAA path completed normally through
% Access-Challenge and terminated explicitly with Access-Reject,
% confirming a negative authorization decision from FreeRADIUS.
% ---------------------------------------------------------------------

\subsection{10.5 Failed Authentication: Authenticator-to-RADIUS Analysis}
\label{sec:failed-authentication-authenticator-to-radius-analysis}

The failed upstream RADIUS-side capture corresponds to:

\begin{verbatim}
failed_8021x_ens39.pcap
\end{verbatim}

This capture was taken on \texttt{ens39}, the routed side of \texttt{access-node} toward \texttt{dist-router} and the FreeRADIUS server. In Wireshark, the failed capture shows the corresponding backend AAA exchange that produced the client-facing failure observed on \texttt{ens38}.

The visible RADIUS transaction sequence is:

\begin{enumerate}
    \item Access-Request
    \item Access-Challenge
    \item Access-Request
    \item Access-Reject
\end{enumerate}

In the packet details for the terminal RADIUS response, the final \texttt{Access-Reject} carries an embedded EAP payload indicating \texttt{EAP Failure}. This is the strongest backend confirmation of the failed test case because it proves that the RADIUS server was reachable, the challenge-based workflow occurred normally, and the final authorization decision was explicitly negative.

This point matters analytically. It confirms that the failed incident was not caused by upstream transport loss or a silent AAA outage. The RADIUS service participated fully in the authentication exchange and returned an authoritative reject decision, which the client-facing side then reflected as \texttt{EAP Failure}.

\begin{figure}[H]
    \centering
    \includegraphics[width=\textwidth]{assets/screenshots/Wireshark_failed_8021x_ens39_pcap.png}
    \caption{Failed wired 802.1X authentication on the authenticator-to-RADIUS path}
    \label{sec:fig-failed-8021x-ens39}
\end{figure}

% ---------------------------------------------------------------------
% SUBSECTION 10.6: SUCCESSFUL REMEDIATION RE-TEST - CLIENT-FACING EAPOL ANALYSIS
% ---------------------------------------------------------------------
% This subsection analyzes the successful post-remediation client-facing
% capture and establishes that the endpoint completed the same general
% IEEE 802.1X workflow as the failed case but terminated with EAP
% Success after credential restoration.
% ---------------------------------------------------------------------

\subsection{10.6 Successful Remediation Re-Test: Client-Facing EAPOL Analysis}
\label{sec:successful-remediation-re-test-client-facing-eapol-analysis}

The successful client-facing remediation re-test capture corresponds to:

\begin{verbatim}
success_8021x_ens38.pcap
\end{verbatim}

This capture was also taken on \texttt{ens38}. In Wireshark, the successful capture shows the same general 802.1X / EAP structure as the failed case, but with a different terminal authorization result.

The visible sequence is:

\begin{enumerate}
    \item EAPOL Start
    \item EAP Request, Identity
    \item EAP Response, Identity
    \item EAP Request, MD5-Challenge
    \item EAP Response, MD5-Challenge
    \item EAP Success
\end{enumerate}

This is an important contrast with the failed capture. The early control-plane stages are substantially the same, which means the remediation did not alter the existence of the authentication workflow itself. Instead, it altered the final authorization result after the known-good credential state was restored.

The successful client-facing capture therefore confirms that the remediated endpoint completed the exchange and received an explicit \texttt{EAP Success}, which is the protocol-level counterpart to the authorized endpoint state observed in the supplicant and hostapd evidence.

\begin{figure}[H]
    \centering
    \includegraphics[width=\textwidth]{assets/screenshots/Wireshark_success_8021x_ens38_pcap.png}
    \caption{Successful wired 802.1X authentication on the client-facing EAPOL segment}
    \label{sec:fig-success-8021x-ens38}
\end{figure}

% ---------------------------------------------------------------------
% SUBSECTION 10.7: SUCCESSFUL REMEDIATION RE-TEST - AUTHENTICATOR-TO-RADIUS ANALYSIS
% ---------------------------------------------------------------------
% This subsection analyzes the successful upstream RADIUS capture and
% demonstrates that the backend AAA workflow completed normally through
% Access-Challenge and terminated explicitly with Access-Accept,
% confirming restoration of successful authentication after remediation.
% ---------------------------------------------------------------------

\subsection{10.7 Successful Remediation Re-Test: Authenticator-to-RADIUS Analysis}
\label{sec:successful-remediation-re-test-authenticator-to-radius-analysis}

The successful upstream RADIUS-side remediation re-test capture corresponds to:

\begin{verbatim}
success_8021x_ens39.pcap
\end{verbatim}

This capture was taken on \texttt{ens39}. In Wireshark, it shows the backend RADIUS exchange that corresponds to the successful client-facing authentication result.

The visible RADIUS transaction sequence is:

\begin{enumerate}
    \item Access-Request
    \item Access-Challenge
    \item Access-Request
    \item Access-Accept
\end{enumerate}

In the packet details for the terminal response, the \texttt{Access-Accept} contains an embedded EAP payload indicating \texttt{EAP Success}. This directly correlates with the client-facing \texttt{EAP Success} on \texttt{ens38} and provides authoritative AAA confirmation that the remediation restored successful backend authentication.

This is the protocol-level counterpart to the operational conclusion already established in the incident section: once the correct credential state was restored, the authentication workflow did not merely proceed farther or behave differently at random. It terminated with a positive authorization decision at the AAA layer.

\begin{figure}[H]
    \centering
    \includegraphics[width=\textwidth]{assets/screenshots/Wireshark_success_8021x_ens39_pcap.png}
    \caption{Successful wired 802.1X authentication on the authenticator-to-RADIUS path}
    \label{sec:fig-success-8021x-ens39}
\end{figure}

% ---------------------------------------------------------------------
% SUBSECTION 10.8: COMPARATIVE PROTOCOL INTERPRETATION
% ---------------------------------------------------------------------
% This subsection compares the failed and successful capture sets and
% establishes that both tests followed the same general control-plane
% workflow but diverged only at the terminal authorization decision,
% which is the key analytical finding of the packet evidence.
% ---------------------------------------------------------------------

\subsection{10.8 Comparative Protocol Interpretation}
\label{sec:comparative-protocol-interpretation}

The most important protocol distinction between the failed and successful capture sets is that both scenarios show a valid and complete 802.1X / EAP control-plane exchange through the Identity and MD5 challenge/response stages, but they terminate differently at both observation points.

In the failed case:
\begin{itemize}
    \item the client-facing \texttt{ens38} capture ends in \texttt{EAP Failure};
    \item the routed \texttt{ens39} capture ends in \texttt{Access-Reject}, carrying an embedded EAP failure indication.
\end{itemize}

In the successful remediation case:
\begin{itemize}
    \item the client-facing \texttt{ens38} capture ends in \texttt{EAP Success};
    \item the routed \texttt{ens39} capture ends in \texttt{Access-Accept}, carrying an embedded EAP success indication.
\end{itemize}

This is analytically important because it demonstrates that the failed test case was not caused by absence of EAPOL negotiation, lack of client participation, or general transport failure between the authenticator and FreeRADIUS. In both cases, the supplicant, authenticator, and AAA service all participated in a normal challenge-based authentication workflow. The decisive difference is the terminal policy result.

The captures therefore visually reinforce the distinction between the pre-authentication control-plane exchange and the final authorization outcome. Both failed and successful tests show substantially similar early negotiation behavior, but they diverge at the terminal decision point where authentication either remains unauthorized or transitions to authorized state.

% ---------------------------------------------------------------------
% SUBSECTION 10.9: PUBLICATION AND REDACTION GUIDANCE
% ---------------------------------------------------------------------
% This subsection records the publication policy for packet
% screenshots and identifies the specific fields visible in the current
% Wireshark screenshots that must be redacted before public release.
% ---------------------------------------------------------------------

\subsection{10.9 Publication and Redaction Guidance}
\label{sec:publication-and-redaction-guidance}

The publication policy for packet screenshots is:

\begin{itemize}
    \item private IP addresses may remain visible;
    \item MAC addresses must be redacted;
    \item passwords, shared secrets, and sensitive authentication blobs must be redacted;
    \item packet payload details that unnecessarily expose sensitive values should be redacted.
\end{itemize}

The following items should be redacted in the publication-safe versions:

\begin{itemize}
    \item all visible MAC addresses in the packet list pane;
    \item all visible MAC addresses in the packet details pane;
    \item MAC addresses exposed again in the hex pane;
    \item RADIUS \texttt{Authenticator} values where shown;
    \item \texttt{Message-Authenticator} values where shown;
\end{itemize}

% ---------------------------------------------------------------------
% SUBSECTION 10.10: SECTION CONCLUSION
% ---------------------------------------------------------------------
% This subsection closes the packet-analysis section by summarizing the
% evidentiary significance of the failed and successful capture pairs
% and restating why the protocol-layer validation materially strengthens
% the technical credibility of the lab.
% ---------------------------------------------------------------------

\subsection{10.11 Section Conclusion}
\label{sec:section-10-conclusion}

The packet capture analysis materially strengthens the technical credibility of the lab. It shows that the incident and remediation were not validated solely through host logs, service states, or application-level reachability tests. They were also validated through deterministic protocol transitions at both the IEEE~802.1X and RADIUS layers.

The failed capture set demonstrates a complete challenge-based authentication workflow ending in \texttt{EAP Failure} and \texttt{Access-Reject}. The successful remediation capture set demonstrates the same workflow ending in \texttt{EAP Success} and \texttt{Access-Accept}. That contrast provides clean protocol-level confirmation that the simulated incident was a true authentication decision failure and that remediation restored normal authorization behavior.

The authoritative capture commands and packet inventory will be referenced again in the appendix material so that the packet-analysis section remains readable while still being fully supported by raw artifact traceability.

% =====================================================================
% SECTION 11: DESIGN CONSTRAINTS AND LIMITATIONS
% =====================================================================
% This section documents the intentional design boundaries of the lab so
% that the document accurately represents what was implemented, what was
% validated, and what was deliberately excluded. The purpose of the
% section is to prevent overstatement, preserve technical credibility,
% and distinguish the proof-of-concept objective from a claim of full
% enterprise production equivalence.
% =====================================================================

\section{11. Design Constraints and Limitations}
\label{sec:design-constraints-and-limitations}

The lab documented in this report is a deliberate proof-of-concept implementation intended to validate core NAC behaviors, incident handling workflow, and operational evidence collection in a constrained virtual environment. It is not presented as a one-to-one reproduction of a full enterprise campus access control deployment. This distinction is important and intentional.

The design succeeds at the objective for which it was built: it demonstrates credential-gated wired access control, protected and guest network segmentation, deterministic incident reproduction, structured remediation, lightweight automation, and packet-level protocol validation. However, it does so through a set of explicit design constraints chosen to keep the lab tractable, reproducible, and appropriate for portfolio documentation.

These limitations should not be interpreted as weaknesses in the proof-of-concept. On the contrary, making them explicit increases the credibility of the document because it clearly separates what was actually implemented and validated from what would exist in a larger or more production-like deployment.

% ---------------------------------------------------------------------
% SUBSECTION 11.1: PLATFORM REALISM BOUNDARIES
% ---------------------------------------------------------------------
% This subsection defines the most important platform-level constraints
% in the lab, including the use of Linux hostapd as a functional wired
% authenticator approximation, the role of pfSense in guest isolation,
% and the treatment of virtualization-based segmentation as a lab
% implementation mechanism rather than a production architecture claim.
% ---------------------------------------------------------------------

\subsection{11.1 Platform Realism Boundaries}
\label{sec:platform-realism-boundaries}

The most important platform constraint in the lab is that the wired 802.1X authenticator function is implemented using Linux \texttt{hostapd} in wired mode rather than a managed enterprise access switch. This design choice is fully appropriate for the proof-of-concept objective because it produces real IEEE~802.1X / EAPOL behavior and a real RADIUS-backed authorization workflow, but it should not be described as equivalent to a vendor campus switching platform.

In practical terms, the lab demonstrates credential-gated access control behavior rather than switch-native NAC feature breadth. It does not attempt to model enterprise access-switch capabilities such as dynamic VLAN assignment, downloadable ACLs, switch vendor posture integration, or centralized policy enforcement engines. The access decision in this lab is still meaningful and technically valid, but the implementation path is intentionally simplified.

Similarly, the guest network is enforced primarily through pfSense interface and rule policy rather than through dynamic 802.1X authorization state. This is an intentional architectural separation: the authenticated access path demonstrates wired 802.1X behavior, while the guest path demonstrates policy-based isolation and controlled internet access. Together they provide a credible functional access model without claiming to be a full enterprise NAC fabric.

Finally, the virtual network segmentation used throughout the lab should be understood as a lab isolation mechanism, not as a substitute for a production campus switching and security architecture. The VMware custom network design is useful because it provides deterministic path control and repeatable separation between segments, but it is not presented as equivalent to enterprise hardware switching, trunking, or controller-driven segmentation.

% ---------------------------------------------------------------------
% SUBSECTION 11.2: AUTHENTICATION MODEL LIMITATIONS
% ---------------------------------------------------------------------
% This subsection documents the limitations of the chosen EAP method and
% explicitly distinguishes the lab's simplified EAP-MD5 demonstration
% from modern enterprise authentication practices such as certificate-
% based EAP methods.
% ---------------------------------------------------------------------

\subsection{11.2 Authentication Model Limitations}
\label{sec:authentication-model-limitations}

The authentication workflow demonstrated in this lab is based on EAP-MD5 challenge/response, as confirmed both by the service logs and by the packet captures analyzed in Section~\ref{sec:packet-capture-analysis-tcpdump-and-wireshark}. This was an intentional choice made for demonstration simplicity and reproducibility in a small virtual lab.

That choice should not be interpreted as a modern production recommendation. EAP-MD5 is useful in this context because it provides a compact and easily observable challenge-based authentication flow that can be validated clearly in both logs and packet capture. It is especially useful for showing the exact distinction between a failed credential state and a successful post-remediation state. However, it does not provide the security properties expected of contemporary enterprise access control.

The lab does not demonstrate certificate-based EAP methods such as EAP-TLS, PEAP, or EAP-TTLS. It also does not attempt to model the certificate lifecycle, trust distribution, supplicant provisioning, or PKI operational dependencies that would normally accompany a more production-aligned 802.1X deployment. These omissions are intentional and should be understood as scope boundaries rather than gaps in execution.

% ---------------------------------------------------------------------
% SUBSECTION 11.3: OBSERVABILITY AND SIEM LIMITATIONS
% ---------------------------------------------------------------------
% This subsection defines the boundary between the implemented logging
% and evidence-collection model and a true centralized SIEM deployment.
% It clarifies that Splunk Universal Forwarder was configured locally,
% while incident validation relied on multiple complementary evidence
% sources rather than centralized analytics.
% ---------------------------------------------------------------------

\subsection{11.3 Observability and SIEM Limitations}
\label{sec:observability-and-siem-limitations}

The observability model in this lab is intentionally practical but limited in scope. As documented in Section~\ref{sec:observability-and-logging-pipeline}, Splunk is represented only by a Splunk Universal Forwarder with local monitored-input configuration. This is sufficient to demonstrate awareness of log source selection and forwarder-oriented ingestion configuration, but it is not a full Splunk search-head, indexer, or analytics deployment.

For that reason, the document does not claim centralized SIEM investigation capability. No centralized search workflow, correlation rules, dashboards, or alerting pipelines were demonstrated as part of the lab. Instead, incident validation relied on a combination of:

\begin{itemize}
    \item native service logs from \texttt{hostapd}, FreeRADIUS, and systemd/journald;
    \item manually captured debug artifacts preserved in the incident evidence directory;
    \item the GLPI incident record used as the operational ticketing system of record; and
    \item packet capture analysis using \texttt{tcpdump} and Wireshark.
\end{itemize}

This multi-source evidence model is entirely appropriate for the lab objective. It produces a defensible incident narrative and a strong audit trail without overstating the presence of centralized SIEM operations.

% ---------------------------------------------------------------------
% SUBSECTION 11.4: AUTOMATION SCOPE LIMITATIONS
% ---------------------------------------------------------------------
% This subsection constrains the claims that can be made about the
% Python and Ansible automation used in the lab by explicitly defining
% them as purpose-built, validation-oriented, and tightly coupled to the
% known lab environment rather than generalized production frameworks.
% ---------------------------------------------------------------------

\subsection{11.4 Automation Scope Limitations}
\label{sec:automation-scope-limitations}

The automation implemented in this lab is intentionally narrow in scope and should be described accordingly. As documented in Section~\ref{sec:automation-workflows}, the Ansible playbook is best characterized as a validation-first workflow with optional service remediation. It is designed to confirm service health, optionally restart key services, and preserve the resulting evidence locally. It is not intended to represent a full production change-management pipeline, orchestration framework, or large-scale infrastructure automation pattern.

The same principle applies to the Python parser. The parser is a purpose-built evidence summarizer that evaluates a known set of captured artifact files in fixed paths and produces a human-readable incident conclusion. It is valuable because it demonstrates structured evidence interpretation and repeatable operator validation, but it is not a generalized incident correlation engine or a reusable NAC analytics platform.

The tools are tightly aligned to the services, filenames, and operational checkpoints that were actually implemented and validated. The document should therefore present the automation as operationally useful and technically appropriate, while avoiding any implication that it is broader or more portable than it was designed to be.

% ---------------------------------------------------------------------
% SUBSECTION 11.5: GUI AND EVIDENCE SCOPE LIMITATIONS
% ---------------------------------------------------------------------
% This subsection records the intentional decision to keep GUI evidence
% narrowly scoped to GLPI and Wireshark, and to rely primarily on
% command-line and protocol artifacts rather than broad administrative
% UI walkthroughs or hypervisor screenshots.
% ---------------------------------------------------------------------

\subsection{11.5 GUI and Evidence Scope Limitations}
\label{sec:gui-and-evidence-scope-limitations}

The production document intentionally limits GUI evidence to two categories:

\begin{itemize}
    \item the GLPI incident record; and
    \item Wireshark packet-analysis screenshots.
\end{itemize}

This is an explicit documentation choice rather than an omission. The lab is primarily a command-line and protocol-driven implementation, and the strongest evidence is found in configuration files, service outputs, debug logs, packet traces, and deterministic test results. Broad GUI screenshot inventories would add visual bulk without materially improving the technical case.

For that reason, the document intentionally avoids dependence on hypervisor settings walkthroughs, appliance UI tours, or other screenshot-heavy administrative paths. The goal is to keep the evidence model focused on raw technical artifacts and protocol behavior while using GUI elements only where they provide uniquely valuable context, specifically:

\begin{itemize}
    \item GLPI as the human-readable incident management record; and
    \item Wireshark as the visual representation of protocol-layer success and failure.
\end{itemize}

% ---------------------------------------------------------------------
% SUBSECTION 11.6: OPERATIONAL SCOPE BOUNDARIES
% ---------------------------------------------------------------------
% This subsection defines the scale and operational limits of the lab,
% emphasizing that the implementation validates a single-endpoint NAC
% incident workflow and does not attempt to model redundancy, scale, or
% enterprise control-plane complexity beyond the proof-of-concept goal.
% ---------------------------------------------------------------------

\subsection{11.6 Operational Scope Boundaries}
\label{sec:operational-scope-boundaries}

Operationally, the lab is centered on single-endpoint incident validation using \texttt{client-1} as the authoritative authenticated endpoint and \texttt{guest-test-1} as the guest-segment validation endpoint. This is sufficient for the design objective because the purpose of the lab is to prove that the access model, failure scenario, remediation, and evidence chain are all coherent and repeatable. It is not intended to simulate multi-endpoint concurrency or campus-scale endpoint churn.

Accordingly, the lab does not attempt to model any of the following:

\begin{itemize}
    \item redundant or load-balanced AAA infrastructure;
    \item high availability or failover behavior;
    \item large-scale endpoint inventory or onboarding workflows;
    \item enterprise certificate lifecycle and PKI operations;
    \item centralized NAC policy engines or posture frameworks;
    \item vendor switch telemetry, downloadable ACLs, or dynamic policy distribution.
\end{itemize}

These exclusions are intentional design boundaries. They do not diminish the proof-of-concept objective, which is to validate a coherent and technically defensible NAC incident workflow in a virtualized environment. In fact, by avoiding artificial breadth, the lab remains focused on a smaller set of claims that are strongly supported by direct evidence.

% ---------------------------------------------------------------------
% SUBSECTION 11.7: CREDIBILITY THROUGH EXPLICIT CONSTRAINTS
% ---------------------------------------------------------------------
% This subsection explains why documenting limitations is itself a
% positive quality attribute of the report and frames the constraints as
% deliberate boundaries that improve accuracy and trustworthiness.
% ---------------------------------------------------------------------

\subsection{11.7 Credibility Through Explicit Constraints}
\label{sec:credibility-through-explicit-constraints}

The constraints documented in this section should be understood as a strength of the report rather than as a disclaimer of weakness. In technical portfolio work, credibility is improved when the implementation is described precisely and bounded accurately. Overstating what a lab represents is more damaging than acknowledging where simplifications were made.

This lab makes strong, defensible claims in the areas where direct evidence exists:

\begin{itemize}
    \item functional wired 802.1X authentication behavior;
    \item RADIUS-backed authorization decisions;
    \item guest versus protected path differentiation;
    \item deterministic incident reproduction and remediation;
    \item operational evidence preservation through logs, automation, GLPI, and packet capture.
\end{itemize}

By explicitly stating what the lab does not attempt to represent, the document becomes more trustworthy. The result is a stronger portfolio artifact because the scope is narrow enough to be believable and broad enough to demonstrate meaningful engineering and operational skill.

% ---------------------------------------------------------------------
% SUBSECTION 11.8: SECTION CONCLUSION
% ---------------------------------------------------------------------
% This subsection closes the limitations section by restating that the
% lab should be read as a disciplined proof-of-concept with intentional
% scope boundaries, not as a claim of full production equivalence.
% ---------------------------------------------------------------------

\subsection{11.8 Section Conclusion}
\label{sec:section-11-conclusion}

This lab should be read as a disciplined proof-of-concept implementation with intentional scope boundaries. It successfully demonstrates the targeted behaviors it set out to validate: wired 802.1X access control, protected and guest path differentiation, incident reproduction, remediation, automation-assisted validation, and packet-level protocol corroboration.

It should not be interpreted as a full production NAC architecture or as a complete enterprise security platform. The design choices documented here were made deliberately to keep the lab reproducible, explainable, and strongly evidenced. Those constraints do not weaken the value of the implementation. They define it accurately.

The next section will use this bounded interpretation as the basis for a final conclusion and for a short discussion of logical future expansion paths that would extend the lab while preserving the same evidence-driven methodology.

% =====================================================================
% SECTION 12: CONCLUSION AND FUTURE EXPANSION OPPORTUNITIES
% =====================================================================
% This section concludes the document by synthesizing the technical
% outcomes of the lab, reaffirming the scope-bounded success criteria,
% and identifying realistic future expansion opportunities that would
% extend the proof-of-concept without overstating what was actually
% implemented in the current build.
% =====================================================================

\section{12. Conclusion and Future Expansion Opportunities}
\label{sec:conclusion-and-future-expansion-opportunities}

The lab documented in this report successfully met its original objective: to implement and validate a multi-segment network access control proof-of-concept that demonstrates wired 802.1X authentication, RADIUS-backed authorization, protected and guest path differentiation, incident simulation and remediation, lightweight automation, and packet-level protocol corroboration. Within the scope defined by the design, each major component was not only configured but also validated through direct operational evidence.

The authenticated path was shown to function through successful 802.1X supplicant behavior, hostapd authenticator processing, and FreeRADIUS acceptance. The guest path was shown to function through DHCP-based onboarding, gateway reachability, continued internet access, and enforced isolation from protected internal segments. The incident workflow further strengthened the implementation by proving that the environment could intentionally reproduce an authentication failure, preserve the resulting evidence, document the event in GLPI, restore correct configuration, and verify successful recovery. This moved the lab beyond simple service installation and into a more operationally credible validation model.

Equally important, the lab demonstrated end-to-end engineering continuity. The document does not merely show isolated services; it shows a coherent chain from design and segmentation through implementation, validation, troubleshooting, remediation, automation-assisted verification, and evidence preservation. That continuity is the strongest portfolio value of the project. It shows that the environment was built with enough rigor to support not only successful outcomes, but also controlled failure analysis and structured recovery.

As documented in Section~\ref{sec:design-constraints-and-limitations}, the lab is intentionally bounded and does not claim full enterprise production equivalence. That is a strength rather than a weakness. By explicitly limiting the claims to what was actually implemented and verified, the document remains technically credible and avoids overstating the realism of the environment. The result is a more trustworthy proof-of-concept and a stronger professional artifact.

% ---------------------------------------------------------------------
% SUBSECTION 12.1: SUMMARY OF ACCOMPLISHMENTS
% ---------------------------------------------------------------------
% This subsection restates the major validated accomplishments of the
% lab in concise form so that the conclusion clearly ties the final
% assessment back to the core implementation and evidence sections.
% ---------------------------------------------------------------------

\subsection{12.1 Summary of Accomplishments}
\label{sec:summary-of-accomplishments}

At a practical level, the lab achieved the following validated outcomes:

\begin{itemize}
    \item successful implementation of a multi-segment virtual network with distinct authenticated, guest, and service paths;
    \item successful wired 802.1X authentication flow using \texttt{wpa\_supplicant}, \texttt{hostapd} in wired mode, and FreeRADIUS;
    \item successful demonstration of guest-network DHCP assignment and enforced isolation from protected internal segments;
    \item successful reproduction of a realistic NAC-style authentication failure caused by incorrect endpoint credentials;
    \item successful remediation of that failure and confirmation of restored authorized access;
    \item successful documentation of the incident in GLPI as an operational record;
    \item successful implementation of lightweight Ansible-based service validation with optional remediation behavior; and
    \item successful packet-level corroboration of both failed and successful authentication states using \texttt{tcpdump} and Wireshark.
\end{itemize}

Taken together, these outcomes demonstrate that the lab was not merely assembled, but systematically tested and evidenced. That distinction is central to the value of the project.

% ---------------------------------------------------------------------
% SUBSECTION 12.2: PROFESSIONAL AND PORTFOLIO VALUE
% ---------------------------------------------------------------------
% This subsection explains why the completed lab functions as a useful
% portfolio artifact by emphasizing the operational coherence, evidence-
% driven methodology, and disciplined scope control demonstrated across
% the report.
% ---------------------------------------------------------------------

\subsection{12.2 Professional and Portfolio Value}
\label{sec:professional-and-portfolio-value}

The completed lab is a useful and credible portfolio piece because it demonstrates practical infrastructure and security-adjacent engineering in a disciplined way. It shows that the environment was designed intentionally, implemented from first principles, validated with repeatable test procedures, and documented with a strong evidentiary chain.

From a professional standpoint, the lab is especially valuable because it combines multiple modes of technical work in one coherent project:

\begin{itemize}
    \item network segmentation and path design;
    \item AAA and access-control service configuration;
    \item endpoint supplicant testing;
    \item troubleshooting and incident-style reasoning;
    \item log and packet evidence interpretation;
    \item automation-assisted operational validation; and
    \item structured technical documentation.
\end{itemize}

This breadth is meaningful because it reflects how real infrastructure work is often evaluated: not by whether a single service starts, but by whether the operator can connect architecture, implementation, validation, and evidence into a coherent operational story. This document does that successfully.

% ---------------------------------------------------------------------
% SUBSECTION 12.3: FUTURE EXPANSION OPPORTUNITIES
% ---------------------------------------------------------------------
% This subsection identifies realistic next-step enhancements that would
% extend the lab while remaining clearly framed as future work rather
% than as features already implemented in the current environment.
% ---------------------------------------------------------------------

\subsection{12.3 Future Expansion Opportunities}
\label{sec:future-expansion-opportunities}

Although the current implementation is complete within its stated scope, it also provides a strong foundation for future expansion. Several next-step enhancements would extend the realism or operational depth of the lab while preserving the same evidence-driven methodology used throughout this report.

A natural first expansion would be to replace the current EAP-MD5 demonstration method with a more production-aligned EAP model such as EAP-TLS, PEAP, or EAP-TTLS. This would significantly increase realism by introducing stronger authentication semantics and a more modern enterprise access-control pattern, while also creating opportunities to document certificate-based supplicant behavior and more advanced troubleshooting.

A second useful expansion would be to explore more dynamic authorization behavior. Within the constraints of a virtual lab, this could include experiments that more closely approximate switch-driven policy changes, segmented access outcomes, or more granular authorization handling beyond the current binary authenticated-versus-guest distinction. The goal would not be to claim full enterprise feature parity, but to extend the lab toward richer policy demonstrations.

A third expansion path would be to deepen observability. Section~\ref{sec:observability-and-logging-pipeline} intentionally documents a local Splunk Universal Forwarder monitored-input configuration rather than a full centralized analytics stack. A future iteration could extend this into a more complete logging pipeline with centralized indexing, search, and event correlation, thereby improving the analytical side of the incident workflow.

The automation layer also provides room for meaningful growth. The current Ansible playbook and Python parser are intentionally narrow, which is appropriate for the present proof-of-concept. Future versions could add richer health checks, more structured output formats, broader artifact collection logic, or more generalized remediation patterns while still preserving the validation-first operating model established in Section~\ref{sec:automation-workflows}.

Scale realism could also be improved by introducing multiple authenticated endpoints or additional controlled failure scenarios. The current single-endpoint incident workflow is sufficient to prove the concept, but repeated validation across multiple clients or multiple failure classes would strengthen the repeatability argument and provide more varied operational evidence.

Finally, the GLPI workflow could be extended from its current manual but effective incident-record role into a more formalized evidence-management process. Future iterations could attach or reference more structured artifacts directly, impose a more consistent incident template, or align the ticketing workflow more closely with the operational outputs generated by the automation layer.

Each of these expansion opportunities is intentionally framed as future work. None of them are necessary to validate the success of the current build. Rather, they demonstrate that the lab has been constructed on a foundation strong enough to support meaningful iterative growth.

% ---------------------------------------------------------------------
% SUBSECTION 12.4: FINAL ASSESSMENT
% ---------------------------------------------------------------------
% This subsection provides the final evaluative judgment of the lab in
% concise form, closing the report with a scope-bounded but affirmative
% statement of technical success and professional usefulness.
% ---------------------------------------------------------------------

\subsection{12.4 Final Assessment}
\label{sec:final-assessment}

In final assessment, the lab should be considered successful both as a technical implementation and as a professional documentation artifact. It achieved the intended proof-of-concept objectives, produced direct evidence for each major claim, and maintained strong internal consistency between architecture, validation, incident handling, automation, and protocol analysis.

Just as importantly, the document avoids the common failure mode of overstatement. It does not claim enterprise completeness where none was implemented. Instead, it presents a bounded, defensible, and well-evidenced engineering exercise that demonstrates practical competence across several adjacent domains.

Within its stated scope, the lab provides a technically defensible and well-documented proof-of-concept for wired NAC validation.

% ---------------------------------------------------------------------
% SUBSECTION 12.5: TRANSITION TO APPENDICES
% ---------------------------------------------------------------------
% This subsection provides a brief transition from the main body to the
% appendices, clarifying that the appendices preserve the raw artifacts
% already referenced throughout the report rather than introducing new
% substantive claims.
% ---------------------------------------------------------------------

\subsection{12.5 Transition to Appendices}
\label{sec:transition-to-appendices}

The appendices that follow preserve the raw configurations, scripts, command outputs, and supporting artifacts referenced throughout the body of the report. They are included to support traceability, reproducibility, and evidentiary completeness, and should be read as the authoritative source for the full raw materials cited in earlier sections.

% =====================================================================
% APPENDIX A: RAW CONFIGURATIONS AND AUTOMATION ARTIFACTS
% =====================================================================
% This appendix preserves the authoritative raw configuration files,
% manually transcribed firewall policy artifacts, and automation source
% files referenced throughout the body of the report. Its purpose is to
% provide full-text traceability for the final implemented state of the
% lab without forcing the main body to carry every raw artifact inline.
% =====================================================================

\appendix

\section{Appendix A. Raw Configurations and Automation Artifacts}
\label{sec:appendix-a-raw-configurations-and-automation-artifacts}

This appendix preserves the authoritative raw configurations and automation artifacts referenced throughout the report. All artifacts in this appendix are treated as final-state materials unless otherwise noted.

% ---------------------------------------------------------------------
% SUBSECTION A.1: HOST NETWORK CONFIGURATIONS
% ---------------------------------------------------------------------
% This subsection preserves the final netplan configurations for the
% principal Linux systems in the lab so that the routing and interface
% design described in the body can be reconstructed directly from raw
% source files.
% ---------------------------------------------------------------------

\subsection{A.1 Host Network Configurations}
\label{sec:appendix-a-host-network-configurations}

% ---------------------------------------------------------------------
% SUBSECTION A.1.1: RADIUS-SPLUNK NETPLAN CONFIGURATION
% ---------------------------------------------------------------------
% This subsection preserves the final radius-splunk netplan file, which
% defines the management and service-side interfaces and static routes
% toward the authenticated and transit network segments.
% ---------------------------------------------------------------------

\subsubsection{A.1.1 radius-splunk Netplan Configuration}
\label{sec:appendix-a-radius-splunk-netplan-config}

\begin{verbatim}
radius-splunk: /etc/netplan/50-cloud-init.yaml
network:
  version: 2
  ethernets:
    ens33:
      dhcp4: true
    ens38:
      dhcp4: false
      addresses:
        - 10.0.30.10/24
      routes:
        - to: 10.0.34.0/30
          via: 10.0.30.1
        - to: 10.0.23.0/24
          via: 10.0.30.1
        - to: 10.0.12.0/24
          via: 10.0.30.1
\end{verbatim}

% ---------------------------------------------------------------------
% SUBSECTION A.1.2: ACCESS-NODE NETPLAN CONFIGURATION
% ---------------------------------------------------------------------
% This subsection preserves the final access-node netplan file, which
% defines the management interface, the authenticated client-facing
% segment, and the routed upstream segment toward dist-router.
% ---------------------------------------------------------------------

\subsubsection{A.1.2 access-node Netplan Configuration}
\label{sec:appendix-a-access-node-netplan-config}

\begin{verbatim}
access-node: /etc/netplan/50-cloud-init.yaml
network:
  version: 2
  ethernets:
    ens33:
      dhcp4: true
    ens38:
      dhcp4: false
      addresses:
        - 10.0.12.1/24
    ens39:
      dhcp4: false
      addresses:
        - 10.0.23.1/24
      routes:
        - to: 10.0.30.0/24
          via: 10.0.23.2
\end{verbatim}

% ---------------------------------------------------------------------
% SUBSECTION A.1.3: DIST-ROUTER NETPLAN CONFIGURATION
% ---------------------------------------------------------------------
% This subsection preserves the final dist-router netplan file, which
% defines the upstream/downstream routing role between access-node,
% pfSense, and the protected services segment.
% ---------------------------------------------------------------------

\subsubsection{A.1.3 dist-router Netplan Configuration}
\label{sec:appendix-a-dist-router-netplan-config}

\begin{verbatim}
dist-router: /etc/netplan/50-cloud-init.yaml
network:
  version: 2
  ethernets:
    ens33:
      dhcp4: true
    ens38:
      dhcp4: false
      addresses:
        - 10.0.23.2/24
      routes:
        - to: 10.0.12.0/24
          via: 10.0.23.1
    ens39:
      dhcp4: false
      addresses:
        - 10.0.34.1/30
      routes:
        - to: 10.0.30.0/24
          via: 10.0.34.2
\end{verbatim}

% ---------------------------------------------------------------------
% SUBSECTION A.1.4: OPS-GLPI NETPLAN CONFIGURATION
% ---------------------------------------------------------------------
% This subsection preserves the final ops-glpi netplan file, which
% defines the management interface, protected services interface, and
% routes back toward the authenticated and transit networks.
% ---------------------------------------------------------------------

\subsubsection{A.1.4 ops-glpi Netplan Configuration}
\label{sec:appendix-a-ops-glpi-netplan-config}

\begin{verbatim}
ops-glpi: /etc/netplan/50-cloud-init.yaml
network:
  version: 2
  ethernets:
    ens33:
      dhcp4: true
    ens38:
      dhcp4: false
      addresses:
        - 10.0.30.20/24
      routes:
        - to: 10.0.34.0/30
          via: 10.0.30.1
        - to: 10.0.23.0/24
          via: 10.0.30.1
        - to: 10.0.12.0/24
          via: 10.0.30.1
\end{verbatim}

% ---------------------------------------------------------------------
% SUBSECTION A.1.5: CLIENT-1 NETPLAN CONFIGURATION
% ---------------------------------------------------------------------
% This subsection preserves the final client-1 netplan file, which
% defines the static authenticated endpoint address and default route
% toward the access-node controlled segment.
% ---------------------------------------------------------------------

\subsubsection{A.1.5 client-1 Netplan Configuration}
\label{sec:appendix-a-client-1-netplan-config}

\begin{verbatim}
client-1: /etc/netplan/50-cloud-init.yaml
network:
  version: 2
  ethernets:
    ens33:
      addresses:
        - 10.0.12.10/24
      routes:
        - to: default
          via: 10.0.12.1
\end{verbatim}

% ---------------------------------------------------------------------
% SUBSECTION A.1.6: GUEST-TEST-1 NETPLAN CONFIGURATION
% ---------------------------------------------------------------------
% This subsection preserves the final guest-test-1 netplan file, which
% demonstrates DHCP-based onboarding on the guest segment.
% ---------------------------------------------------------------------

\subsubsection{A.1.6 guest-test-1 Netplan Configuration}
\label{sec:appendix-a-guest-test-1-netplan-config}

\begin{verbatim}
guest-test-1: /etc/netplan/50-cloud-init.yaml
network:
  version: 2
  ethernets:
    ens33:
      dhcp4: true
\end{verbatim}

% ---------------------------------------------------------------------
% SUBSECTION A.2: AUTHENTICATOR, SUPPLICANT, AND AAA CONFIGURATIONS
% ---------------------------------------------------------------------
% This subsection preserves the raw configuration artifacts that define
% the wired 802.1X control path: hostapd in wired mode, the systemd
% service unit used to run it, the client-side wpa_supplicant
% configuration, and the FreeRADIUS client and user definitions.
% ---------------------------------------------------------------------

\subsection{A.2 Authenticator, Supplicant, and AAA Configurations}
\label{sec:appendix-a-authenticator-supplicant-and-aaa-configurations}

% ---------------------------------------------------------------------
% SUBSECTION A.2.1: HOSTAPD WIRED CONFIGURATION
% ---------------------------------------------------------------------
% This subsection preserves the final wired-mode hostapd configuration
% used by access-node to act as the lab authenticator for client-1.
% ---------------------------------------------------------------------

\subsubsection{A.2.1 hostapd Wired Configuration}
\label{sec:appendix-a-access-node-hostapd-wired-config}

\begin{verbatim}
access-node: /etc/hostapd/lab/hostapd-wired.conf
interface=ens38
driver=wired

ieee8021x=1
eapol_version=2
eap_reauth_period=3600

auth_server_addr=10.0.30.10
auth_server_port=1812
auth_server_shared_secret=<REDACTED>

own_ip_addr=10.0.23.1

logger_stdout=-1
logger_stdout_level=0
logger_syslog=-1
logger_syslog_level=0

ctrl_interface=/var/run/hostapd
ctrl_interface_group=0
\end{verbatim}

% ---------------------------------------------------------------------
% SUBSECTION A.2.2: HOSTAPD WIRED SYSTEMD UNIT
% ---------------------------------------------------------------------
% This subsection preserves the systemd unit used to run hostapd in
% wired 802.1X mode on access-node.
% ---------------------------------------------------------------------

\subsubsection{A.2.2 hostapd Wired systemd Unit}
\label{sec:appendix-a-access-node-hostapd-wired-systemd-unit}

\begin{verbatim}
access-node: /etc/systemd/system/hostapd-wired.service
[Unit]
Description=hostapd wired 802.1X authenticator (lab)
After=network-online.target
Wants=network-online.target

[Service]
Type=simple
ExecStart=/usr/sbin/hostapd -dd /etc/hostapd/lab/hostapd-wired.conf
Restart=on-failure
RestartSec=2

[Install]
WantedBy=multi-user.target
\end{verbatim}

% ---------------------------------------------------------------------
% SUBSECTION A.2.3: CLIENT-1 WPA_SUPPLICANT WIRED CONFIGURATION
% ---------------------------------------------------------------------
% This subsection preserves the final wired supplicant configuration
% used by client-1 during successful authenticated operation.
% ---------------------------------------------------------------------

\subsubsection{A.2.3 client-1 wpa\_supplicant Wired Configuration}
\label{sec:appendix-a-client-1-wpa-supplicant-wired-8021x-config}

\begin{verbatim}
client-1: /etc/wpa_supplicant/lab/wired-8021x.conf
ctrl_interface=/run/wpa_supplicant
ap_scan=0
eapol_version=2

network={
    key_mgmt=IEEE8021X
    eap=MD5
    identity="nacuser"
    password="<REDACTED>"
}
\end{verbatim}

% ---------------------------------------------------------------------
% SUBSECTION A.2.4: FREERADIUS CLIENTS CONFIGURATION
% ---------------------------------------------------------------------
% This subsection preserves the authoritative FreeRADIUS clients
% configuration used in the lab, including the explicitly defined
% access-node client stanza.
% ---------------------------------------------------------------------

\subsubsection{A.2.4 FreeRADIUS Clients Configuration}
\label{sec:appendix-a-radius-splunk-freeradius-clients-config}

\begin{verbatim}
radius-splunk: /etc/freeradius/3.0/clients.conf
# -*- text -*-
##
## clients.conf -- client configuration directives
##
##	$Id: <REDACTED> $

#######################################################################
#
#  Define RADIUS clients (usually a NAS, Access Point, etc.).
#
#  [commentary omitted in summary: file preserved exactly as provided]
#
client localhost {
	ipaddr = 127.0.0.1
	proto = *
	secret = testing123
	nas_type = other
	limit {
		max_connections = 16
		lifetime = 0
		idle_timeout = 30
	}
}

client localhost_ipv6 {
	ipv6addr	= ::1
	secret		= testing123
}

client access-node {
    ipaddr = 10.0.23.1
    secret = <REDACTED>
    shortname = access-node
}
\end{verbatim}

% ---------------------------------------------------------------------
% SUBSECTION A.2.5: FREERADIUS USER AUTHORIZATION FILE
% ---------------------------------------------------------------------
% This subsection preserves the authoritative FreeRADIUS file-backed
% user definition used for nacuser in the lab.
% ---------------------------------------------------------------------

\subsubsection{A.2.5 FreeRADIUS User Authorization File}
\label{sec:appendix-a-radius-splunk-freeradius-users-config}

\begin{verbatim}
radius-splunk: /etc/freeradius/3.0/mods-config/files/authorize
#
# 	Configuration file for the rlm_files module.
# 	Please see rlm_files(5) manpage for more information.
#
# 	[commentary and example entries preserved as part of the file]
#

DEFAULT	Framed-Protocol == PPP
	Framed-Protocol = PPP,
	Framed-Compression = Van-Jacobson-TCP-IP

DEFAULT	Hint == "CSLIP"
	Framed-Protocol = SLIP,
	Framed-Compression = Van-Jacobson-TCP-IP

DEFAULT	Hint == "SLIP"
	Framed-Protocol = SLIP

#########################################################
# You should add test accounts to the TOP of this file! #
# See the example user "bob" above.                     #
#########################################################

nacuser Cleartext-Password := "<REDACTED>"
\end{verbatim}

% ---------------------------------------------------------------------
% SUBSECTION A.3: PFSENSE MANUAL CONFIGURATION TRANSCRIPTION
% ---------------------------------------------------------------------
% This subsection preserves the authoritative manual transcription of
% the pfSense interface addressing, guest DHCP scope, guest firewall
% policy, and outbound NAT mode used in the final lab.
% ---------------------------------------------------------------------

\subsection{A.3 pfSense Manual Configuration Transcription}
\label{sec:appendix-a-pfsense-manual-configuration-transcription}

% ---------------------------------------------------------------------
% SUBSECTION A.3.1: PFSENSE AUTHENTICATED / PROTECTED INTERFACE SUMMARY
% ---------------------------------------------------------------------
% This subsection preserves the final authoritative interface summary
% for the pfSense VM as manually transcribed from the implemented state.
% ---------------------------------------------------------------------

\subsubsection{A.3.1 pfSense Authenticated / Protected Interface Summary}
\label{sec:appendix-a-pfsense-authenticated-interface-summary}

\begin{verbatim}
WAN (wan) -> em0 -> v4/DHCP4: 192.168.195.145/24
LAN (lan) -> em1 -> v4: 10.0.34.2/30
OPT1 (opt1) -> em2 -> v4: 10.0.30.1/24
GUEST (opt2) -> em3 -> v4: 10.0.20.254/24
\end{verbatim}

% ---------------------------------------------------------------------
% SUBSECTION A.3.2: PFSENSE GUEST INTERFACE SUMMARY
% ---------------------------------------------------------------------
% This subsection preserves the authoritative guest-segment DHCP and
% addressing summary as manually transcribed from pfSense.
% ---------------------------------------------------------------------

\subsubsection{A.3.2 pfSense Guest Interface Summary}
\label{sec:appendix-a-pfsense-guest-interface-summary}

\begin{verbatim}
Enabled (Enable DHCP server on GUEST interface)
Subnet: 10.0.20.0/24
Subnet Range: 10.0.20.1 - 10.0.20.254
Address Pool Range: 10.0.20.100 - 10.0.20.199
\end{verbatim}

% ---------------------------------------------------------------------
% SUBSECTION A.3.3: PFSENSE GUEST FIREWALL RULES
% ---------------------------------------------------------------------
% This subsection preserves the final top-to-bottom guest firewall rule
% order that enforces isolation between the guest segment and protected
% internal networks while allowing internet-bound traffic.
% ---------------------------------------------------------------------

\subsubsection{A.3.3 pfSense Guest Firewall Rules}
\label{sec:appendix-a-pfsense-guest-firewall-rules}

\begin{verbatim}
Pass  | IPv4* | Source: 10.0.20.0/24 | Destination: This Firewall (Self) 
        | ALLOW-GUEST-TO-PFSENSE
Block | IPv4* | Source: 10.0.20.0/24 | Destination: 10.0.30.0/24       
        | BLOCK-GUEST-TO-SERVICES
Block | IPv4* | Source: 10.0.20.0/24 | Destination: 10.0.34.0/30       
        | BLOCK-GUEST-TO-TRANSIT
Block | IPv4* | Source: 10.0.20.0/24 | Destination: 10.0.23.0/24       
        | BLOCK-GUEST-TO-DIST
Block | IPv4* | Source: 10.0.20.0/24 | Destination: 10.0.12.0/24       
        | BLOCK-GUEST-TO-AUTH-ACCESS
Pass  | IPv4* | Source: 10.0.20.0/24 | Destination: *                  
        | ALLOW-GUEST-TO-INTERNET-AND-OTHER-NON-INTERNAL
\end{verbatim}

% ---------------------------------------------------------------------
% SUBSECTION A.3.4: PFSENSE OUTBOUND NAT SUMMARY
% ---------------------------------------------------------------------
% This subsection preserves the authoritative outbound NAT mode
% confirmation for the pfSense VM.
% ---------------------------------------------------------------------

\subsubsection{A.3.4 pfSense Outbound NAT Summary}
\label{sec:appendix-a-pfsense-outbound-nat-summary}

\begin{verbatim}
Automatic
\end{verbatim}

% ---------------------------------------------------------------------
% SUBSECTION A.4: AUTOMATION ARTIFACTS
% ---------------------------------------------------------------------
% This subsection preserves the final authoritative automation source
% files and supporting inventory/listing artifacts used by the Ansible
% and Python workflows described in the body of the report.
% ---------------------------------------------------------------------

\subsection{A.4 Automation Artifacts}
\label{sec:appendix-a-automation-artifacts}

% ---------------------------------------------------------------------
% SUBSECTION A.4.1: ANSIBLE INVENTORY
% ---------------------------------------------------------------------
% This subsection preserves the authoritative Ansible inventory used for
% the validation-first with optional remediation workflow.
% ---------------------------------------------------------------------

\subsubsection{A.4.1 Ansible Inventory}
\label{sec:appendix-a-ansible-inventory-hosts-ini}

\begin{verbatim}
radius-splunk: ~/ansible-nac/inventory/hosts.ini
[nac_nodes]
access-node ansible_host=10.0.23.1 ansible_user=nac-lab
radius-splunk ansible_host=127.0.0.1 ansible_connection=local

[authenticator]
access-node

[radius]
radius-splunk
\end{verbatim}

% ---------------------------------------------------------------------
% SUBSECTION A.4.2: ANSIBLE NAC VALIDATION PLAYBOOK
% ---------------------------------------------------------------------
% This subsection preserves the final authoritative Ansible playbook
% used for service validation, optional remediation, and local artifact
% capture.
% ---------------------------------------------------------------------

\subsubsection{A.4.2 Ansible NAC Validation Playbook}
\label{sec:appendix-a-ansible-nac-validation-playbook}

\begin{verbatim}
radius-splunk: ~/ansible-nac/playbooks/nac_validation.yml
---
- name: Validate NAC services and optionally remediate
  hosts: nac_nodes
  gather_facts: false
  become: false

  vars:
    remediate: false
    artifact_dir: "/home/nac-lab/ansible-nac/artifacts"

  tasks:
    - name: Collect basic hostname
      command: hostname
      register: hostname_out
      changed_when: false

    - name: Check hostapd-wired service on access-node
      become: true
      command: systemctl status hostapd-wired.service --no-pager
      register: hostapd_status_before
      changed_when: false
      failed_when: false
      when: inventory_hostname == "access-node"

    - name: Check freeradius service on radius-splunk
      become: true
      command: systemctl status freeradius --no-pager
      register: freeradius_status_before
      changed_when: false
      failed_when: false
      when: inventory_hostname == "radius-splunk"

    - name: Restart hostapd-wired service when remediation is enabled
      become: true
      command: systemctl restart hostapd-wired.service
      register: hostapd_restart
      changed_when: true
      failed_when: false
      when:
        - inventory_hostname == "access-node"
        - remediate | bool

    - name: Restart freeradius service when remediation is enabled
      become: true
      command: systemctl restart freeradius
      register: freeradius_restart
      changed_when: true
      failed_when: false
      when:
        - inventory_hostname == "radius-splunk"
        - remediate | bool

    - name: Re-check hostapd-wired service on access-node
      become: true
      command: systemctl status hostapd-wired.service --no-pager
      register: hostapd_status_after
      changed_when: false
      failed_when: false
      when: inventory_hostname == "access-node"

    - name: Re-check freeradius service on radius-splunk
      become: true
      command: systemctl status freeradius --no-pager
      register: freeradius_status_after
      changed_when: false
      failed_when: false
      when: inventory_hostname == "radius-splunk"

    - name: Collect recent hostapd logs
      become: true
      command: journalctl -u hostapd-wired.service -n 40 --no-pager
      register: hostapd_logs
      changed_when: false
      failed_when: false
      when: inventory_hostname == "access-node"

    - name: Collect recent freeradius service logs
      become: true
      command: journalctl -u freeradius -n 40 --no-pager
      register: freeradius_logs
      changed_when: false
      failed_when: false
      when: inventory_hostname == "radius-splunk"

    - name: Collect recent freeradius application log
      become: true
      command: tail -n 40 /var/log/freeradius/radius.log
      register: freeradius_radiuslog
      changed_when: false
      failed_when: false
      when: inventory_hostname == "radius-splunk"

    - name: Save hostapd pre-check status locally
      become: false
      delegate_to: localhost
      copy:
        content: "{{ hostapd_status_before.stdout | default('N/A') }}"
        dest: "{{ artifact_dir }}/{{ inventory_hostname }}_hostapd_status_before.txt"
      when: inventory_hostname == "access-node"

    - name: Save hostapd post-check status locally
      become: false
      delegate_to: localhost
      copy:
        content: "{{ hostapd_status_after.stdout | default('N/A') }}"
        dest: "{{ artifact_dir }}/{{ inventory_hostname }}_hostapd_status_after.txt"
      when: inventory_hostname == "access-node"

    - name: Save hostapd restart result locally
      become: false
      delegate_to: localhost
      copy:
        content: "{{ hostapd_restart.stdout | default('remediation not requested') }}"
        dest: "{{ artifact_dir }}/{{ inventory_hostname }}_hostapd_restart_result.txt"
      when: inventory_hostname == "access-node"

    - name: Save hostapd logs locally
      become: false
      delegate_to: localhost
      copy:
        content: "{{ hostapd_logs.stdout | default('N/A') }}"
        dest: "{{ artifact_dir }}/{{ inventory_hostname }}_hostapd_logs.txt"
      when: inventory_hostname == "access-node"

    - name: Save freeradius pre-check status locally
      become: false
      delegate_to: localhost
      copy:
        content: "{{ freeradius_status_before.stdout | default('N/A') }}"
        dest: "{{ artifact_dir }}/{{ inventory_hostname }}_freeradius_status_before.txt"
      when: inventory_hostname == "radius-splunk"

    - name: Save freeradius post-check status locally
      become: false
      delegate_to: localhost
      copy:
        content: "{{ freeradius_status_after.stdout | default('N/A') }}"
        dest: "{{ artifact_dir }}/{{ inventory_hostname }}_freeradius_status_after.txt"
      when: inventory_hostname == "radius-splunk"

    - name: Save freeradius restart result locally
      become: false
      delegate_to: localhost
      copy:
        content: "{{ freeradius_restart.stdout | default('remediation not requested') }}"
        dest: "{{ artifact_dir }}/{{ inventory_hostname }}_freeradius_restart_result.txt"
      when: inventory_hostname == "radius-splunk"

    - name: Save freeradius service logs locally
      become: false
      delegate_to: localhost
      copy:
        content: "{{ freeradius_logs.stdout | default('N/A') }}"
        dest: "{{ artifact_dir }}/{{ inventory_hostname }}_freeradius_journal.txt"
      when: inventory_hostname == "radius-splunk"

    - name: Save freeradius application log locally
      become: false
      delegate_to: localhost
      copy:
        content: "{{ freeradius_radiuslog.stdout | default('N/A') }}"
        dest: "{{ artifact_dir }}/{{ inventory_hostname }}_freeradius_radiuslog.txt"
      when: inventory_hostname == "radius-splunk"
\end{verbatim}

% ---------------------------------------------------------------------
% SUBSECTION A.4.3: PYTHON NAC INCIDENT PARSER
% ---------------------------------------------------------------------
% This subsection preserves the final Python parser used to summarize
% the incident evidence set.
% ---------------------------------------------------------------------

\subsubsection{A.4.3 Python NAC Incident Parser}
\label{sec:appendix-a-python-nac-incident-parser}

\begin{verbatim}
radius-splunk: ~/incident-evidence/nac_incident_parser.py
#!/usr/bin/env python3

from pathlib import Path

BASE = Path.home() / "incident-evidence"

FILES = {
    "client_failed": BASE / "wpa_supplicant_failed_auth_incident.txt",
    "client_success": BASE / "wpa_supplicant_success_retest_v2.txt",
    "hostapd_failed": BASE / "hostapd_failed_auth_incident.txt",
    "hostapd_success": BASE / "hostapd_success_retest_v2.txt",
    "radius_failed": BASE / "freeradius_failed_auth_incident.txt",
}

def read_text(path):
    try:
        return path.read_text(errors="ignore")
    except FileNotFoundError:
        return None

def present(text, needles):
    if text is None:
        return False
    return any(n in text for n in needles)

def summarize():
    data = {name: read_text(path) for name, path in FILES.items()}

    results = {
        "client_failed_unauthorized": present(data["client_failed"], [
            "Supplicant port status: Unauthorized"
        ]),
        "client_failed_connected": present(data["client_failed"], [
            "CTRL-EVENT-CONNECTED"
        ]),
        "client_success_connected": present(data["client_success"], [
            "CTRL-EVENT-CONNECTED"
        ]),
        "client_success_authorized": present(data["client_success"], [
            "Supplicant port status: Authorized"
        ]),
        "hostapd_failed_attempt_seen": present(data["hostapd_failed"], [
            "IEEE 802.1X: start authentication",
            "RADIUS Sending RADIUS message to authentication server",
            "CTRL-EVENT-EAP-STARTED",
            "recvmsg [RADIUS]"
        ]),
        "hostapd_success_authorized": present(data["hostapd_success"], [
            "AP-STA-CONNECTED",
            "authorizing port",
            "CTRL-EVENT-EAP-SUCCESS"
        ]),
        "radius_failed_reject": present(data["radius_failed"], [
            "Failed to authenticate the user",
            "Using Post-Auth-Type Reject",
            "EAP Failure",
            "Sending EAP Failure"
        ]),
    }

    print("=== NAC Incident Evidence Summary ===")
    print()

    print("[1] Failed authentication simulation")
    print(f"- Client remained unauthorized: {'YES' if results['client_failed_unauthorized']
            else 'NO'}")
    print(f"- Client showed successful connection during failed test: {'YES' if 
            results['client_failed_connected'] else 'NO'}")
    print(f"- hostapd saw authentication attempt: {'YES' if 
            results['hostapd_failed_attempt_seen'] else 'NO'}")
    print(f"- FreeRADIUS recorded failed auth / reject path: {'YES' if 
            results['radius_failed_reject'] else 'NO'}")
    print()

    print("[2] Remediation / successful re-test")
    print(f"- Client reached CTRL-EVENT-CONNECTED: {'YES' if 
            results['client_success_connected'] else 'NO'}")
    print(f"- Client reached Authorized state: {'YES' if 
            results['client_success_authorized'] else 'NO'}")
    print(f"- hostapd showed successful authorization indicators: {'YES' if 
            results['hostapd_success_authorized'] else 'NO'}")
    print()

    if (
        results["client_failed_unauthorized"]
        and results["radius_failed_reject"]
        and results["client_success_connected"]
        and results["client_success_authorized"]
        and results["hostapd_success_authorized"]
    ):
        print("[3] Incident conclusion")
        print("- SIMULATED NAC INCIDENT SUCCESSFULLY REPRODUCED AND REMEDIATED")
    else:
        print("[3] Incident conclusion")
        print("- Evidence is incomplete or mixed; review raw logs for additional detail")

if __name__ == "__main__":
    summarize()
\end{verbatim}

% ---------------------------------------------------------------------
% SUBSECTION A.4.4: ANSIBLE ARTIFACT DIRECTORY LISTING
% ---------------------------------------------------------------------
% This subsection preserves the authoritative post-run Ansible artifact
% directory listing referenced in the automation section.
% ---------------------------------------------------------------------

\subsubsection{A.4.4 Ansible Artifact Directory Listing}
\label{sec:appendix-a-ansible-artifact-directory-listing}

\begin{verbatim}
radius-splunk: ls -lh ~/ansible-nac/artifacts
total 44K
-rw-r--r-- 1 nac-lab nac-lab 3.5K Mar 24 12:23 access-node_hostapd_logs.txt
-rw-r--r-- 1 nac-lab nac-lab   25 Mar 24 12:43 access-node_hostapd_restart_result.txt
-rw-r--r-- 1 nac-lab nac-lab 1.4K Mar 24 12:43 access-node_hostapd_status_after.txt
-rw-r--r-- 1 nac-lab nac-lab 1.4K Mar 24 12:43 access-node_hostapd_status_before.txt
-rw-r--r-- 1 root    root    1.4K Mar 23 19:30 access-node_hostapd_status.txt
-rw-r--r-- 1 nac-lab nac-lab 4.0K Mar 24 12:23 radius-splunk_freeradius_journal.txt
-rw-r--r-- 1 nac-lab nac-lab 3.1K Mar 24 12:23 radius-splunk_freeradius_radiuslog.txt
-rw-r--r-- 1 nac-lab nac-lab   25 Mar 24 12:43 radius-splunk_freeradius_restart_result.txt
-rw-r--r-- 1 nac-lab nac-lab 1.8K Mar 24 12:43 radius-splunk_freeradius_status_after.txt
-rw-r--r-- 1 nac-lab nac-lab 1.8K Mar 24 12:43 radius-splunk_freeradius_status_before.txt
-rw-r--r-- 1 root    root    1.8K Mar 23 19:30 radius-splunk_freeradius_status.txt
\end{verbatim}

% ---------------------------------------------------------------------
% SUBSECTION A.4.5: INCIDENT EVIDENCE DIRECTORY LISTING
% ---------------------------------------------------------------------
% This subsection preserves the authoritative incident evidence
% directory listing referenced in the incident and observability
% sections.
% ---------------------------------------------------------------------

\subsubsection{A.4.5 Incident Evidence Directory Listing}
\label{sec:appendix-a-incident-evidence-directory-listing}

\begin{verbatim}
radius-splunk: ls -lh ~/incident-evidence
total 92K
-rw-rw-r-- 1 nac-lab nac-lab  32K Mar 23 04:02 freeradius_failed_auth_incident.txt
-rw-rw-r-- 1 nac-lab nac-lab  11K Mar 23 04:01 hostapd_failed_auth_incident.txt
-rw-rw-r-- 1 nac-lab nac-lab  19K Mar 23 04:01 hostapd_success_retest_v2.txt
-rw-r--r-- 1 root    root    3.5K Mar 24 09:52 nac_incident_parser.py
-rw-rw-r-- 1 nac-lab nac-lab  535 Mar 23 04:28 nac_incident_parser_success.txt
-rw-rw-r-- 1 nac-lab nac-lab 7.3K Mar 23 03:59 wpa_supplicant_failed_auth_incident.txt
-rw-rw-r-- 1 nac-lab nac-lab 8.4K Mar 23 03:59 wpa_supplicant_success_retest_v2.txt
\end{verbatim}

% ---------------------------------------------------------------------
% SUBSECTION A.4.6: SPLUNK UNIVERSAL FORWARDER INPUTS CONFIGURATION
% ---------------------------------------------------------------------
% This subsection preserves the authoritative Splunk Universal
% Forwarder monitored-input configuration referenced in the
% observability section.
% ---------------------------------------------------------------------

\subsubsection{A.4.6 Splunk Universal Forwarder inputs.conf}
\label{sec:appendix-a-splunk-uf-inputs-conf}

\begin{verbatim}
radius-splunk: /opt/splunkforwarder/etc/apps/lab_nac_inputs/local/inputs.conf
[monitor:///var/log/syslog]
disabled = false
index = main
sourcetype = syslog

[monitor:///var/log/auth.log]
disabled = false
index = main
sourcetype = linux_secure

[monitor:///var/log/freeradius/radius.log]
disabled = false
index = main
sourcetype = freeradius
\end{verbatim}

% =====================================================================
% APPENDIX B: INCIDENT, VALIDATION, AND EVIDENCE ARTIFACTS
% =====================================================================
% This appendix preserves the raw incident artifacts, validation
% outputs, selected observability excerpts, packet-capture inventory,
% and figure-mapping metadata referenced throughout the body of the
% report. Its purpose is to provide evidentiary completeness without
% requiring the main body to inline every raw log and command output.
% =====================================================================
\appendix

\section{Appendix B. Incident, Validation, and Evidence Artifacts}
\label{sec:appendix-b-incident-validation-and-evidence-artifacts}

This appendix preserves the authoritative raw evidence artifacts referenced throughout the report. The materials included here are grouped by function: incident reproduction evidence, GLPI incident record content, guest-path validation outputs, automation execution evidence, selected observability excerpts, and packet-capture metadata.

% ---------------------------------------------------------------------
% SUBSECTION B.1: INCIDENT EVIDENCE FILES
% ---------------------------------------------------------------------
% This subsection preserves the raw captured evidence files associated
% with the failed authentication simulation and the successful
% remediation re-test.
% ---------------------------------------------------------------------

\subsection{B.1 Incident Evidence Files}
\label{sec:appendix-b-incident-evidence-files}

% ---------------------------------------------------------------------
% SUBSECTION B.1.1: WPA_SUPPLICANT FAILED AUTHENTICATION LOG
% ---------------------------------------------------------------------
% This subsection preserves the raw failed-authentication capture taken
% from the manual client-side wpa_supplicant debug run.
% ---------------------------------------------------------------------

\subsubsection{B.1.1 wpa\_supplicant Failed Authentication Log}
\label{sec:appendix-b-wpa-supplicant-failed-auth-incident-log}

\begin{verbatim}
radius-splunk: ~/incident-evidence/wpa_supplicant_failed_auth_incident.txt
wpa_supplicant v2.10
random: getrandom() support available
Successfully initialized wpa_supplicant
Initializing interface 'ens33' conf '/etc/wpa_supplicant/lab/wired-8021x.conf' driver 
        'wired' ctrl_interface 'N/A' bridge 'N/A'
Configuration file '/etc/wpa_supplicant/lab/wired-8021x.conf' -> '/etc/wpa_supplicant/lab/
        wired-8021x.conf'
Reading configuration file '/etc/wpa_supplicant/lab/wired-8021x.conf'
ctrl_interface='/run/wpa_supplicant'
ap_scan=0
eapol_version=2
Line: 5 - start of a new network block
key_mgmt: 0x8
eap methods - hexdump(len=16): 00 00 00 00 04 00 00 00 00 00 00 00 00 00 00 00
identity - hexdump_ascii(len=7):
     6e 61 63 75 73 65 72                              nacuser         
password - hexdump_ascii(len=13): [REMOVED]
Priority group 0
   id=0 ssid=''
driver_wired_init_common: Added multicast membership with packet socket
Add interface ens33 to a new radio N/A
ens33: Own MAC address: <REDACTED>
ens33: RSN: flushing PMKID list in the driver
ens33: Setting scan request: 0.100000 sec
TDLS: TDLS operation not supported by driver
TDLS: Driver uses internal link setup
TDLS: Driver does not support TDLS channel switching
ens33: WPS: UUID based on MAC address: <REDACTED>
ENGINE: Loading builtin engines
ENGINE: Loading builtin engines
EAPOL: SUPP_PAE entering state DISCONNECTED
EAPOL: Supplicant port status: Unauthorized
EAPOL: KEY_RX entering state NO_KEY_RECEIVE
EAPOL: SUPP_BE entering state INITIALIZE
EAP: EAP entering state DISABLED
MBO: Update non-preferred channels, non_pref_chan=N/A
ens33: Added interface ens33
ens33: State: DISCONNECTED -> DISCONNECTED
EAPOL: External notification - EAP success=0
EAPOL: External notification - EAP fail=0
EAPOL: External notification - portControl=Auto
ens33: Already associated with a configured network - generating associated event
ens33: Event ASSOC (0) received
ens33: Association info event
ens33: State: DISCONNECTED -> ASSOCIATED
ens33: Associated to a new BSS: BSSID=<REDACTED>
ens33: Select network based on association information
ens33: Network configuration found for the current AP
ens33: WPA: clearing AP WPA IE
ens33: WPA: clearing AP RSN IE
ens33: WPA: clearing AP RSNXE
ens33: WPA: clearing own WPA/RSN IE
ens33: RSN: clearing own RSNXE
ens33: Failed to get scan results
EAPOL: External notification - EAP success=0
EAPOL: External notification - EAP fail=0
EAPOL: External notification - portControl=Auto
ens33: Associated with <REDACTED>
ens33: WPA: Association event - clear replay counter
ens33: WPA: Clear old PTK
TDLS: Remove peers on association
EAPOL: External notification - portEnabled=0
EAPOL: External notification - portValid=0
EAPOL: External notification - portEnabled=1
EAPOL: SUPP_PAE entering state CONNECTING
EAPOL: SUPP_BE entering state IDLE
EAP: EAP entering state INITIALIZE
EAP: EAP entering state IDLE
ens33: Cancelling scan request
ens33: CTRL-EVENT-SUBNET-STATUS-UPDATE status=0
EAPOL: startWhen --> 0
EAPOL: SUPP_PAE entering state CONNECTING
EAPOL: txStart
TX EAPOL: dst=01:80:c2:00:00:03
TX EAPOL - hexdump(len=4): 02 01 00 00
l2_packet_receive: src=<REDACTED> len=46
ens33: RX EAPOL from <REDACTED>
RX EAPOL - hexdump(len=46): 02 00 00 05 01 19 00 05 01 00 00 00 00 00 00 00 00 00 00 00 00
         00 00 00 00 00 00 00 00 00 00 00 00 00 00 00 00 00 00 00 00 00 00 00 00 00
EAPOL: Received EAP-Packet frame
EAPOL: SUPP_PAE entering state RESTART
EAP: EAP entering state INITIALIZE
EAP: EAP entering state IDLE
EAPOL: SUPP_PAE entering state AUTHENTICATING
EAPOL: SUPP_BE entering state REQUEST
EAPOL: getSuppRsp
EAP: EAP entering state RECEIVED
EAP: Received EAP-Request id=25 method=1 vendor=0 vendorMethod=0
EAP: EAP entering state IDENTITY
ens33: CTRL-EVENT-EAP-STARTED EAP authentication started
EAP: Status notification: started (param=)
EAP: EAP-Request Identity data - hexdump_ascii(len=0):
EAP: using real identity - hexdump_ascii(len=7):
     6e 61 63 75 73 65 72                              nacuser         
EAP: EAP entering state SEND_RESPONSE
EAP: EAP entering state IDLE
EAPOL: SUPP_BE entering state RESPONSE
EAPOL: txSuppRsp
TX EAPOL: dst=<REDACTED>
TX EAPOL - hexdump(len=16): 02 00 00 0c 02 19 00 0c 01 6e 61 63 75 73 65 72
EAPOL: SUPP_BE entering state RECEIVE
l2_packet_receive: src=<REDACTED> len=46
ens33: RX EAPOL from <REDACTED>
RX EAPOL - hexdump(len=46): 02 00 00 16 01 1a 00 16 04 10 d8 d1 ca 68 87 96 28 22 03 d8 a5
         44 92 bb 84 f9 00 00 00 00 00 00 00 00 00 00 00 00 00 00 00 00 00 00 00 00
EAPOL: Received EAP-Packet frame
EAPOL: SUPP_BE entering state REQUEST
EAPOL: getSuppRsp
EAP: EAP entering state RECEIVED
EAP: Received EAP-Request id=26 method=4 vendor=0 vendorMethod=0
EAP: EAP entering state GET_METHOD
ens33: CTRL-EVENT-EAP-PROPOSED-METHOD vendor=0 method=4
EAP: Status notification: accept proposed method (param=MD5)
EAP: Initialize selected EAP method: vendor 0 method 4 (MD5)
ens33: CTRL-EVENT-EAP-METHOD EAP vendor 0 method 4 (MD5) selected
EAP: EAP entering state METHOD
EAP-MD5: Challenge - hexdump(len=16): d8 d1 ca 68 87 96 28 22 03 d8 a5 44 92 bb 84 f9
EAP-MD5: Generating Challenge Response
EAP-MD5: Response - hexdump(len=16): 41 b2 c0 1c cf a9 4d 12 fc 49 1a 04 d9 5c d9 35
EAP: method process -> ignore=FALSE methodState=DONE decision=COND_SUCC eapRespData=0x62e7
        14bc0420
EAP: EAP entering state SEND_RESPONSE
EAP: EAP entering state IDLE
EAPOL: SUPP_BE entering state RESPONSE
EAPOL: txSuppRsp
TX EAPOL: dst=<REDACTED>
TX EAPOL - hexdump(len=26): 02 00 00 16 02 1a 00 16 04 10 41 b2 c0 1c cf a9 4d 12 fc 49 1a
         04 d9 5c d9 35
EAPOL: SUPP_BE entering state RECEIVE
ens33: Removing interface ens33
ens33: Request to deauthenticate - bssid=<REDACTED> pending_bssid=00:00:00:00:00:00 reason
        =3 (DEAUTH_LEAVING) state=ASSOCIATED
TDLS: Tear down peers
ens33: Event DEAUTH (11) received
ens33: Deauthentication notification
ens33:  * reason 3 (DEAUTH_LEAVING) locally_generated=1
Deauthentication frame IE(s) - hexdump(len=0): [NULL]
ens33: CTRL-EVENT-DISCONNECTED bssid=<REDACTED> reason=3 locally_generated=1
ens33: Auto connect disabled: do not try to re-connect
ens33: Ignore connection failure indication since interface has been put into disconnected
         state
TDLS: Remove peers on disassociation
ens33: WPA: Clear old PMK and PTK
ens33: Disconnect event - remove keys
ens33: State: ASSOCIATED -> DISCONNECTED
EAPOL: External notification - portEnabled=0
EAPOL: SUPP_PAE entering state DISCONNECTED
EAPOL: Supplicant port status: Unauthorized
EAPOL: SUPP_BE entering state INITIALIZE
EAP: EAP entering state DISABLED
EAPOL: External notification - portValid=0
ens33: State: DISCONNECTED -> DISCONNECTED
EAPOL: External notification - portEnabled=0
EAPOL: External notification - portValid=0
QM: Clear all active DSCP policies
ens33: CTRL-EVENT-DSCP-POLICY clear_all
EAP: deinitialize previously used EAP method (4, MD5) at EAP deinit
ens33: WPA: Clear old PMK and PTK
ens33: Cancelling scan request
ens33: Cancelling authentication timeout
Off-channel: Clear pending Action frame TX (pending_action_tx=(nil)
HS20: Delete all stored icons
Off-channel: Action frame sequence done notification: pending_action_tx=(nil) drv_offchan_
        tx=0 action_tx_wait_time=0 off_channel_freq=0 roc_waiting_drv_freq=0
QM: Clear all active DSCP policies
ens33: CTRL-EVENT-DSCP-POLICY clear_all
Remove interface ens33 from radio 
Remove radio 
ens33: CTRL-EVENT-TERMINATING 
\end{verbatim}

% ---------------------------------------------------------------------
% SUBSECTION B.1.2: HOSTAPD FAILED AUTHENTICATION LOG
% ---------------------------------------------------------------------
% This subsection preserves the raw failed-authentication evidence
% captured from the hostapd side during the simulated incident.
% ---------------------------------------------------------------------

\subsubsection{B.1.2 hostapd Failed Authentication Log}
\label{sec:appendix-b-hostapd-failed-auth-incident-log}

\begin{verbatim}
radius-splunk: ~/incident-evidence/hostapd_failed_auth_incident.txt
Mar 22 22:32:33 access-node hostapd[2052]: EAP-Identity: Peer identity - hexdump_ascii(len
        =7):
Mar 22 22:32:33 access-node hostapd[2052]:      6e 61 63 75 73 65 72                      
                nacuser
Mar 22 22:32:33 access-node hostapd[2052]: EAP: EAP entering state SELECT_ACTION
Mar 22 22:32:33 access-node hostapd[2052]: EAP: getDecision: -> PASSTHROUGH
Mar 22 22:32:33 access-node hostapd[2052]: EAP: EAP entering state INITIALIZE_PASSTHROUGH
Mar 22 22:32:33 access-node hostapd[2052]: EAP: EAP entering state AAA_REQUEST
Mar 22 22:32:33 access-node hostapd[2052]: EAP: EAP entering state AAA_IDLE
Mar 22 22:32:33 access-node hostapd[2052]: ens38: STA <REDACTED> IEEE 802.1X: STA identity
         'nacuser'
Mar 22 22:32:33 access-node hostapd[2052]: Encapsulating EAP message into a RADIUS packet
Mar 22 22:32:33 access-node hostapd[2052]: ens38: STA <REDACTED> IEEE 802.1X: unauthorizin
        g port
Mar 22 22:32:33 access-node hostapd[2052]: ens38: STA <REDACTED> IEEE 802.1X: Sending EAP 
        Packet (identifier 25)
Mar 22 22:32:33 access-node hostapd[2052]: ens38: STA <REDACTED> IEEE 802.1X: received EAP
         packet (code=2 id=25 len=12) from STA: EAP Response-Identity (1)
Mar 22 22:32:33 access-node hostapd[2052]: ens38: STA <REDACTED> IEEE 802.1X: STA identity
         'nacuser'
Mar 22 22:32:33 access-node hostapd[2052]: ens38: RADIUS Sending RADIUS message to authent
        ication server
Mar 22 22:32:33 access-node hostapd[2052]: RADIUS message: code=1 (Access-Request) identif
        ier=0 length=165
Mar 22 22:32:33 access-node hostapd[2052]:    Attribute 1 (User-Name) length=9
Mar 22 22:32:33 access-node hostapd[2052]:       Value: 'nacuser'
Mar 22 22:32:33 access-node hostapd[2052]:    Attribute 4 (NAS-IP-Address) length=6
Mar 22 22:32:33 access-node hostapd[2052]:       Value: 10.0.23.1
Mar 22 22:32:33 access-node hostapd[2052]:    Attribute 30 (Called-Station-Id) length=20
Mar 22 22:32:33 access-node hostapd[2052]:       Value: '<REDACTED>:'
Mar 22 22:32:33 access-node hostapd[2052]:    Attribute 61 (NAS-Port-Type) length=6
Mar 22 22:32:33 access-node hostapd[2052]:       Value: 19
Mar 22 22:32:33 access-node hostapd[2052]:    Attribute 6 (Service-Type) length=6
Mar 22 22:32:33 access-node hostapd[2052]:       Value: 2
Mar 22 22:32:33 access-node hostapd[2052]:    Attribute 31 (Calling-Station-Id) length=19
Mar 22 22:32:33 access-node hostapd[2052]:       Value: '<REDACTED>'
Mar 22 22:32:33 access-node hostapd[2052]:    Attribute 77 (Connect-Info) length=23
Mar 22 22:32:33 access-node hostapd[2052]:       Value: 'CONNECT 0Mbps 802.11b'
Mar 22 22:32:33 access-node hostapd[2052]:    Attribute 44 (Acct-Session-Id) length=18
Mar 22 22:32:33 access-node hostapd[2052]:       Value: '1A082DB0BEE2BD34'
Mar 22 22:32:33 access-node hostapd[2052]:    Attribute 12 (Framed-MTU) length=6
Mar 22 22:32:33 access-node hostapd[2052]:       Value: 1400
Mar 22 22:32:33 access-node hostapd[2052]:    Attribute 79 (EAP-Message) length=14
Mar 22 22:32:33 access-node hostapd[2052]:       Value: 0219000c016e616375736572
Mar 22 22:32:33 access-node hostapd[2052]:    Attribute 80 (Message-Authenticator) length=
        18
Mar 22 22:32:33 access-node hostapd[2052]:       Value: 71423639f013148f1ce7a71576679bbe
Mar 22 22:32:33 access-node hostapd[2052]: ens38: RADIUS Next RADIUS client retransmit in 
        3 seconds
Mar 22 22:32:33 access-node hostapd[2052]: ens38: RADIUS Sending RADIUS message to authent
        ication server
Mar 22 22:32:33 access-node hostapd[2052]: ens38: RADIUS Next RADIUS client retransmit in 
        3 seconds
Mar 22 22:32:33 access-node hostapd[2052]: ens38: RADIUS Received 80 bytes from RADIUS ser
        ver
Mar 22 22:32:33 access-node hostapd[2052]: ens38: RADIUS Received RADIUS message
Mar 22 22:32:33 access-node hostapd[2052]: RADIUS message: code=11 (Access-Challenge) iden
        tifier=0 length=80
Mar 22 22:32:33 access-node hostapd[2052]:    Attribute 80 (Message-Authenticator) length=18
Mar 22 22:32:33 access-node hostapd[2052]:       Value: e72a3944e289c8a2d9b97adfbceb639a
Mar 22 22:32:33 access-node hostapd[2052]:    Attribute 79 (EAP-Message) length=24
Mar 22 22:32:33 access-node hostapd[2052]:       Value: 011a00160410d8d1ca688796282203d8a5
        4492bb84f9
Mar 22 22:32:33 access-node hostapd[2052]:    Attribute 24 (State) length=18
Mar 22 22:32:33 access-node hostapd[2052]:       Value: bb99b235bb83b6a0b74fb036b1c91a02
Mar 22 22:32:33 access-node hostapd[2052]: ens38: STA <REDACTED> RADIUS: Received RADIUS p
        acket matched with a pending request, round trip time 0.00 sec
Mar 22 22:32:33 access-node hostapd[2052]: ens38: RADIUS Received 80 bytes from RADIUS ser
        ver
Mar 22 22:32:33 access-node hostapd[2052]: ens38: RADIUS Received RADIUS message
Mar 22 22:32:33 access-node hostapd[2052]: ens38: STA <REDACTED> RADIUS: Received RADIUS p
        acket matched with a pending request, round trip time 0.00 sec
Mar 22 22:32:33 access-node hostapd[2052]: ens38: STA <REDACTED> IEEE 802.1X: decapsulated
         EAP packet (code=1 id=26 len=22) from RADIUS server: EAP-Request-MD5 (4)
Mar 22 22:32:33 access-node hostapd[2052]: RADIUS packet matching with station 00:0c:29:68
        :53:33
Mar 22 22:32:33 access-node hostapd[2052]: ens38: STA <REDACTED> IEEE 802.1X: decapsulated
         EAP packet (code=1 id=26 len=22) from RADIUS server: EAP-Request-MD5 (4)
Mar 22 22:32:33 access-node hostapd[2052]: EAP: EAP entering state AAA_RESPONSE
Mar 22 22:32:33 access-node hostapd[2052]: EAP: getId: id=26
Mar 22 22:32:33 access-node hostapd[2052]: EAP: EAP entering state SEND_REQUEST2
Mar 22 22:32:33 access-node hostapd[2052]: EAP: EAP entering state IDLE2
Mar 22 22:32:33 access-node hostapd[2052]: EAP: retransmit timeout 3 seconds (from dynamic
         back off; retransCount=0)
Mar 22 22:32:33 access-node hostapd[2052]: IEEE 802.1X: <REDACTED> BE_AUTH entering state 
        REQUEST
Mar 22 22:32:33 access-node hostapd[2052]: ens38: STA <REDACTED> IEEE 802.1X: Sending EAP 
        Packet (identifier 26)
Mar 22 22:32:33 access-node hostapd[2052]: ens38: STA <REDACTED> IEEE 802.1X: Sending EAP 
        Packet (identifier 26)
Mar 22 22:32:33 access-node hostapd[2052]: Received EAPOL packet
Mar 22 22:32:33 access-node hostapd[2052]: ens38: Event NEW_STA (22) received
Mar 22 22:32:33 access-node hostapd[2052]: ens38: Event EAPOL_RX (23) received
Mar 22 22:32:33 access-node hostapd[2052]: IEEE 802.1X: 46 bytes from 00:0c:29:68:53:33
Mar 22 22:32:33 access-node hostapd[2052]:    IEEE 802.1X: version=2 type=0 length=22
Mar 22 22:32:33 access-node hostapd[2052]:    ignoring 20 extra octets after IEEE 802.1X p
        acket
Mar 22 22:32:33 access-node hostapd[2052]: EAP: code=2 (response) identifier=26 length=22
Mar 22 22:32:33 access-node hostapd[2052]: ens38: STA <REDACTED> IEEE 802.1X: received EAP
         packet (code=2 id=26 len=22) from STA: EAP Response-MD5 (4)
Mar 22 22:32:33 access-node hostapd[2052]: IEEE 802.1X: <REDACTED> BE_AUTH entering state 
        RESPONSE
Mar 22 22:32:33 access-node hostapd[2052]: EAP: EAP entering state RECEIVED2
Mar 22 22:32:33 access-node hostapd[2052]: EAP: parseEapResp: rxResp=1 rxInitiate=0 respId
        =26 respMethod=4 respVendor=0 respVendorMethod=0
Mar 22 22:32:33 access-node hostapd[2052]: EAP: EAP entering state AAA_REQUEST
Mar 22 22:32:33 access-node hostapd[2052]: EAP: EAP entering state AAA_IDLE
Mar 22 22:32:33 access-node hostapd[2052]: Encapsulating EAP message into a RADIUS packet
Mar 22 22:32:33 access-node hostapd[2052]: Copied RADIUS State Attribute
Mar 22 22:32:33 access-node hostapd[2052]: ens38: RADIUS Sending RADIUS message to authent
        ication server
Mar 22 22:32:33 access-node hostapd[2052]: RADIUS message: code=1 (Access-Request) identif
        ier=1 length=193
Mar 22 22:32:33 access-node hostapd[2052]:    Attribute 1 (User-Name) length=9
Mar 22 22:32:33 access-node hostapd[2052]:       Value: 'nacuser'
Mar 22 22:32:33 access-node hostapd[2052]:    Attribute 4 (NAS-IP-Address) length=6
Mar 22 22:32:33 access-node hostapd[2052]:       Value: 10.0.23.1
Mar 22 22:32:33 access-node hostapd[2052]:    Attribute 30 (Called-Station-Id) length=20
Mar 22 22:32:33 access-node hostapd[2052]:       Value: '<REDACTED>:'
Mar 22 22:32:33 access-node hostapd[2052]:    Attribute 61 (NAS-Port-Type) length=6
Mar 22 22:32:33 access-node hostapd[2052]:       Value: 19
Mar 22 22:32:33 access-node hostapd[2052]:    Attribute 6 (Service-Type) length=6
Mar 22 22:32:33 access-node hostapd[2052]:       Value: 2
Mar 22 22:32:33 access-node hostapd[2052]:    Attribute 31 (Calling-Station-Id) length=19
Mar 22 22:32:33 access-node hostapd[2052]:       Value: '<REDACTED>'
Mar 22 22:32:33 access-node hostapd[2052]:    Attribute 77 (Connect-Info) length=23
Mar 22 22:32:33 access-node hostapd[2052]:       Value: 'CONNECT 0Mbps 802.11b'
Mar 22 22:32:33 access-node hostapd[2052]:    Attribute 44 (Acct-Session-Id) length=18
Mar 22 22:32:33 access-node hostapd[2052]:       Value: '1A082DB0BEE2BD34'
Mar 22 22:32:33 access-node hostapd[2052]:    Attribute 12 (Framed-MTU) length=6
Mar 22 22:32:33 access-node hostapd[2052]:       Value: 1400
Mar 22 22:32:33 access-node hostapd[2052]:    Attribute 79 (EAP-Message) length=24
Mar 22 22:32:33 access-node hostapd[2052]:       Value: 021a0016041041b2c01ccfa94d12fc491a
        04d95cd935
Mar 22 22:32:33 access-node hostapd[2052]:    Attribute 24 (State) length=18
Mar 22 22:32:33 access-node hostapd[2052]:       Value: bb99b235bb83b6a0b74fb036b1c91a02
Mar 22 22:32:33 access-node hostapd[2052]:    Attribute 80 (Message-Authenticator) length=
        18
Mar 22 22:32:33 access-node hostapd[2052]:       Value: ee00aa03a7391d212502d3e4595c510c
Mar 22 22:32:33 access-node hostapd[2052]: ens38: RADIUS Next RADIUS client retransmit in 
        3 seconds
Mar 22 22:32:33 access-node hostapd[2052]: ens38: STA <REDACTED> IEEE 802.1X: received EAP
         packet (code=2 id=26 len=22) from STA: EAP Response-MD5 (4)
Mar 22 22:32:33 access-node hostapd[2052]: ens38: RADIUS Sending RADIUS message to authent
        ication server
Mar 22 22:32:33 access-node hostapd[2052]: ens38: RADIUS Next RADIUS client retransmit in 
        3 seconds
Mar 22 22:32:36 access-node hostapd[2052]: IEEE 802.1X: <REDACTED> - (EAP) retransWhile --
        > 0
Mar 22 22:32:36 access-node hostapd[2052]: ens38: STA <REDACTED> RADIUS: Resending RADIUS 
        message (id=1)
Mar 22 22:32:36 access-node hostapd[2052]: ens38: STA <REDACTED> RADIUS: Resending RADIUS 
        message (id=1)
Mar 22 22:32:36 access-node hostapd[2052]: ens38: RADIUS Next RADIUS client retransmit in 
        6 seconds
Mar 22 22:32:36 access-node hostapd[2052]: ens38: RADIUS Next RADIUS client retransmit in 
        6 seconds
Mar 22 22:32:36 access-node hostapd[2052]: recvmsg[RADIUS]: Connection refused
Mar 22 22:32:42 access-node hostapd[2052]: ens38: STA <REDACTED> RADIUS: Resending RADIUS 
        message (id=1)
Mar 22 22:32:42 access-node hostapd[2052]: ens38: RADIUS Next RADIUS client retransmit in 
        12 seconds
Mar 22 22:32:42 access-node hostapd[2052]: ens38: STA <REDACTED> RADIUS: Resending RADIUS 
        message (id=1)
Mar 22 22:32:42 access-node hostapd[2052]: ens38: RADIUS Next RADIUS client retransmit in 
        12 seconds
Mar 22 22:32:42 access-node hostapd[2052]: recvmsg[RADIUS]: Connection refused
\end{verbatim}

% ---------------------------------------------------------------------
% SUBSECTION B.1.3: FREERADIUS FAILED AUTHENTICATION LOG
% ---------------------------------------------------------------------
% This subsection preserves the raw failed-authentication evidence
% captured from the FreeRADIUS debug run during the simulated incident.
% ---------------------------------------------------------------------

\subsubsection{B.1.3 FreeRADIUS Failed Authentication Log}
\label{sec:appendix-b-freeradius-failed-auth-incident-log}

\begin{verbatim}
radius-splunk: ~/incident-evidence/freeradius_failed_auth_incident.txt
FreeRADIUS Version 3.2.5
Copyright (C) 1999-2023 The FreeRADIUS server project and contributors
There is NO warranty; not even for MERCHANTABILITY or FITNESS FOR A
PARTICULAR PURPOSE
You may redistribute copies of FreeRADIUS under the terms of the
GNU General Public License
For more information about these matters, see the file named COPYRIGHT
Starting - reading configuration files ...
including dictionary file /usr/share/freeradius/dictionary
including dictionary file /usr/share/freeradius/dictionary.dhcp
including dictionary file /usr/share/freeradius/dictionary.vqp
including dictionary file /etc/freeradius/3.0/dictionary
including configuration file /etc/freeradius/3.0/radiusd.conf
including configuration file /etc/freeradius/3.0/proxy.conf
including configuration file /etc/freeradius/3.0/clients.conf
including files in directory /etc/freeradius/3.0/mods-enabled/
including configuration file /etc/freeradius/3.0/mods-enabled/always
including configuration file /etc/freeradius/3.0/mods-enabled/mschap
including configuration file /etc/freeradius/3.0/mods-enabled/radutmp
including configuration file /etc/freeradius/3.0/mods-enabled/exec
including configuration file /etc/freeradius/3.0/mods-enabled/sradutmp
including configuration file /etc/freeradius/3.0/mods-enabled/ntlm_auth
including configuration file /etc/freeradius/3.0/mods-enabled/preprocess
including configuration file /etc/freeradius/3.0/mods-enabled/detail
including configuration file /etc/freeradius/3.0/mods-enabled/soh
including configuration file /etc/freeradius/3.0/mods-enabled/echo
including configuration file /etc/freeradius/3.0/mods-enabled/expr
including configuration file /etc/freeradius/3.0/mods-enabled/unpack
including configuration file /etc/freeradius/3.0/mods-enabled/logintime
including configuration file /etc/freeradius/3.0/mods-enabled/realm
including configuration file /etc/freeradius/3.0/mods-enabled/attr_filter
including configuration file /etc/freeradius/3.0/mods-enabled/digest
including configuration file /etc/freeradius/3.0/mods-enabled/expiration
including configuration file /etc/freeradius/3.0/mods-enabled/utf8
including configuration file /etc/freeradius/3.0/mods-enabled/chap
including configuration file /etc/freeradius/3.0/mods-enabled/pap
including configuration file /etc/freeradius/3.0/mods-enabled/unix
including configuration file /etc/freeradius/3.0/mods-enabled/passwd
including configuration file /etc/freeradius/3.0/mods-enabled/detail.log
including configuration file /etc/freeradius/3.0/mods-enabled/eap
including configuration file /etc/freeradius/3.0/mods-enabled/replicate
including configuration file /etc/freeradius/3.0/mods-enabled/linelog
including configuration file /etc/freeradius/3.0/mods-enabled/dynamic_clients
including configuration file /etc/freeradius/3.0/mods-enabled/files
including files in directory /etc/freeradius/3.0/policy.d/
including configuration file /etc/freeradius/3.0/policy.d/canonicalization
including configuration file /etc/freeradius/3.0/policy.d/operator-name
including configuration file /etc/freeradius/3.0/policy.d/cui
including configuration file /etc/freeradius/3.0/policy.d/moonshot-targeted-ids
including configuration file /etc/freeradius/3.0/policy.d/accounting
including configuration file /etc/freeradius/3.0/policy.d/abfab-tr
including configuration file /etc/freeradius/3.0/policy.d/control
including configuration file /etc/freeradius/3.0/policy.d/filter
including configuration file /etc/freeradius/3.0/policy.d/debug
including configuration file /etc/freeradius/3.0/policy.d/eap
including configuration file /etc/freeradius/3.0/policy.d/rfc7542
including configuration file /etc/freeradius/3.0/policy.d/dhcp
including files in directory /etc/freeradius/3.0/sites-enabled/
including configuration file /etc/freeradius/3.0/sites-enabled/inner-tunnel
including configuration file /etc/freeradius/3.0/sites-enabled/default
main {
 security {
 	user = "freerad"
 	group = "freerad"
 	allow_core_dumps = no
 }
	name = "freeradius"
	prefix = "/usr"
	localstatedir = "/var"
	logdir = "/var/log/freeradius"
	run_dir = "/var/run/freeradius"
}
main {
	name = "freeradius"
	prefix = "/usr"
	localstatedir = "/var"
	sbindir = "/usr/sbin"
	logdir = "/var/log/freeradius"
	run_dir = "/var/run/freeradius"
	libdir = "/usr/lib/freeradius"
	radacctdir = "/var/log/freeradius/radacct"
	hostname_lookups = no
	max_request_time = 30
	proxy_dedup_window = 1
	cleanup_delay = 5
	max_requests = 16384
	max_fds = 512
	postauth_client_lost = no
	pidfile = "/var/run/freeradius/freeradius.pid"
	checkrad = "/usr/sbin/checkrad"
	debug_level = 0
	proxy_requests = yes
 log {
 	stripped_names = no
 	auth = no
 	auth_badpass = no
 	auth_goodpass = no
 	colourise = yes
 	msg_denied = "You are already logged in - access denied"
 }
 resources {
 }
 security {
 	max_attributes = 200
 	reject_delay = 1.000000
 	status_server = yes
 	require_message_authenticator = "auto"
 	limit_proxy_state = "auto"
 }
}
radiusd: #### Loading Realms and Home Servers ####
 proxy server {
 	retry_delay = 5
 	retry_count = 3
 	default_fallback = no
 	dead_time = 120
 	wake_all_if_all_dead = no
 }
 home_server localhost {
 	nonblock = no
 	ipaddr = 127.0.0.1
 	port = 1812
 	type = "auth"
 	secret = <<< secret >>>
 	response_window = 20.000000
 	response_timeouts = 1
 	max_outstanding = 65536
 	zombie_period = 40
 	status_check = "status-server"
 	ping_interval = 30
 	check_interval = 30
 	check_timeout = 4
 	num_answers_to_alive = 3
 	revive_interval = 120
  limit {
  	max_connections = 16
  	max_requests = 0
  	lifetime = 0
  	idle_timeout = 0
  }
  coa {
  	irt = 2
  	mrt = 16
  	mrc = 5
  	mrd = 30
  }
 }
 home_server_pool my_auth_failover {
	type = fail-over
	home_server = localhost
 }
 realm example.com {
	auth_pool = my_auth_failover
 }
 realm LOCAL {
 }
radiusd: #### Loading Clients ####
 client localhost {
 	ipaddr = 127.0.0.1
 	secret = <<< secret >>>
 	nas_type = "other"
 	proto = "*"
  limit {
  	max_connections = 16
  	lifetime = 0
  	idle_timeout = 30
  }
 }
 client localhost_ipv6 {
 	ipv6addr = ::1
 	secret = <<< secret >>>
  limit {
  	max_connections = 16
  	lifetime = 0
  	idle_timeout = 30
  }
 }
 client access-node {
 	ipaddr = 10.0.23.1
 	secret = <<< secret >>>
 	shortname = "access-node"
  limit {
  	max_connections = 16
  	lifetime = 0
  	idle_timeout = 30
  }
 }
Debugger not attached
systemd watchdog is disabled
 # Creating Auth-Type = mschap
 # Creating Auth-Type = eap
 # Creating Auth-Type = PAP
 # Creating Auth-Type = CHAP
 # Creating Auth-Type = MS-CHAP
 # Creating Auth-Type = digest
 # Creating Autz-Type = New-TLS-Connection
radiusd: #### Instantiating modules ####
 modules {
  # Loaded module rlm_always
  # Loading module "reject" from file /etc/freeradius/3.0/mods-enabled/always
  always reject {
  	rcode = "reject"
  	simulcount = 0
  	mpp = no
  }
  # Loading module "fail" from file /etc/freeradius/3.0/mods-enabled/always
  always fail {
  	rcode = "fail"
  	simulcount = 0
  	mpp = no
  }
  # Loading module "ok" from file /etc/freeradius/3.0/mods-enabled/always
  always ok {
  	rcode = "ok"
  	simulcount = 0
  	mpp = no
  }
  # Loading module "handled" from file /etc/freeradius/3.0/mods-enabled/always
  always handled {
  	rcode = "handled"
  	simulcount = 0
  	mpp = no
  }
  # Loading module "invalid" from file /etc/freeradius/3.0/mods-enabled/always
  always invalid {
  	rcode = "invalid"
  	simulcount = 0
  	mpp = no
  }
  # Loading module "userlock" from file /etc/freeradius/3.0/mods-enabled/always
  always userlock {
  	rcode = "userlock"
  	simulcount = 0
  	mpp = no
  }
  # Loading module "notfound" from file /etc/freeradius/3.0/mods-enabled/always
  always notfound {
  	rcode = "notfound"
  	simulcount = 0
  	mpp = no
  }
  # Loading module "noop" from file /etc/freeradius/3.0/mods-enabled/always
  always noop {
  	rcode = "noop"
  	simulcount = 0
  	mpp = no
  }
  # Loading module "updated" from file /etc/freeradius/3.0/mods-enabled/always
  always updated {
  	rcode = "updated"
  	simulcount = 0
  	mpp = no
  }
  # Loaded module rlm_mschap
  # Loading module "mschap" from file /etc/freeradius/3.0/mods-enabled/mschap
  mschap {
  	use_mppe = yes
  	require_encryption = no
  	require_strong = no
  	with_ntdomain_hack = yes
   passchange {
   }
  	allow_retry = yes
  	winbind_retry_with_normalised_username = no
  }
  # Loaded module rlm_radutmp
  # Loading module "radutmp" from file /etc/freeradius/3.0/mods-enabled/radutmp
  radutmp {
  	filename = "/var/log/freeradius/radutmp"
  	username = "%{User-Name}"
  	case_sensitive = yes
  	check_with_nas = yes
  	permissions = 384
  	caller_id = yes
  }
  # Loaded module rlm_exec
  # Loading module "exec" from file /etc/freeradius/3.0/mods-enabled/exec
  exec {
  	wait = no
  	input_pairs = "request"
  	shell_escape = yes
  	timeout = 10
  }
  # Loading module "sradutmp" from file /etc/freeradius/3.0/mods-enabled/sradutmp
  radutmp sradutmp {
  	filename = "/var/log/freeradius/sradutmp"
  	username = "%{User-Name}"
  	case_sensitive = yes
  	check_with_nas = yes
  	permissions = 420
  	caller_id = no
  }
  # Loading module "ntlm_auth" from file /etc/freeradius/3.0/mods-enabled/ntlm_auth
  exec ntlm_auth {
  	wait = yes
  	program = "/path/to/ntlm_auth --request-nt-key --domain=MYDOMAIN --username=%{mschap:Us
        er-Name} --password=%{User-Password}"
  	shell_escape = yes
  }
  # Loaded module rlm_preprocess
  # Loading module "preprocess" from file /etc/freeradius/3.0/mods-enabled/preprocess
  preprocess {
  	huntgroups = "/etc/freeradius/3.0/mods-config/preprocess/huntgroups"
  	hints = "/etc/freeradius/3.0/mods-config/preprocess/hints"
  	with_ascend_hack = no
  	ascend_channels_per_line = 23
  	with_ntdomain_hack = no
  	with_specialix_jetstream_hack = no
  	with_cisco_vsa_hack = no
  	with_alvarion_vsa_hack = no
  }
  # Loaded module rlm_detail
  # Loading module "detail" from file /etc/freeradius/3.0/mods-enabled/detail
  detail {
  	filename = "/var/log/freeradius/radacct/%{%{Packet-Src-IP-Address}:-%{Packet-Src-IPv6-A
        ddress}}/detail-%Y%m%d"
  	header = "%t"
  	permissions = 384
  	locking = no
  	dates_as_integer = no
  	escape_filenames = no
  	log_packet_header = no
  }
  # Loaded module rlm_soh
  # Loading module "soh" from file /etc/freeradius/3.0/mods-enabled/soh
  soh {
  	dhcp = yes
  }
  # Loading module "echo" from file /etc/freeradius/3.0/mods-enabled/echo
  exec echo {
  	wait = yes
  	program = "/bin/echo %{User-Name}"
  	input_pairs = "request"
  	output_pairs = "reply"
  	shell_escape = yes
  }
  # Loaded module rlm_expr
  # Loading module "expr" from file /etc/freeradius/3.0/mods-enabled/expr
  expr {
  	safe_characters = "@abcdefghijklmnopqrstuvwxyzABCDEFGHIJKLMNOPQRSTUVWXYZ0123456789.-_: 
        /äéöüàâæçèéêëîïôœùûüaÿÄÉÖÜßÀÂÆÇÈÉÊËÎÏÔŒÙÛÜŸ"
  }
  # Loaded module rlm_unpack
  # Loading module "unpack" from file /etc/freeradius/3.0/mods-enabled/unpack
  # Loaded module rlm_logintime
  # Loading module "logintime" from file /etc/freeradius/3.0/mods-enabled/logintime
  logintime {
  	minimum_timeout = 60
  }
  # Loaded module rlm_realm
  # Loading module "IPASS" from file /etc/freeradius/3.0/mods-enabled/realm
  realm IPASS {
  	format = "prefix"
  	delimiter = "/"
  	ignore_default = no
  	ignore_null = no
  }
  # Loading module "suffix" from file /etc/freeradius/3.0/mods-enabled/realm
  realm suffix {
  	format = "suffix"
  	delimiter = "@"
  	ignore_default = no
  	ignore_null = no
  }
  # Loading module "bangpath" from file /etc/freeradius/3.0/mods-enabled/realm
  realm bangpath {
  	format = "prefix"
  	delimiter = "!"
  	ignore_default = no
  	ignore_null = no
  }
  # Loading module "realmpercent" from file /etc/freeradius/3.0/mods-enabled/realm
  realm realmpercent {
  	format = "suffix"
  	delimiter = "%"
  	ignore_default = no
  	ignore_null = no
  }
  # Loading module "ntdomain" from file /etc/freeradius/3.0/mods-enabled/realm
  realm ntdomain {
  	format = "prefix"
  	delimiter = "\"
  	ignore_default = no
  	ignore_null = no
  }
  # Loaded module rlm_attr_filter
  # Loading module "attr_filter.post-proxy" from file /etc/freeradius/3.0/mods-enabled/
          attr_filter
  attr_filter attr_filter.post-proxy {
  	filename = "/etc/freeradius/3.0/mods-config/attr_filter/post-proxy"
  	key = "%{Realm}"
  	relaxed = no
  }
  # Loading module "attr_filter.pre-proxy" from file /etc/freeradius/3.0/mods-enabled/
          attr_filter
  attr_filter attr_filter.pre-proxy {
  	filename = "/etc/freeradius/3.0/mods-config/attr_filter/pre-proxy"
  	key = "%{Realm}"
  	relaxed = no
  }
  # Loading module "attr_filter.access_reject" from file /etc/freeradius/3.0/
          mods-enabled/attr_filter
  attr_filter attr_filter.access_reject {
  	filename = "/etc/freeradius/3.0/mods-config/attr_filter/access_reject"
  	key = "%{User-Name}"
  	relaxed = no
  }
  # Loading module "attr_filter.access_challenge" from file /etc/freeradius/3.0/
          mods-enabled/attr_filter
  attr_filter attr_filter.access_challenge {
  	filename = "/etc/freeradius/3.0/mods-config/attr_filter/access_challenge"
  	key = "%{User-Name}"
  	relaxed = no
  }
  # Loading module "attr_filter.accounting_response" from file /etc/freeradius/3.0/
          mods-enabled/attr_filter
  attr_filter attr_filter.accounting_response {
  	filename = "/etc/freeradius/3.0/mods-config/attr_filter/accounting_response"
  	key = "%{User-Name}"
  	relaxed = no
  }
  # Loading module "attr_filter.coa" from file /etc/freeradius/3.0/
          mods-enabled/attr_filter
  attr_filter attr_filter.coa {
  	filename = "/etc/freeradius/3.0/mods-config/attr_filter/coa"
  	key = "%{User-Name}"
  	relaxed = no
  }
  # Loaded module rlm_digest
  # Loading module "digest" from file /etc/freeradius/3.0/mods-enabled/digest
  # Loaded module rlm_expiration
  # Loading module "expiration" from file /etc/freeradius/3.0/mods-enabled/expiration
  # Loaded module rlm_utf8
  # Loading module "utf8" from file /etc/freeradius/3.0/mods-enabled/utf8
  # Loaded module rlm_chap
  # Loading module "chap" from file /etc/freeradius/3.0/mods-enabled/chap
  # Loaded module rlm_pap
  # Loading module "pap" from file /etc/freeradius/3.0/mods-enabled/pap
  pap {
  	normalise = yes
  }
  # Loaded module rlm_unix
  # Loading module "unix" from file /etc/freeradius/3.0/mods-enabled/unix
  unix {
  	radwtmp = "/var/log/freeradius/radwtmp"
  }
Creating attribute Unix-Group
  # Loaded module rlm_passwd
  # Loading module "etc_passwd" from file /etc/freeradius/3.0/mods-enabled/passwd
  passwd etc_passwd {
  	filename = "/etc/passwd"
  	format = "*User-Name:Crypt-Password:"
  	delimiter = ":"
  	ignore_nislike = no
  	ignore_empty = yes
  	allow_multiple_keys = no
  	hash_size = 100
  }
  # Loading module "auth_log" from file /etc/freeradius/3.0/mods-enabled/detail.log
  detail auth_log {
  	filename = "/var/log/freeradius/radacct/%{%{Packet-Src-IP-Address}:-%{Packet-Src-IPv6-A
        ddress}}/auth-detail-%Y%m%d"
  	header = "%t"
  	permissions = 384
  	locking = no
  	dates_as_integer = no
  	escape_filenames = no
  	log_packet_header = no
  }
  # Loading module "reply_log" from file /etc/freeradius/3.0/mods-enabled/detail.log
  detail reply_log {
  	filename = "/var/log/freeradius/radacct/%{%{Packet-Src-IP-Address}:-%{Packet-Src-IPv6-A
        ddress}}/reply-detail-%Y%m%d"
  	header = "%t"
  	permissions = 384
  	locking = no
  	dates_as_integer = no
  	escape_filenames = no
  	log_packet_header = no
  }
  # Loading module "pre_proxy_log" from file /etc/freeradius/3.0/mods-enabled/detail.log
  detail pre_proxy_log {
  	filename = "/var/log/freeradius/radacct/%{%{Packet-Src-IP-Address}:-%{Packet-Src-IPv6-A
        ddress}}/pre-proxy-detail-%Y%m%d"
  	header = "%t"
  	permissions = 384
  	locking = no
  	dates_as_integer = no
  	escape_filenames = no
  	log_packet_header = no
  }
  # Loading module "post_proxy_log" from file /etc/freeradius/3.0/mods-enabled/detail.log
  detail post_proxy_log {
  	filename = "/var/log/freeradius/radacct/%{%{Packet-Src-IP-Address}:-%{Packet-Src-IPv6-A
        ddress}}/post-proxy-detail-%Y%m%d"
  	header = "%t"
  	permissions = 384
  	locking = no
  	dates_as_integer = no
  	escape_filenames = no
  	log_packet_header = no
  }
  # Loaded module rlm_eap
  # Loading module "eap" from file /etc/freeradius/3.0/mods-enabled/eap
  eap {
  	default_eap_type = "md5"
  	timer_expire = 60
  	max_eap_type = 52
  	ignore_unknown_eap_types = no
  	cisco_accounting_username_bug = no
  	max_sessions = 16384
  	dedup_key = ""
  }
  # Loaded module rlm_replicate
  # Loading module "replicate" from file /etc/freeradius/3.0/mods-enabled/replicate
  # Loaded module rlm_linelog
  # Loading module "linelog" from file /etc/freeradius/3.0/mods-enabled/linelog
  linelog {
  	filename = "/var/log/freeradius/linelog"
  	escape_filenames = no
  	syslog_severity = "info"
  	permissions = 384
  	format = "This is a log message for %{User-Name}"
  	reference = "messages.%{%{reply:Packet-Type}:-default}"
  }
  # Loading module "log_accounting" from file /etc/freeradius/3.0/mods-enabled/linelog
  linelog log_accounting {
  	filename = "/var/log/freeradius/linelog-accounting"
  	escape_filenames = no
  	syslog_severity = "info"
  	permissions = 384
  	format = ""
  	reference = "Accounting-Request.%{%{Acct-Status-Type}:-unknown}"
  }
  # Loaded module rlm_dynamic_clients
  # Loading module "dynamic_clients" from file /etc/freeradius/3.0/mods-enabled/dynamic_cl
        ients
  # Loaded module rlm_files
  # Loading module "files" from file /etc/freeradius/3.0/mods-enabled/files
  files {
  	filename = "/etc/freeradius/3.0/mods-config/files/authorize"
  	acctusersfile = "/etc/freeradius/3.0/mods-config/files/accounting"
  	preproxy_usersfile = "/etc/freeradius/3.0/mods-config/files/pre-proxy"
  }
  instantiate {
  }
  # Instantiating module "reject" from file /etc/freeradius/3.0/mods-enabled/always
  # Instantiating module "fail" from file /etc/freeradius/3.0/mods-enabled/always
  # Instantiating module "ok" from file /etc/freeradius/3.0/mods-enabled/always
  # Instantiating module "handled" from file /etc/freeradius/3.0/mods-enabled/always
  # Instantiating module "invalid" from file /etc/freeradius/3.0/mods-enabled/always
  # Instantiating module "userlock" from file /etc/freeradius/3.0/mods-enabled/always
  # Instantiating module "notfound" from file /etc/freeradius/3.0/mods-enabled/always
  # Instantiating module "noop" from file /etc/freeradius/3.0/mods-enabled/always
  # Instantiating module "updated" from file /etc/freeradius/3.0/mods-enabled/always
  # Instantiating module "mschap" from file /etc/freeradius/3.0/mods-enabled/mschap
rlm_mschap (mschap): using internal authentication
  # Instantiating module "preprocess" from file /etc/freeradius/3.0/mods-enabled/
        preprocess
reading pairlist file /etc/freeradius/3.0/mods-config/preprocess/huntgroups
reading pairlist file /etc/freeradius/3.0/mods-config/preprocess/hints
  # Instantiating module "detail" from file /etc/freeradius/3.0/mods-enabled/detail
  # Instantiating module "logintime" from file /etc/freeradius/3.0/mods-enabled/logintime
  # Instantiating module "IPASS" from file /etc/freeradius/3.0/mods-enabled/realm
  # Instantiating module "suffix" from file /etc/freeradius/3.0/mods-enabled/realm
  # Instantiating module "bangpath" from file /etc/freeradius/3.0/mods-enabled/realm
  # Instantiating module "realmpercent" from file /etc/freeradius/3.0/mods-enabled/realm
  # Instantiating module "ntdomain" from file /etc/freeradius/3.0/mods-enabled/realm
  # Instantiating module "attr_filter.post-proxy" from file /etc/freeradius/3.0/mods-enabl
        ed/attr_filter
reading pairlist file /etc/freeradius/3.0/mods-config/attr_filter/post-proxy
  # Instantiating module "attr_filter.pre-proxy" from file /etc/freeradius/3.0/mods-enable
        d/attr_filter
reading pairlist file /etc/freeradius/3.0/mods-config/attr_filter/pre-proxy
  # Instantiating module "attr_filter.access_reject" from file /etc/freeradius/3.0/mods-en
        abled/attr_filter
reading pairlist file /etc/freeradius/3.0/mods-config/attr_filter/access_reject
  # Instantiating module "attr_filter.access_challenge" from file /etc/freeradius/3.0/mods
        -enabled/attr_filter
reading pairlist file /etc/freeradius/3.0/mods-config/attr_filter/access_challenge
  # Instantiating module "attr_filter.accounting_response" from file /etc/freeradius/3.0/m
        ods-enabled/attr_filter
reading pairlist file /etc/freeradius/3.0/mods-config/attr_filter/accounting_response
  # Instantiating module "attr_filter.coa" from file /etc/freeradius/3.0/mods-enabled/attr
        _filter
reading pairlist file /etc/freeradius/3.0/mods-config/attr_filter/coa
  # Instantiating module "expiration" from file /etc/freeradius/3.0/mods-enabled/expiratio
        n
  # Instantiating module "pap" from file /etc/freeradius/3.0/mods-enabled/pap
  # Instantiating module "etc_passwd" from file /etc/freeradius/3.0/mods-enabled/passwd
rlm_passwd: nfields: 3 keyfield 0(User-Name) listable: no
  # Instantiating module "auth_log" from file /etc/freeradius/3.0/mods-enabled/detail.log
rlm_detail (auth_log): 'User-Password' suppressed, will not appear in detail output
  # Instantiating module "reply_log" from file /etc/freeradius/3.0/mods-enabled/detail.log
  # Instantiating module "pre_proxy_log" from file /etc/freeradius/3.0/mods-enabled/detail
        .log
  # Instantiating module "post_proxy_log" from file /etc/freeradius/3.0/mods-enabled/detai
        l.log
  # Instantiating module "eap" from file /etc/freeradius/3.0/mods-enabled/eap
   # Linked to sub-module rlm_eap_md5
   # Linked to sub-module rlm_eap_gtc
   gtc {
   	challenge = "Password: "
   	auth_type = "PAP"
   }
   # Linked to sub-module rlm_eap_tls
   tls {
   	tls = "tls-common"
   }
   tls-config tls-common {
   	verify_depth = 0
   	ca_path = "/etc/freeradius/3.0/certs"
   	pem_file_type = yes
   	private_key_file = "/etc/ssl/private/ssl-cert-snakeoil.key"
   	certificate_file = "/etc/ssl/certs/ssl-cert-snakeoil.pem"
   	ca_file = "/etc/ssl/certs/ca-certificates.crt"
   	private_key_password = <<< secret >>>
   	fragment_size = 1024
   	include_length = yes
   	auto_chain = yes
   	check_crl = no
   	check_all_crl = no
   	ca_path_reload_interval = 0
   	cipher_list = "DEFAULT"
   	cipher_server_preference = no
   	reject_unknown_intermediate_ca = no
   	ecdh_curve = ""
   	tls_max_version = "1.2"
   	tls_min_version = "1.2"
    cache {
    	enable = no
    	lifetime = 24
    	max_entries = 255
    }
    verify {
    	skip_if_ocsp_ok = no
    }
    ocsp {
    	enable = no
    	override_cert_url = yes
    	url = "http://127.0.0.1/ocsp/"
    	use_nonce = yes
    	timeout = 0
    	softfail = no
    }
   }
   # Linked to sub-module rlm_eap_ttls
   ttls {
   	tls = "tls-common"
   	default_eap_type = "md5"
   	copy_request_to_tunnel = no
   	use_tunneled_reply = no
   	virtual_server = "inner-tunnel"
   	include_length = yes
   	require_client_cert = no
   }
tls: Using cached TLS configuration from previous invocation
   # Linked to sub-module rlm_eap_peap
   peap {
   	tls = "tls-common"
   	default_eap_type = "mschapv2"
   	copy_request_to_tunnel = no
   	use_tunneled_reply = no
   	proxy_tunneled_request_as_eap = yes
   	virtual_server = "inner-tunnel"
   	soh = no
   	require_client_cert = no
   }
tls: Using cached TLS configuration from previous invocation
   # Linked to sub-module rlm_eap_mschapv2
   mschapv2 {
   	with_ntdomain_hack = no
   	send_error = no
   }
  # Instantiating module "linelog" from file /etc/freeradius/3.0/mods-enabled/linelog
  # Instantiating module "log_accounting" from file /etc/freeradius/3.0/mods-enabled/linel
        og
  # Instantiating module "files" from file /etc/freeradius/3.0/mods-enabled/files
reading pairlist file /etc/freeradius/3.0/mods-config/files/authorize
reading pairlist file /etc/freeradius/3.0/mods-config/files/accounting
reading pairlist file /etc/freeradius/3.0/mods-config/files/pre-proxy
 } # modules
radiusd: #### Loading Virtual Servers ####
server { # from file /etc/freeradius/3.0/radiusd.conf
} # server
server inner-tunnel { # from file /etc/freeradius/3.0/sites-enabled/inner-tunnel
 # Loading authenticate {...}
Compiling Auth-Type PAP for attr Auth-Type
Compiling Auth-Type CHAP for attr Auth-Type
Compiling Auth-Type MS-CHAP for attr Auth-Type
 # Loading authorize {...}
Ignoring "sql" (see raddb/mods-available/README.rst)
Ignoring "ldap" (see raddb/mods-available/README.rst)
 # Loading session {...}
 # Loading post-proxy {...}
 # Loading post-auth {...}
 # Skipping contents of 'if' as it is always 'false' -- /etc/freeradius/3.0/sites-enabled/
        inner-tunnel:366
Compiling Post-Auth-Type REJECT for attr Post-Auth-Type
} # server inner-tunnel
server default { # from file /etc/freeradius/3.0/sites-enabled/default
 # Loading authenticate {...}
Compiling Auth-Type PAP for attr Auth-Type
Compiling Auth-Type CHAP for attr Auth-Type
Compiling Auth-Type MS-CHAP for attr Auth-Type
 # Loading authorize {...}
Compiling Autz-Type New-TLS-Connection for attr Autz-Type
 # Loading preacct {...}
 # Loading accounting {...}
 # Loading post-proxy {...}
 # Loading post-auth {...}
Compiling Post-Auth-Type REJECT for attr Post-Auth-Type
Compiling Post-Auth-Type Challenge for attr Post-Auth-Type
Compiling Post-Auth-Type Client-Lost for attr Post-Auth-Type
} # server default
radiusd: #### Opening IP addresses and Ports ####
listen {
  	type = "auth"
  	ipaddr = 127.0.0.1
  	port = 18120
}
listen {
  	type = "auth"
  	ipaddr = *
  	port = 0
   limit {
   	max_connections = 16
   	lifetime = 0
   	idle_timeout = 30
   }
}
listen {
  	type = "acct"
  	ipaddr = *
  	port = 0
   limit {
   	max_connections = 16
   	lifetime = 0
   	idle_timeout = 30
   }
}
listen {
  	type = "auth"
  	ipv6addr = ::
  	port = 0
   limit {
   	max_connections = 16
   	lifetime = 0
   	idle_timeout = 30
   }
}
listen {
  	type = "acct"
  	ipv6addr = ::
  	port = 0
   limit {
   	max_connections = 16
   	lifetime = 0
   	idle_timeout = 30
   }
}
Listening on auth address 127.0.0.1 port 18120 bound to server inner-tunnel
Listening on auth address * port 1812 bound to server default
Listening on acct address * port 1813 bound to server default
Listening on auth address :: port 1812 bound to server default
Listening on acct address :: port 1813 bound to server default
Listening on proxy address * port 48578
Listening on proxy address :: port 42007
Ready to process requests
(0) Received Access-Request Id 0 from 10.0.23.1:55171 to 10.0.30.10:1812 length 165
(0)   User-Name = "nacuser"
(0)   NAS-IP-Address = 10.0.23.1
(0)   Called-Station-Id = "<REDACTED>:"
(0)   NAS-Port-Type = Wireless-802.11
(0)   Service-Type = Framed-User
(0)   Calling-Station-Id = "<REDACTED>"
(0)   Connect-Info = "CONNECT 0Mbps 802.11b"
(0)   Acct-Session-Id = "1A082DB0BEE2BD34"
(0)   Framed-MTU = 1400
(0)   EAP-Message = 0x0219000c016e616375736572
(0)   Message-Authenticator = 0x71423639f013148f1ce7a71576679bbe
(0) # Executing section authorize from file /etc/freeradius/3.0/sites-enabled/default
(0)   authorize {
(0)     policy filter_username {
(0)       if (&User-Name) {
(0)       if (&User-Name)  -> TRUE
(0)       if (&User-Name)  {
(0)         if (&User-Name =~ / /) {
(0)         if (&User-Name =~ / /)  -> FALSE
(0)         if (&User-Name =~ /@[^@]*@/ ) {
(0)         if (&User-Name =~ /@[^@]*@/ )  -> FALSE
(0)         if (&User-Name =~ /\.\./ ) {
(0)         if (&User-Name =~ /\.\./ )  -> FALSE
(0)         if ((&User-Name =~ /@/) && (&User-Name !~ /@(.+)\.(.+)$/))  {
(0)         if ((&User-Name =~ /@/) && (&User-Name !~ /@(.+)\.(.+)$/))   -> FALSE
(0)         if (&User-Name =~ /\.$/)  {
(0)         if (&User-Name =~ /\.$/)   -> FALSE
(0)         if (&User-Name =~ /@\./)  {
(0)         if (&User-Name =~ /@\./)   -> FALSE
(0)       } # if (&User-Name)  = notfound
(0)     } # policy filter_username = notfound
(0)     [preprocess] = ok
(0)     [chap] = noop
(0)     [mschap] = noop
(0)     [digest] = noop
(0) suffix: Checking for suffix after "@"
(0) suffix: No '@' in User-Name = "nacuser", looking up realm NULL
(0) suffix: No such realm "NULL"
(0)     [suffix] = noop
(0) eap: Peer sent EAP Response (code 2) ID 25 length 12
(0) eap: EAP-Identity reply, returning 'ok' so we can short-circuit the rest of authorize
(0)     [eap] = ok
(0)   } # authorize = ok
(0) Found Auth-Type = eap
(0) # Executing group from file /etc/freeradius/3.0/sites-enabled/default
(0)   authenticate {
(0) eap: Peer sent packet with method EAP Identity (1)
(0) eap: Calling submodule eap_md5 to process data
(0) eap_md5: Issuing MD5 Challenge
(0) eap: Sending EAP Request (code 1) ID 26 length 22
(0) eap: EAP session adding &reply:State = 0xbb99b235bb83b6a0
(0)     [eap] = handled
(0)   } # authenticate = handled
(0) Using Post-Auth-Type Challenge
(0) # Executing group from file /etc/freeradius/3.0/sites-enabled/default
(0)   Challenge { ... } # empty sub-section is ignored
(0) Sent Access-Challenge Id 0 from 10.0.30.10:1812 to 10.0.23.1:55171 length 80
(0)   EAP-Message = 0x011a00160410d8d1ca688796282203d8a54492bb84f9
(0)   Message-Authenticator = 0x00000000000000000000000000000000
(0)   State = 0xbb99b235bb83b6a0b74fb036b1c91a02
(0) Finished request
Waking up in 4.9 seconds.
(1) Received Access-Request Id 1 from 10.0.23.1:55171 to 10.0.30.10:1812 length 193
(1)   User-Name = "nacuser"
(1)   NAS-IP-Address = 10.0.23.1
(1)   Called-Station-Id = "<REDACTED>:"
(1)   NAS-Port-Type = Wireless-802.11
(1)   Service-Type = Framed-User
(1)   Calling-Station-Id = "<REDACTED>"
(1)   Connect-Info = "CONNECT 0Mbps 802.11b"
(1)   Acct-Session-Id = "1A082DB0BEE2BD34"
(1)   Framed-MTU = 1400
(1)   EAP-Message = 0x021a0016041041b2c01ccfa94d12fc491a04d95cd935
(1)   State = 0xbb99b235bb83b6a0b74fb036b1c91a02
(1)   Message-Authenticator = 0xee00aa03a7391d212502d3e4595c510c
(1) session-state: No cached attributes
(1) # Executing section authorize from file /etc/freeradius/3.0/sites-enabled/default
(1)   authorize {
(1)     policy filter_username {
(1)       if (&User-Name) {
(1)       if (&User-Name)  -> TRUE
(1)       if (&User-Name)  {
(1)         if (&User-Name =~ / /) {
(1)         if (&User-Name =~ / /)  -> FALSE
(1)         if (&User-Name =~ /@[^@]*@/ ) {
(1)         if (&User-Name =~ /@[^@]*@/ )  -> FALSE
(1)         if (&User-Name =~ /\.\./ ) {
(1)         if (&User-Name =~ /\.\./ )  -> FALSE
(1)         if ((&User-Name =~ /@/) && (&User-Name !~ /@(.+)\.(.+)$/))  {
(1)         if ((&User-Name =~ /@/) && (&User-Name !~ /@(.+)\.(.+)$/))   -> FALSE
(1)         if (&User-Name =~ /\.$/)  {
(1)         if (&User-Name =~ /\.$/)   -> FALSE
(1)         if (&User-Name =~ /@\./)  {
(1)         if (&User-Name =~ /@\./)   -> FALSE
(1)       } # if (&User-Name)  = notfound
(1)     } # policy filter_username = notfound
(1)     [preprocess] = ok
(1)     [chap] = noop
(1)     [mschap] = noop
(1)     [digest] = noop
(1) suffix: Checking for suffix after "@"
(1) suffix: No '@' in User-Name = "nacuser", looking up realm NULL
(1) suffix: No such realm "NULL"
(1)     [suffix] = noop
(1) eap: Peer sent EAP Response (code 2) ID 26 length 22
(1) eap: No EAP Start, assuming it's an on-going EAP conversation
(1)     [eap] = updated
(1) files: users: Matched entry nacuser at line 207
(1)     [files] = ok
(1)     [expiration] = noop
(1)     [logintime] = noop
(1) pap: WARNING: Auth-Type already set.  Not setting to PAP
(1)     [pap] = noop
(1)   } # authorize = updated
(1) Found Auth-Type = eap
(1) # Executing group from file /etc/freeradius/3.0/sites-enabled/default
(1)   authenticate {
(1) eap: Removing EAP session with state 0xbb99b235bb83b6a0
(1) eap: Previous EAP request found for state 0xbb99b235bb83b6a0, released from the list
(1) eap: Peer sent packet with method EAP MD5 (4)
(1) eap: Calling submodule eap_md5 to process data
EAP-MD5 digests do not match.
(1) eap: Sending EAP Failure (code 4) ID 26 length 4
(1) eap: Freeing handler
(1)     [eap] = reject
(1)   } # authenticate = reject
(1) Failed to authenticate the user
(1) Using Post-Auth-Type Reject
(1) # Executing group from file /etc/freeradius/3.0/sites-enabled/default
(1)   Post-Auth-Type REJECT {
(1) attr_filter.access_reject: EXPAND %{User-Name}
(1) attr_filter.access_reject:    --> nacuser
(1) attr_filter.access_reject: Matched entry DEFAULT at line 11
(1)     [attr_filter.access_reject] = updated
(1)     [eap] = noop
(1)     policy remove_reply_message_if_eap {
(1)       if (&reply:EAP-Message && &reply:Reply-Message) {
(1)       if (&reply:EAP-Message && &reply:Reply-Message)  -> FALSE
(1)       else {
(1)         [noop] = noop
(1)       } # else = noop
(1)     } # policy remove_reply_message_if_eap = noop
(1)   } # Post-Auth-Type REJECT = updated
(1) Delaying response for 1.000000 seconds
Waking up in 0.3 seconds.
Waking up in 0.1 seconds.
Signalled to terminate
Exiting normally
\end{verbatim}

% ---------------------------------------------------------------------
% SUBSECTION B.1.4: WPA_SUPPLICANT SUCCESS RE-TEST LOG
% ---------------------------------------------------------------------
% This subsection preserves the raw successful remediation re-test
% capture taken from the client-side manual wpa_supplicant run.
% ---------------------------------------------------------------------

\subsubsection{B.1.4 wpa\_supplicant Success Re-Test Log}
\label{sec:appendix-b-wpa-supplicant-success-retest-log}

\begin{verbatim}
radius-splunk: ~/incident-evidence/wpa_supplicant_success_retest_v2.txt
wpa_supplicant v2.10
random: getrandom() support available
Successfully initialized wpa_supplicant
Initializing interface 'ens33' conf '/etc/wpa_supplicant/lab/wired-8021x.conf' driver 'wir
        ed' ctrl_interface 'N/A' bridge 'N/A'
Configuration file '/etc/wpa_supplicant/lab/wired-8021x.conf' -> '/etc/wpa_supplicant/lab/
        wired-8021x.conf'
Reading configuration file '/etc/wpa_supplicant/lab/wired-8021x.conf'
ctrl_interface='/run/wpa_supplicant'
ap_scan=0
eapol_version=2
Line: 5 - start of a new network block
key_mgmt: 0x8
eap methods - hexdump(len=16): 00 00 00 00 04 00 00 00 00 00 00 00 00 00 00 00
identity - hexdump_ascii(len=7):
     6e 61 63 75 73 65 72                              nacuser         
password - hexdump_ascii(len=11): [REMOVED]
Priority group 0
   id=0 ssid=''
driver_wired_init_common: Added multicast membership with packet socket
Add interface ens33 to a new radio N/A
ens33: Own MAC address: <REDACTED>
ens33: RSN: flushing PMKID list in the driver
ens33: Setting scan request: 0.100000 sec
TDLS: TDLS operation not supported by driver
TDLS: Driver uses internal link setup
TDLS: Driver does not support TDLS channel switching
ens33: WPS: UUID based on MAC address: f6fbf57e-3bcd-5a3f-9f55-30679786050a
ENGINE: Loading builtin engines
ENGINE: Loading builtin engines
EAPOL: SUPP_PAE entering state DISCONNECTED
EAPOL: Supplicant port status: Unauthorized
EAPOL: KEY_RX entering state NO_KEY_RECEIVE
EAPOL: SUPP_BE entering state INITIALIZE
EAP: EAP entering state DISABLED
MBO: Update non-preferred channels, non_pref_chan=N/A
ens33: Added interface ens33
ens33: State: DISCONNECTED -> DISCONNECTED
EAPOL: External notification - EAP success=0
EAPOL: External notification - EAP fail=0
EAPOL: External notification - portControl=Auto
ens33: Already associated with a configured network - generating associated event
ens33: Event ASSOC (0) received
ens33: Association info event
ens33: State: DISCONNECTED -> ASSOCIATED
ens33: Associated to a new BSS: BSSID=<REDACTED>
ens33: Select network based on association information
ens33: Network configuration found for the current AP
ens33: WPA: clearing AP WPA IE
ens33: WPA: clearing AP RSN IE
ens33: WPA: clearing AP RSNXE
ens33: WPA: clearing own WPA/RSN IE
ens33: RSN: clearing own RSNXE
ens33: Failed to get scan results
EAPOL: External notification - EAP success=0
EAPOL: External notification - EAP fail=0
EAPOL: External notification - portControl=Auto
ens33: Associated with <REDACTED>
ens33: WPA: Association event - clear replay counter
ens33: WPA: Clear old PTK
TDLS: Remove peers on association
EAPOL: External notification - portEnabled=0
EAPOL: External notification - portValid=0
EAPOL: External notification - portEnabled=1
EAPOL: SUPP_PAE entering state CONNECTING
EAPOL: SUPP_BE entering state IDLE
EAP: EAP entering state INITIALIZE
EAP: EAP entering state IDLE
ens33: Cancelling scan request
ens33: CTRL-EVENT-SUBNET-STATUS-UPDATE status=0
EAPOL: startWhen --> 0
EAPOL: SUPP_PAE entering state CONNECTING
EAPOL: txStart
TX EAPOL: dst=<REDACTED>
TX EAPOL - hexdump(len=4): 02 01 00 00
l2_packet_receive: src=<REDACTED> len=46
ens33: RX EAPOL from <REDACTED>
RX EAPOL - hexdump(len=46): 02 00 00 05 01 fe 00 05 01 00 00 00 00 00 00 00 00 00 00 00 00
         00 00 00 00 00 00 00 00 00 00 00 00 00 00 00 00 00 00 00 00 00 00 00 00 00
EAPOL: Received EAP-Packet frame
EAPOL: SUPP_PAE entering state RESTART
EAP: EAP entering state INITIALIZE
EAP: EAP entering state IDLE
EAPOL: SUPP_PAE entering state AUTHENTICATING
EAPOL: SUPP_BE entering state REQUEST
EAPOL: getSuppRsp
EAP: EAP entering state RECEIVED
EAP: Received EAP-Request id=254 method=1 vendor=0 vendorMethod=0
EAP: EAP entering state IDENTITY
ens33: CTRL-EVENT-EAP-STARTED EAP authentication started
EAP: Status notification: started (param=)
EAP: EAP-Request Identity data - hexdump_ascii(len=0):
EAP: using real identity - hexdump_ascii(len=7):
     6e 61 63 75 73 65 72                              nacuser         
EAP: EAP entering state SEND_RESPONSE
EAP: EAP entering state IDLE
EAPOL: SUPP_BE entering state RESPONSE
EAPOL: txSuppRsp
TX EAPOL: dst=<REDACTED>
TX EAPOL - hexdump(len=16): 02 00 00 0c 02 fe 00 0c 01 6e 61 63 75 73 65 72
EAPOL: SUPP_BE entering state RECEIVE
l2_packet_receive: src=<REDACTED> len=46
ens33: RX EAPOL from <REDACTED>
RX EAPOL - hexdump(len=46): 02 00 00 16 01 ff 00 16 04 10 90 ab 89 1e 63 b1 ce 73 b1 28 33
         17 e8 9c c4 1b 00 00 00 00 00 00 00 00 00 00 00 00 00 00 00 00 00 00 00 00
EAPOL: Received EAP-Packet frame
EAPOL: SUPP_BE entering state REQUEST
EAPOL: getSuppRsp
EAP: EAP entering state RECEIVED
EAP: Received EAP-Request id=255 method=4 vendor=0 vendorMethod=0
EAP: EAP entering state GET_METHOD
ens33: CTRL-EVENT-EAP-PROPOSED-METHOD vendor=0 method=4
EAP: Status notification: accept proposed method (param=MD5)
EAP: Initialize selected EAP method: vendor 0 method 4 (MD5)
ens33: CTRL-EVENT-EAP-METHOD EAP vendor 0 method 4 (MD5) selected
EAP: EAP entering state METHOD
EAP-MD5: Challenge - hexdump(len=16): 90 ab 89 1e 63 b1 ce 73 b1 28 33 17 e8 9c c4 1b
EAP-MD5: Generating Challenge Response
EAP-MD5: Response - hexdump(len=16): 61 59 ca e5 bd 45 91 62 54 d8 5a 8e 12 d8 60 af
EAP: method process -> ignore=FALSE methodState=DONE decision=COND_SUCC eapRespData=0x623c
        dc395420
EAP: EAP entering state SEND_RESPONSE
EAP: EAP entering state IDLE
EAPOL: SUPP_BE entering state RESPONSE
EAPOL: txSuppRsp
TX EAPOL: dst=<REDACTED>
TX EAPOL - hexdump(len=26): 02 00 00 16 02 ff 00 16 04 10 61 59 ca e5 bd 45 91 62 54 d8 5a
         8e 12 d8 60 af
EAPOL: SUPP_BE entering state RECEIVE
l2_packet_receive: src=<REDACTED> len=46
ens33: RX EAPOL from <REDACTED>
RX EAPOL - hexdump(len=46): 02 00 00 04 03 ff 00 04 00 00 00 00 00 00 00 00 00 00 00 00 00
         00 00 00 00 00 00 00 00 00 00 00 00 00 00 00 00 00 00 00 00 00 00 00 00 00
EAPOL: Received EAP-Packet frame
EAPOL: SUPP_BE entering state REQUEST
EAPOL: getSuppRsp
EAP: EAP entering state RECEIVED
EAP: Received EAP-Success
EAP: Status notification: completion (param=success)
EAP: EAP entering state SUCCESS
ens33: CTRL-EVENT-EAP-SUCCESS EAP authentication completed successfully
EAPOL: IEEE 802.1X for plaintext connection; no EAPOL-Key frames required
ens33: WPA: EAPOL processing complete
ens33: Cancelling authentication timeout
ens33: State: ASSOCIATED -> COMPLETED
ens33: CTRL-EVENT-CONNECTED - Connection to <REDACTED> completed [id=0 id_str=]
EAPOL: SUPP_PAE entering state AUTHENTICATED
EAPOL: Supplicant port status: Authorized
EAPOL: SUPP_BE entering state RECEIVE
EAPOL: SUPP_BE entering state SUCCESS
EAPOL: SUPP_BE entering state IDLE
EAPOL authentication completed - result=SUCCESS
ens33: Removing interface ens33
ens33: Request to deauthenticate - bssid=<REDACTED> pending_bssid=00:00:00:00:00:00 reason
        =3 (DEAUTH_LEAVING) state=COMPLETED
TDLS: Tear down peers
ens33: Event DEAUTH (11) received
ens33: Deauthentication notification
ens33:  * reason 3 (DEAUTH_LEAVING) locally_generated=1
Deauthentication frame IE(s) - hexdump(len=0): [NULL]
ens33: CTRL-EVENT-DISCONNECTED bssid=<REDACTED> reason=3 locally_generated=1
ens33: Auto connect disabled: do not try to re-connect
ens33: Ignore connection failure indication since interface has been put into disconnected
         state
TDLS: Remove peers on disassociation
ens33: WPA: Clear old PMK and PTK
ens33: Disconnect event - remove keys
ens33: State: COMPLETED -> DISCONNECTED
EAPOL: External notification - portEnabled=0
EAPOL: SUPP_PAE entering state DISCONNECTED
EAPOL: Supplicant port status: Unauthorized
EAPOL: SUPP_BE entering state INITIALIZE
EAP: EAP entering state DISABLED
EAPOL: External notification - portValid=0
ens33: State: DISCONNECTED -> DISCONNECTED
EAPOL: External notification - portEnabled=0
EAPOL: External notification - portValid=0
QM: Clear all active DSCP policies
ens33: CTRL-EVENT-DSCP-POLICY clear_all
EAP: deinitialize previously used EAP method (4, MD5) at EAP deinit
ens33: WPA: Clear old PMK and PTK
ens33: Cancelling scan request
ens33: Cancelling authentication timeout
Off-channel: Clear pending Action frame TX (pending_action_tx=(nil)
HS20: Delete all stored icons
Off-channel: Action frame sequence done notification: pending_action_tx=(nil) drv_offchan_
        tx=0 action_tx_wait_time=0 off_channel_freq=0 roc_waiting_drv_freq=0
QM: Clear all active DSCP policies
ens33: CTRL-EVENT-DSCP-POLICY clear_all
Remove interface ens33 from radio 
Remove radio 
ens33: CTRL-EVENT-TERMINATING 
\end{verbatim}

% ---------------------------------------------------------------------
% SUBSECTION B.1.5: HOSTAPD SUCCESS RE-TEST LOG
% ---------------------------------------------------------------------
% This subsection preserves the raw successful remediation re-test
% evidence captured from the hostapd side.
% ---------------------------------------------------------------------

\subsubsection{B.1.5 hostapd Success Re-Test Log}
\label{sec:appendix-b-hostapd-success-retest-log}

\begin{verbatim}
radius-splunk: ~/incident-evidence/hostapd_success_retest_v2.txt
Mar 23 02:49:00 access-node hostapd[2052]: Received EAPOL packet
Mar 23 02:49:00 access-node hostapd[2052]: ens38: Event NEW_STA (22) received
Mar 23 02:49:00 access-node hostapd[2052]: Data frame from unknown STA <REDACTED> - adding
         a new STA
Mar 23 02:49:00 access-node hostapd[2052]:   New STA
Mar 23 02:49:00 access-node hostapd[2052]: ap_sta_add: register ap_handle_timer timeout fo
        r <REDACTED> (300 seconds - ap_max_inactivity)
Mar 23 02:49:00 access-node hostapd[2052]: ens38: STA <REDACTED> IEEE 802.1X: start authen
        tication
Mar 23 02:49:00 access-node hostapd[2052]: EAP: Server state machine created
Mar 23 02:49:00 access-node hostapd[2052]: IEEE 802.1X: <REDACTED> BE_AUTH entering state 
        IDLE
Mar 23 02:49:00 access-node hostapd[2052]: IEEE 802.1X: <REDACTED> CTRL_DIR entering state
         FORCE_BOTH
Mar 23 02:49:00 access-node hostapd[2052]: ens38: hostapd_new_assoc_sta: canceled wired ap
        _handle_timer timeout for <REDACTED>
Mar 23 02:49:00 access-node hostapd[2052]: ens38: Event EAPOL_RX (23) received
Mar 23 02:49:00 access-node hostapd[2052]: IEEE 802.1X: 46 bytes from <REDACTED>
Mar 23 02:49:00 access-node hostapd[2052]:    IEEE 802.1X: version=2 type=1 length=0
Mar 23 02:49:00 access-node hostapd[2052]:    ignoring 42 extra octets after IEEE 802.1X p
        acket
Mar 23 02:49:00 access-node hostapd[2052]: ens38: STA <REDACTED> IEEE 802.1X: received EAP
        OL-Start from STA
Mar 23 02:49:00 access-node hostapd[2052]: IEEE 802.1X: <REDACTED> AUTH_PAE entering state
         DISCONNECTED
Mar 23 02:49:00 access-node hostapd[2052]: ens38: STA <REDACTED> IEEE 802.1X: unauthorizin
        g port
Mar 23 02:49:00 access-node hostapd[2052]: IEEE 802.1X: <REDACTED> AUTH_PAE entering state
         RESTART
Mar 23 02:49:00 access-node hostapd[2052]: EAP: EAP entering state INITIALIZE
Mar 23 02:49:00 access-node hostapd[2052]: ens38: CTRL-EVENT-EAP-STARTED <REDACTED>
Mar 23 02:49:00 access-node hostapd[2052]: EAP: EAP entering state SELECT_ACTION
Mar 23 02:49:00 access-node hostapd[2052]: EAP: getDecision: no identity known yet -> CONT
        INUE
Mar 23 02:49:00 access-node hostapd[2052]: EAP: EAP entering state PROPOSE_METHOD
Mar 23 02:49:00 access-node hostapd[2052]: EAP: getNextMethod: vendor 0 type 1
Mar 23 02:49:00 access-node hostapd[2052]: ens38: CTRL-EVENT-EAP-PROPOSED-METHOD vendor=0 
        method=1
Mar 23 02:49:00 access-node hostapd[2052]: EAP: EAP entering state METHOD_REQUEST
Mar 23 02:49:00 access-node hostapd[2052]: EAP: building EAP-Request: Identifier 254
Mar 23 02:49:00 access-node hostapd[2052]: EAP: EAP entering state SEND_REQUEST
Mar 23 02:49:00 access-node hostapd[2052]: EAP: EAP entering state IDLE
Mar 23 02:49:00 access-node hostapd[2052]: EAP: retransmit timeout 3 seconds (from dynamic
         back off; retransCount=0)
Mar 23 02:49:00 access-node hostapd[2052]: IEEE 802.1X: <REDACTED> AUTH_PAE entering state
         CONNECTING
Mar 23 02:49:00 access-node hostapd[2052]: IEEE 802.1X: <REDACTED> AUTH_PAE entering state
         AUTHENTICATING
Mar 23 02:49:00 access-node hostapd[2052]: IEEE 802.1X: <REDACTED> BE_AUTH entering state 
        REQUEST
Mar 23 02:49:00 access-node hostapd[2052]: ens38: STA <REDACTED> IEEE 802.1X: Sending EAP 
        Packet (identifier 254)
Mar 23 02:49:00 access-node hostapd[2052]: Received EAPOL packet
Mar 23 02:49:00 access-node hostapd[2052]: ens38: Event NEW_STA (22) received
Mar 23 02:49:00 access-node hostapd[2052]: ens38: Event EAPOL_RX (23) received
Mar 23 02:49:00 access-node hostapd[2052]: IEEE 802.1X: 46 bytes from <REDACTED>
Mar 23 02:49:00 access-node hostapd[2052]:    IEEE 802.1X: version=2 type=0 length=12
Mar 23 02:49:00 access-node hostapd[2052]:    ignoring 30 extra octets after IEEE 802.1X p
        acket
Mar 23 02:49:00 access-node hostapd[2052]: EAP: code=2 (response) identifier=254 length=12
Mar 23 02:49:00 access-node hostapd[2052]: ens38: STA <REDACTED> IEEE 802.1X: start authen
        tication
Mar 23 02:49:00 access-node hostapd[2052]: ens38: STA <REDACTED> IEEE 802.1X: received EAP
         packet (code=2 id=254 len=12) from STA: EAP Response-Identity (1)
Mar 23 02:49:00 access-node hostapd[2052]: IEEE 802.1X: <REDACTED> BE_AUTH entering state 
        RESPONSE
Mar 23 02:49:00 access-node hostapd[2052]: EAP: EAP entering state RECEIVED
Mar 23 02:49:00 access-node hostapd[2052]: EAP: parseEapResp: rxResp=1 rxInitiate=0 respId
        =254 respMethod=1 respVendor=0 respVendorMethod=0
Mar 23 02:49:00 access-node hostapd[2052]: EAP: EAP entering state INTEGRITY_CHECK
Mar 23 02:49:00 access-node hostapd[2052]: EAP: EAP entering state METHOD_RESPONSE
Mar 23 02:49:00 access-node hostapd[2052]: EAP-Identity: Peer identity - hexdump_ascii(len
        =7):
Mar 23 02:49:00 access-node hostapd[2052]:      6e 61 63 75 73 65 72                      
                nacuser
Mar 23 02:49:00 access-node hostapd[2052]: EAP: EAP entering state SELECT_ACTION
Mar 23 02:49:00 access-node hostapd[2052]: EAP: getDecision: -> PASSTHROUGH
Mar 23 02:49:00 access-node hostapd[2052]: EAP: EAP entering state INITIALIZE_PASSTHROUGH
Mar 23 02:49:00 access-node hostapd[2052]: EAP: EAP entering state AAA_REQUEST
Mar 23 02:49:00 access-node hostapd[2052]: EAP: EAP entering state AAA_IDLE
Mar 23 02:49:00 access-node hostapd[2052]: ens38: STA <REDACTED> IEEE 802.1X: STA identity
         'nacuser'
Mar 23 02:49:00 access-node hostapd[2052]: Encapsulating EAP message into a RADIUS packet
Mar 23 02:49:00 access-node hostapd[2052]: ens38: RADIUS Sending RADIUS message to authent
        ication server
Mar 23 02:49:00 access-node hostapd[2052]: RADIUS message: code=1 (Access-Request) identif
        ier=3 length=165
Mar 23 02:49:00 access-node hostapd[2052]:    Attribute 1 (User-Name) length=9
Mar 23 02:49:00 access-node hostapd[2052]:       Value: 'nacuser'
Mar 23 02:49:00 access-node hostapd[2052]:    Attribute 4 (NAS-IP-Address) length=6
Mar 23 02:49:00 access-node hostapd[2052]:       Value: 10.0.23.1
Mar 23 02:49:00 access-node hostapd[2052]:    Attribute 30 (Called-Station-Id) length=20
Mar 23 02:49:00 access-node hostapd[2052]:       Value: '<REDACTED>:'
Mar 23 02:49:00 access-node hostapd[2052]:    Attribute 61 (NAS-Port-Type) length=6
Mar 23 02:49:00 access-node hostapd[2052]:       Value: 19
Mar 23 02:49:00 access-node hostapd[2052]:    Attribute 6 (Service-Type) length=6
Mar 23 02:49:00 access-node hostapd[2052]:       Value: 2
Mar 23 02:49:00 access-node hostapd[2052]:    Attribute 31 (Calling-Station-Id) length=19
Mar 23 02:49:00 access-node hostapd[2052]:       Value: '<REDACTED>'
Mar 23 02:49:00 access-node hostapd[2052]:    Attribute 77 (Connect-Info) length=23
Mar 23 02:49:00 access-node hostapd[2052]:       Value: 'CONNECT 0Mbps 802.11b'
Mar 23 02:49:00 access-node hostapd[2052]:    Attribute 44 (Acct-Session-Id) length=18
Mar 23 02:49:00 access-node hostapd[2052]:       Value: '2268479955CB85A0'
Mar 23 02:49:00 access-node hostapd[2052]:    Attribute 12 (Framed-MTU) length=6
Mar 23 02:49:00 access-node hostapd[2052]:       Value: 1400
Mar 23 02:49:00 access-node hostapd[2052]:    Attribute 79 (EAP-Message) length=14
Mar 23 02:49:00 access-node hostapd[2052]:       Value: 02fe000c016e616375736572
Mar 23 02:49:00 access-node hostapd[2052]:    Attribute 80 (Message-Authenticator) length=
        18
Mar 23 02:49:00 access-node hostapd[2052]:       Value: ede40768f3affc833f5ecb655d5f6878
Mar 23 02:49:00 access-node hostapd[2052]: ens38: RADIUS Next RADIUS client retransmit in 
        3 seconds
Mar 23 02:49:00 access-node hostapd[2052]: ens38: STA <REDACTED> IEEE 802.1X: received EAP
        OL-Start from STA
Mar 23 02:49:00 access-node hostapd[2052]: ens38: RADIUS Received 80 bytes from RADIUS ser
        ver
Mar 23 02:49:00 access-node hostapd[2052]: ens38: RADIUS Received RADIUS message
Mar 23 02:49:00 access-node hostapd[2052]: RADIUS message: code=11 (Access-Challenge) iden
        tifier=3 length=80
Mar 23 02:49:00 access-node hostapd[2052]:    Attribute 80 (Message-Authenticator) length=
        18
Mar 23 02:49:00 access-node hostapd[2052]:       Value: e9f4ffb9acaefad6cd9fe021f6937580
Mar 23 02:49:00 access-node hostapd[2052]:    Attribute 79 (EAP-Message) length=24
Mar 23 02:49:00 access-node hostapd[2052]:       Value: 01ff0016041090ab891e63b1ce73b12833
        17e89cc41b
Mar 23 02:49:00 access-node hostapd[2052]:    Attribute 24 (State) length=18
Mar 23 02:49:00 access-node hostapd[2052]:       Value: 33b1db94334edf7cde63d907d83a901f
Mar 23 02:49:00 access-node hostapd[2052]: ens38: STA <REDACTED> RADIUS: Received RADIUS p
        acket matched with a pending request, round trip time 0.00 sec
Mar 23 02:49:00 access-node hostapd[2052]: RADIUS packet matching with station <REDACTED>
Mar 23 02:49:00 access-node hostapd[2052]: ens38: STA <REDACTED> IEEE 802.1X: decapsulated
         EAP packet (code=1 id=255 len=22) from RADIUS server: EAP-Request-MD5 (4)
Mar 23 02:49:00 access-node hostapd[2052]: EAP: EAP entering state AAA_RESPONSE
Mar 23 02:49:00 access-node hostapd[2052]: EAP: getId: id=255
Mar 23 02:49:00 access-node hostapd[2052]: EAP: EAP entering state SEND_REQUEST2
Mar 23 02:49:00 access-node hostapd[2052]: EAP: EAP entering state IDLE2
Mar 23 02:49:00 access-node hostapd[2052]: EAP: retransmit timeout 3 seconds (from dynamic
         back off; retransCount=0)
Mar 23 02:49:00 access-node hostapd[2052]: IEEE 802.1X: <REDACTED> BE_AUTH entering state 
        REQUEST
Mar 23 02:49:00 access-node hostapd[2052]: ens38: STA <REDACTED> IEEE 802.1X: Sending EAP 
        Packet (identifier 255)
Mar 23 02:49:00 access-node hostapd[2052]: Received EAPOL packet
Mar 23 02:49:00 access-node hostapd[2052]: ens38: Event NEW_STA (22) received
Mar 23 02:49:00 access-node hostapd[2052]: ens38: Event EAPOL_RX (23) received
Mar 23 02:49:00 access-node hostapd[2052]: IEEE 802.1X: 46 bytes from <REDACTED>
Mar 23 02:49:00 access-node hostapd[2052]:    IEEE 802.1X: version=2 type=0 length=22
Mar 23 02:49:00 access-node hostapd[2052]:    ignoring 20 extra octets after IEEE 802.1X p
        acket
Mar 23 02:49:00 access-node hostapd[2052]: EAP: code=2 (response) identifier=255 length=22
Mar 23 02:49:00 access-node hostapd[2052]: ens38: STA <REDACTED> IEEE 802.1X: received EAP
         packet (code=2 id=255 len=22) from STA: EAP Response-MD5 (4)
Mar 23 02:49:00 access-node hostapd[2052]: IEEE 802.1X: <REDACTED> BE_AUTH entering state 
        RESPONSE
Mar 23 02:49:00 access-node hostapd[2052]: EAP: EAP entering state RECEIVED2
Mar 23 02:49:00 access-node hostapd[2052]: EAP: parseEapResp: rxResp=1 rxInitiate=0 respId
        =255 respMethod=4 respVendor=0 respVendorMethod=0
Mar 23 02:49:00 access-node hostapd[2052]: EAP: EAP entering state AAA_REQUEST
Mar 23 02:49:00 access-node hostapd[2052]: EAP: EAP entering state AAA_IDLE
Mar 23 02:49:00 access-node hostapd[2052]: Encapsulating EAP message into a RADIUS packet
Mar 23 02:49:00 access-node hostapd[2052]: Copied RADIUS State Attribute
Mar 23 02:49:00 access-node hostapd[2052]: ens38: RADIUS Sending RADIUS message to authent
        ication server
Mar 23 02:49:00 access-node hostapd[2052]: RADIUS message: code=1 (Access-Request) identif
        ier=4 length=193
Mar 23 02:49:00 access-node hostapd[2052]:    Attribute 1 (User-Name) length=9
Mar 23 02:49:00 access-node hostapd[2052]:       Value: 'nacuser'
Mar 23 02:49:00 access-node hostapd[2052]:    Attribute 4 (NAS-IP-Address) length=6
Mar 23 02:49:00 access-node hostapd[2052]:       Value: 10.0.23.1
Mar 23 02:49:00 access-node hostapd[2052]: ens38: STA <REDACTED> IEEE 802.1X: unauthorizin
        g port
Mar 23 02:49:00 access-node hostapd[2052]:    Attribute 30 (Called-Station-Id) length=20
Mar 23 02:49:00 access-node hostapd[2052]:       Value: '<REDACTED>:'
Mar 23 02:49:00 access-node hostapd[2052]:    Attribute 61 (NAS-Port-Type) length=6
Mar 23 02:49:00 access-node hostapd[2052]:       Value: 19
Mar 23 02:49:00 access-node hostapd[2052]:    Attribute 6 (Service-Type) length=6
Mar 23 02:49:00 access-node hostapd[2052]:       Value: 2
Mar 23 02:49:00 access-node hostapd[2052]:    Attribute 31 (Calling-Station-Id) length=19
Mar 23 02:49:00 access-node hostapd[2052]:       Value: '<REDACTED>'
Mar 23 02:49:00 access-node hostapd[2052]:    Attribute 77 (Connect-Info) length=23
Mar 23 02:49:00 access-node hostapd[2052]:       Value: 'CONNECT 0Mbps 802.11b'
Mar 23 02:49:00 access-node hostapd[2052]:    Attribute 44 (Acct-Session-Id) length=18
Mar 23 02:49:00 access-node hostapd[2052]:       Value: '2268479955CB85A0'
Mar 23 02:49:00 access-node hostapd[2052]:    Attribute 12 (Framed-MTU) length=6
Mar 23 02:49:00 access-node hostapd[2052]:       Value: 1400
Mar 23 02:49:00 access-node hostapd[2052]:    Attribute 79 (EAP-Message) length=24
Mar 23 02:49:00 access-node hostapd[2052]:       Value: 02ff001604106159cae5bd45916254d85a
        8e12d860af
Mar 23 02:49:00 access-node hostapd[2052]:    Attribute 24 (State) length=18
Mar 23 02:49:00 access-node hostapd[2052]:       Value: 33b1db94334edf7cde63d907d83a901f
Mar 23 02:49:00 access-node hostapd[2052]:    Attribute 80 (Message-Authenticator) length=
        18
Mar 23 02:49:00 access-node hostapd[2052]:       Value: b10e50120d8d1718ff54b589c41af645
Mar 23 02:49:00 access-node hostapd[2052]: ens38: RADIUS Next RADIUS client retransmit in 
        3 seconds
Mar 23 02:49:00 access-node hostapd[2052]: ens38: RADIUS Received 53 bytes from RADIUS ser
        ver
Mar 23 02:49:00 access-node hostapd[2052]: ens38: RADIUS Received RADIUS message
Mar 23 02:49:00 access-node hostapd[2052]: RADIUS message: code=2 (Access-Accept) identifi
        er=4 length=53
Mar 23 02:49:00 access-node hostapd[2052]:    Attribute 80 (Message-Authenticator) length=
        18
Mar 23 02:49:00 access-node hostapd[2052]:       Value: e995ad8e4287ff371bdeaf07853dceac
Mar 23 02:49:00 access-node hostapd[2052]:    Attribute 79 (EAP-Message) length=6
Mar 23 02:49:00 access-node hostapd[2052]:       Value: 03ff0004
Mar 23 02:49:00 access-node hostapd[2052]:    Attribute 1 (User-Name) length=9
Mar 23 02:49:00 access-node hostapd[2052]:       Value: 'nacuser'
Mar 23 02:49:00 access-node hostapd[2052]: ens38: STA <REDACTED> RADIUS: Received RADIUS p
        acket matched with a pending request, round trip time 0.00 sec
Mar 23 02:49:00 access-node hostapd[2052]: RADIUS packet matching with station <REDACTED>
Mar 23 02:49:00 access-node hostapd[2052]: MS-MPPE: 1x_get_keys, could not get keys: 0x5cc
        575f92a40  send: (nil)  recv: (nil)
Mar 23 02:49:00 access-node hostapd[2052]: ens38: STA <REDACTED> IEEE 802.1X: old identity
         'nacuser' updated with User-Name from Access-Accept 'nacuser'
Mar 23 02:49:00 access-node hostapd[2052]: ens38: STA <REDACTED> IEEE 802.1X: decapsulated
         EAP packet (code=3 id=255 len=4) from RADIUS server: EAP Success
Mar 23 02:49:00 access-node hostapd[2052]: EAP: EAP entering state SUCCESS2
Mar 23 02:49:00 access-node hostapd[2052]: ens38: CTRL-EVENT-EAP-SUCCESS2 <REDACTED>
Mar 23 02:49:00 access-node hostapd[2052]: IEEE 802.1X: <REDACTED> BE_AUTH entering state 
        SUCCESS
Mar 23 02:49:00 access-node hostapd[2052]: ens38: STA <REDACTED> IEEE 802.1X: Sending EAP 
        Packet (identifier 255)
Mar 23 02:49:00 access-node hostapd[2052]: IEEE 802.1X: <REDACTED> AUTH_PAE entering state
         AUTHENTICATED
Mar 23 02:49:00 access-node hostapd[2052]: ens38: AP-STA-CONNECTED <REDACTED>
Mar 23 02:49:00 access-node hostapd[2052]: ens38: STA <REDACTED> IEEE 802.1X: Sending EAP 
        Packet (identifier 254)
Mar 23 02:49:00 access-node hostapd[2052]: ens38: STA <REDACTED> IEEE 802.1X: authorizing 
        port
Mar 23 02:49:00 access-node hostapd[2052]: ens38: STA <REDACTED> RADIUS: starting accounti
        ng session 2268479955CB85A0
Mar 23 02:49:00 access-node hostapd[2052]: ens38: STA <REDACTED> IEEE 802.1X: authenticate
        d - EAP type: 4 (MD5)
Mar 23 02:49:00 access-node hostapd[2052]: IEEE 802.1X: <REDACTED> BE_AUTH entering state 
        IDLE
Mar 23 02:49:00 access-node hostapd[2052]: ens38: STA <REDACTED> IEEE 802.1X: received EAP
         packet (code=2 id=254 len=12) from STA: EAP Response-Identity (1)
Mar 23 02:49:00 access-node hostapd[2052]: ens38: STA <REDACTED> IEEE 802.1X: STA identity
         'nacuser'
Mar 23 02:49:00 access-node hostapd[2052]: ens38: RADIUS Sending RADIUS message to authent
        ication server
Mar 23 02:49:00 access-node hostapd[2052]: ens38: RADIUS Next RADIUS client retransmit in 
        3 seconds
Mar 23 02:49:00 access-node hostapd[2052]: ens38: RADIUS Received 80 bytes from RADIUS ser
        ver
Mar 23 02:49:00 access-node hostapd[2052]: ens38: RADIUS Received RADIUS message
Mar 23 02:49:00 access-node hostapd[2052]: ens38: STA <REDACTED> RADIUS: Received RADIUS p
        acket matched with a pending request, round trip time 0.00 sec
Mar 23 02:49:00 access-node hostapd[2052]: ens38: STA <REDACTED> IEEE 802.1X: decapsulated
         EAP packet (code=1 id=255 len=22) from RADIUS server: EAP-Request-MD5 (4)
Mar 23 02:49:00 access-node hostapd[2052]: ens38: STA <REDACTED> IEEE 802.1X: Sending EAP 
        Packet (identifier 255)
Mar 23 02:49:00 access-node hostapd[2052]: ens38: STA <REDACTED> IEEE 802.1X: received EAP
         packet (code=2 id=255 len=22) from STA: EAP Response-MD5 (4)
Mar 23 02:49:00 access-node hostapd[2052]: ens38: RADIUS Sending RADIUS message to authent
        ication server
Mar 23 02:49:00 access-node hostapd[2052]: ens38: RADIUS Next RADIUS client retransmit in 
        3 seconds
Mar 23 02:49:00 access-node hostapd[2052]: ens38: RADIUS Received 53 bytes from RADIUS ser
        ver
Mar 23 02:49:00 access-node hostapd[2052]: ens38: RADIUS Received RADIUS message
Mar 23 02:49:00 access-node hostapd[2052]: ens38: STA <REDACTED> RADIUS: Received RADIUS p
        acket matched with a pending request, round trip time 0.00 sec
Mar 23 02:49:00 access-node hostapd[2052]: ens38: STA <REDACTED> IEEE 802.1X: old identity
         'nacuser' updated with User-Name from Access-Accept 'nacuser'
Mar 23 02:49:00 access-node hostapd[2052]: ens38: STA <REDACTED> IEEE 802.1X: decapsulated
         EAP packet (code=3 id=255 len=4) from RADIUS server: EAP Success
Mar 23 02:49:00 access-node hostapd[2052]: ens38: STA <REDACTED> IEEE 802.1X: Sending EAP 
        Packet (identifier 255)
Mar 23 02:49:00 access-node hostapd[2052]: ens38: STA <REDACTED> IEEE 802.1X: authorizing 
        port
Mar 23 02:49:00 access-node hostapd[2052]: ens38: STA <REDACTED> RADIUS: starting accounti
        ng session 2268479955CB85A0
Mar 23 02:49:00 access-node hostapd[2052]: ens38: STA <REDACTED> IEEE 802.1X: authenticate
        d - EAP type: 4 (MD5)
Mar 23 02:49:03 access-node hostapd[2052]: IEEE 802.1X: <REDACTED> - (EAP) retransWhile --
        > 0
Mar 23 02:49:30 access-node hostapd[2052]: IEEE 802.1X: <REDACTED> - aWhile --> 0
\end{verbatim}

% ---------------------------------------------------------------------
% SUBSECTION B.1.6: NAC INCIDENT PARSER SUCCESS OUTPUT
% ---------------------------------------------------------------------
% This subsection preserves the parser output file demonstrating that
% the collected incident evidence was successfully summarized into a
% deterministic incident conclusion.
% ---------------------------------------------------------------------

\subsubsection{B.1.6 NAC Incident Parser Success Output}
\label{sec:appendix-b-nac-incident-parser-success-output}

\begin{verbatim}
radius-splunk: ~/incident-evidence/nac_incident_parser_success.txt
=== NAC Incident Evidence Summary ===

[1] Failed authentication simulation
- Client remained unauthorized: YES
- Client showed successful connection during failed test: NO
- hostapd saw authentication attempt: YES
- FreeRADIUS recorded failed auth / reject path: YES

[2] Remediation / successful re-test
- Client reached CTRL-EVENT-CONNECTED: YES
- Client reached Authorized state: YES
- hostapd showed successful authorization indicators: YES

[3] Incident conclusion
- SIMULATED NAC INCIDENT SUCCESSFULLY REPRODUCED AND REMEDIATED
\end{verbatim}

% ---------------------------------------------------------------------
% SUBSECTION B.2: GLPI INCIDENT RECORD
% ---------------------------------------------------------------------
% This subsection preserves the exact GLPI incident title, status, raw
% incident text, and the figure-mapping note for the GUI screenshot
% referenced in the body of the report.
% ---------------------------------------------------------------------

\subsection{B.2 GLPI Incident Record}
\label{sec:appendix-b-glpi-incident-record}

% ---------------------------------------------------------------------
% SUBSECTION B.2.1: GLPI INCIDENT TITLE
% ---------------------------------------------------------------------
% This subsection preserves the exact title of the GLPI incident ticket
% created to document the wired 802.1X failure and remediation event.
% ---------------------------------------------------------------------

\subsubsection{B.2.1 GLPI Incident Title}
\label{sec:appendix-b-glpi-incident-title}

\begin{verbatim}
Wired 802.1X authentication failure and remediation for client-1
\end{verbatim}

% ---------------------------------------------------------------------
% SUBSECTION B.2.2: GLPI INCIDENT STATUS
% ---------------------------------------------------------------------
% This subsection preserves the final status of the GLPI incident
% ticket after remediation and validation were completed.
% ---------------------------------------------------------------------

\subsubsection{B.2.2 GLPI Incident Status}
\label{sec:appendix-b-glpi-incident-status}

\begin{verbatim}
Solved
\end{verbatim}

% ---------------------------------------------------------------------
% SUBSECTION B.2.3: GLPI INCIDENT RECORD RAW TEXT
% ---------------------------------------------------------------------
% This subsection preserves the full incident text entered into GLPI so
% that the GUI record referenced in the body remains textually
% reproducible in the appendix.
% ---------------------------------------------------------------------

\subsubsection{B.2.3 GLPI Incident Record Raw Text}
\label{sec:appendix-b-glpi-incident-record-raw-text}

\begin{verbatim}
Incident summary:
client-1 failed wired 802.1X authentication during NAC validation after
the supplicant password was intentionally changed to an incorrect value
to simulate a credential issue.

Observed behavior:
- client-1 remained unauthorized during the failed test
- FreeRADIUS showed a failed authentication / reject path
- hostapd on access-node recorded the 802.1X authentication attempt

Investigation:
- Verified FreeRADIUS service health on radius-splunk
- Confirmed the known-good credentials still produced Access-Accept
  using radtest
- Revalidated the client wpa_supplicant configuration

Remediation:
- Restored the known-good wpa_supplicant credentials on client-1
- Re-ran the wired 802.1X authentication test

Recovery validation:
- client-1 showed CTRL-EVENT-CONNECTED
- client-1 supplicant port reached Authorized state
- access-node hostapd logs showed AP-STA-CONNECTED and port
  authorization

Impact:
This simulated a realistic enterprise NAC incident involving endpoint
credential failure and demonstrated successful troubleshooting and
service restoration.
\end{verbatim}

% ---------------------------------------------------------------------
% SUBSECTION B.3: GUEST NETWORK VALIDATION OUTPUTS
% ---------------------------------------------------------------------
% This subsection preserves the selected raw outputs used to validate
% DHCP onboarding, guest isolation, and allowed internet access on the
% guest segment.
% ---------------------------------------------------------------------

\subsection{B.3 Guest Network Validation Outputs}
\label{sec:appendix-b-guest-network-validation-outputs}

% ---------------------------------------------------------------------
% SUBSECTION B.3.1: GUEST INTERFACE AND ROUTE VALIDATION
% ---------------------------------------------------------------------
% This subsection preserves the raw interface and routing state of the
% guest validation endpoint after successful DHCP onboarding.
% ---------------------------------------------------------------------

\subsubsection{B.3.1 Guest Interface and Route Validation}
\label{sec:appendix-b-guest-interface-and-route-validation}

\begin{verbatim}
guest-test-1: ip -br a
lo               UNKNOWN        127.0.0.1/8 ::1/128
ens33            UP             10.0.20.137/24 fe80::250:56ff:feab:cdef/64

guest-test-1: ip route
default via 10.0.20.254 dev ens33 proto dhcp src 10.0.20.137 metric 100
10.0.20.0/24 dev ens33 proto kernel scope link src 10.0.20.137
10.0.20.254 dev ens33 proto dhcp scope link src 10.0.20.137 metric 100
\end{verbatim}

% ---------------------------------------------------------------------
% SUBSECTION B.3.2: GUEST DHCP RENEWAL VALIDATION
% ---------------------------------------------------------------------
% This subsection preserves the authoritative DHCP renewal output on
% the guest validation endpoint.
% ---------------------------------------------------------------------

\subsubsection{B.3.2 Guest DHCP Renewal Validation}
\label{sec:appendix-b-guest-dhcp-renewal-validation}

\begin{verbatim}
guest-test-1: sudo dhclient -v ens33
Internet Systems Consortium DHCP Client 4.4.1
Copyright 2004-2018 Internet Systems Consortium.
All rights reserved.
For info, please visit https://www.isc.org/software/dhcp/

Listening on LPF/ens33/<REDACTED>
Sending on   LPF/ens33/<REDACTED>
Sending on   Socket/fallback
DHCPDISCOVER on ens33 to 255.255.255.255 port 67 interval 3 (xid=0x4a3f2c1d)
DHCPDISCOVER on ens33 to 255.255.255.255 port 67 interval 6 (xid=0x4a3f2c1d)
DHCPOFFER of 10.0.20.137 from 10.0.20.254
DHCPREQUEST for 10.0.20.137 on ens33 to 255.255.255.255 port 67 (xid=0x4a3f2c1d)
DHCPACK of 10.0.20.137 from 10.0.20.254 (xid=0x4a3f2c1d)
bound to 10.0.20.137 -- renewal in 3600 seconds.
\end{verbatim}

% ---------------------------------------------------------------------
% SUBSECTION B.3.3: GUEST GATEWAY REACHABILITY VALIDATION
% ---------------------------------------------------------------------
% This subsection preserves the raw ping output confirming that the
% guest endpoint can reach the pfSense guest gateway.
% ---------------------------------------------------------------------

\subsubsection{B.3.3 Guest Gateway Reachability Validation}
\label{sec:appendix-b-guest-gateway-reachability-validation}

\begin{verbatim}
guest-test-1: ping -c 4 10.0.20.254
PING 10.0.20.254 (10.0.20.254) 56(84) bytes of data.
64 bytes from 10.0.20.254: icmp_seq=1 ttl=64 time=1.13 ms
64 bytes from 10.0.20.254: icmp_seq=2 ttl=64 time=0.710 ms
64 bytes from 10.0.20.254: icmp_seq=3 ttl=64 time=0.728 ms
64 bytes from 10.0.20.254: icmp_seq=4 ttl=64 time=3.25 ms

--- 10.0.20.254 ping statistics ---
4 packets transmitted, 4 received, 0% packet loss, time 3038ms
rtt min/avg/max/mdev = 0.710/1.454/3.247/1.048 ms
\end{verbatim}

% ---------------------------------------------------------------------
% SUBSECTION B.3.4: GUEST PROTECTED NETWORK ISOLATION VALIDATION
% ---------------------------------------------------------------------
% This subsection preserves the failed reachability tests proving that
% the guest endpoint was isolated from protected internal networks.
% ---------------------------------------------------------------------

\subsubsection{B.3.4 Guest Protected Network Isolation Validation}
\label{sec:appendix-b-guest-protected-network-isolation-validation}

\begin{verbatim}
guest-test-1: ping -c 4 10.0.30.10
PING 10.0.30.10 (10.0.30.10) 56(84) bytes of data.

--- 10.0.30.10 ping statistics ---
4 packets transmitted, 0 received, 100% packet loss, time 3074ms

guest-test-1: ping -c 4 10.0.30.20
PING 10.0.30.20 (10.0.30.20) 56(84) bytes of data.

--- 10.0.30.20 ping statistics ---
4 packets transmitted, 0 received, 100% packet loss, time 3071ms

guest-test-1: ping -c 4 10.0.34.1
PING 10.0.34.1 (10.0.34.1) 56(84) bytes of data.

--- 10.0.34.1 ping statistics ---
4 packets transmitted, 0 received, 100% packet loss, time 3080ms

guest-test-1: ping -c 4 10.0.23.1
PING 10.0.23.1 (10.0.23.1) 56(84) bytes of data.

--- 10.0.23.1 ping statistics ---
4 packets transmitted, 0 received, 100% packet loss, time 3079ms

guest-test-1: ping -c 4 10.0.12.1
PING 10.0.12.1 (10.0.12.1) 56(84) bytes of data.

--- 10.0.12.1 ping statistics ---
4 packets transmitted, 0 received, 100% packet loss, time 3050ms
\end{verbatim}

% ---------------------------------------------------------------------
% SUBSECTION B.3.5: GUEST HTTP-TO-GLPI ISOLATION VALIDATION
% ---------------------------------------------------------------------
% This subsection preserves the HTTP timeout result showing that the
% guest endpoint could not reach the protected GLPI service.
% ---------------------------------------------------------------------

\subsubsection{B.3.5 Guest HTTP-to-GLPI Isolation Validation}
\label{sec:appendix-b-guest-http-to-glpi-isolation-validation}

\begin{verbatim}
guest-test-1: curl -I --max-time 5 http://10.0.30.20/
curl: (28) Connection timed out after 5001 milliseconds
\end{verbatim}

% ---------------------------------------------------------------------
% SUBSECTION B.3.6: GUEST INTERNET-ALLOWED VALIDATION
% ---------------------------------------------------------------------
% This subsection preserves the raw outputs showing that the guest
% endpoint retained permitted internet-bound access.
% ---------------------------------------------------------------------

\subsubsection{B.3.6 Guest Internet-Allowed Validation}
\label{sec:appendix-b-guest-internet-allowed-validation}

\begin{verbatim}
guest-test-1: ping -c 4 1.1.1.1
PING 1.1.1.1 (1.1.1.1) 56(84) bytes of data.
64 bytes from 1.1.1.1: icmp_seq=1 ttl=58 time=8.04 ms
64 bytes from 1.1.1.1: icmp_seq=2 ttl=58 time=10.6 ms
64 bytes from 1.1.1.1: icmp_seq=3 ttl=58 time=8.98 ms
64 bytes from 1.1.1.1: icmp_seq=4 ttl=58 time=11.4 ms

--- 1.1.1.1 ping statistics ---
4 packets transmitted, 4 received, 0% packet loss, time 3004ms
rtt min/avg/max/mdev = 8.043/9.758/11.444/1.326 ms

guest-test-1: curl -I --max-time 5 http://www.google.com
HTTP/1.1 200 OK
Content-Type: text/html; charset=ISO-8859-1
Content-Security-Policy-Report-Only: object-src 'none';base-uri 'self';script-src 'nonce-Q
        UPlZ9-ttgVc2d-VTB3gXA' 'strict-dynamic' 'report-sample' 'unsafe-eval' 'unsafe-inline' 
        https: http:;report-uri https://csp.withgoogle.com/csp/gws/other-hp
Reporting-Endpoints: default="//www.google.com/httpservice/retry/jserror?ei=n1PCacmWHbav5N
        oPtbr14A0&cad=crash&error=Page%20Crash&jsel=1&bver=2405&dpf=
        nIfF8NgQdNWxHThvch26jsRSkc4DpZUm6gVT80MTHMY"
P3P: CP="This is not a P3P policy! See g.co/p3phelp for more info."
Date: Tue, 24 Mar 2026 09:04:31 GMT
Server: gws
X-XSS-Protection: 0
X-Frame-Options: SAMEORIGIN
Expires: Tue, 24 Mar 2026 09:04:31 GMT
Cache-Control: private
Set-Cookie: <REDACTED>
Set-Cookie: <REDACTED>
Transfer-Encoding: chunked
\end{verbatim}

% ---------------------------------------------------------------------
% SUBSECTION B.4: AUTOMATION EXECUTION EVIDENCE
% ---------------------------------------------------------------------
% This subsection preserves the raw Ansible run outputs and post-run
% artifact preview evidence used in the automation section.
% ---------------------------------------------------------------------

\subsection{B.4 Automation Execution Evidence}
\label{sec:appendix-b-automation-execution-evidence}

% ---------------------------------------------------------------------
% SUBSECTION B.4.1: ANSIBLE VALIDATION-ONLY RUN OUTPUT
% ---------------------------------------------------------------------
% This subsection preserves the authoritative validation-only Ansible
% run output.
% ---------------------------------------------------------------------

\subsubsection{B.4.1 Ansible Validation-Only Run Output}
\label{sec:appendix-b-ansible-validation-only-run-output}

\begin{verbatim}
radius-splunk: ansible-playbook -i inventory/hosts.ini playbooks/nac_validation.yml -K

PLAY [Validate NAC services and optionally remediate] **************************

TASK [Collect basic hostname] **************************************************
ok: [radius-splunk]
ok: [access-node]

TASK [Check hostapd-wired service on access-node] ******************************
skipping: [radius-splunk]
ok: [access-node]

TASK [Check freeradius service on radius-splunk] *******************************
skipping: [access-node]
ok: [radius-splunk]

TASK [Restart hostapd-wired service when remediation is enabled] ***************
skipping: [radius-splunk]
skipping: [access-node]

TASK [Restart freeradius service when remediation is enabled] ******************
skipping: [access-node]
skipping: [radius-splunk]

TASK [Re-check hostapd-wired service on access-node] ***************************
skipping: [radius-splunk]
ok: [access-node]

TASK [Re-check freeradius service on radius-splunk] ****************************
skipping: [access-node]
ok: [radius-splunk]

TASK [Collect recent hostapd logs] *********************************************
skipping: [radius-splunk]
ok: [access-node]

TASK [Collect recent freeradius service logs] **********************************
skipping: [access-node]
ok: [radius-splunk]

TASK [Collect recent freeradius application log] *******************************
skipping: [access-node]
ok: [radius-splunk]

TASK [Save hostapd pre-check status locally] ***********************************
skipping: [radius-splunk]
changed: [access-node -> localhost]

TASK [Save hostapd post-check status locally] **********************************
skipping: [radius-splunk]
changed: [access-node -> localhost]

TASK [Save hostapd restart result locally] *************************************
skipping: [radius-splunk]
changed: [access-node -> localhost]

TASK [Save hostapd logs locally] ***********************************************
skipping: [radius-splunk]
ok: [access-node -> localhost]

TASK [Save freeradius pre-check status locally] ********************************
skipping: [access-node]
changed: [radius-splunk -> localhost]

TASK [Save freeradius post-check status locally] *******************************
skipping: [access-node]
changed: [radius-splunk -> localhost]

TASK [Save freeradius restart result locally] **********************************
skipping: [access-node]
changed: [radius-splunk -> localhost]

TASK [Save freeradius service logs locally] ************************************
skipping: [access-node]
ok: [radius-splunk -> localhost]

TASK [Save freeradius application log locally] *********************************
skipping: [access-node]
ok: [radius-splunk -> localhost]

PLAY RECAP *********************************************************************
access-node                : ok=8    changed=3    unreachable=0    failed=0    skipped=11   
        rescued=0    ignored=0
radius-splunk              : ok=10   changed=3    unreachable=0    failed=0    skipped=9    
        rescued=0    ignored=0
\end{verbatim}

% ---------------------------------------------------------------------
% SUBSECTION B.4.2: ANSIBLE REMEDIATION-ENABLED RUN OUTPUT
% ---------------------------------------------------------------------
% This subsection preserves the authoritative remediation-enabled
% Ansible run output.
% ---------------------------------------------------------------------

\subsubsection{B.4.2 Ansible Remediation-Enabled Run Output}
\label{sec:appendix-b-ansible-remediation-enabled-run-output}

\begin{verbatim}
radius-splunk: ansible-playbook -i inventory/hosts.ini playbooks/nac_validation.yml -K -e remediate=true

PLAY [Validate NAC services and optionally remediate] **************************

TASK [Collect basic hostname] **************************************************
ok: [radius-splunk]
ok: [access-node]

TASK [Check hostapd-wired service on access-node] ******************************
skipping: [radius-splunk]
ok: [access-node]

TASK [Check freeradius service on radius-splunk] *******************************
skipping: [access-node]
ok: [radius-splunk]

TASK [Restart hostapd-wired service when remediation is enabled] ***************
skipping: [radius-splunk]
changed: [access-node]

TASK [Restart freeradius service when remediation is enabled] ******************
skipping: [access-node]
changed: [radius-splunk]

TASK [Re-check hostapd-wired service on access-node] ***************************
skipping: [radius-splunk]
ok: [access-node]

TASK [Re-check freeradius service on radius-splunk] ****************************
skipping: [access-node]
ok: [radius-splunk]

TASK [Collect recent hostapd logs] *********************************************
skipping: [radius-splunk]
ok: [access-node]

TASK [Collect recent freeradius service logs] **********************************
skipping: [access-node]
ok: [radius-splunk]

TASK [Collect recent freeradius application log] *******************************
skipping: [access-node]
ok: [radius-splunk]

TASK [Save hostapd pre-check status locally] ***********************************
skipping: [radius-splunk]
changed: [access-node -> localhost]

TASK [Save hostapd post-check status locally] **********************************
skipping: [radius-splunk]
changed: [access-node -> localhost]

TASK [Save hostapd restart result locally] *************************************
skipping: [radius-splunk]
ok: [access-node -> localhost]

TASK [Save hostapd logs locally] ***********************************************
skipping: [radius-splunk]
changed: [access-node -> localhost]

TASK [Save freeradius pre-check status locally] ********************************
skipping: [access-node]
changed: [radius-splunk -> localhost]

TASK [Save freeradius post-check status locally] *******************************
skipping: [access-node]
changed: [radius-splunk -> localhost]

TASK [Save freeradius restart result locally] **********************************
skipping: [access-node]
ok: [radius-splunk -> localhost]

TASK [Save freeradius service logs locally] ************************************
skipping: [access-node]
changed: [radius-splunk -> localhost]

TASK [Save freeradius application log locally] *********************************
skipping: [access-node]
changed: [radius-splunk -> localhost]

PLAY RECAP *********************************************************************
access-node                : ok=9    changed=4    unreachable=0    failed=0    skipped=10   
        rescued=0    ignored=0
radius-splunk              : ok=11   changed=5    unreachable=0    failed=0    skipped=8    
        rescued=0    ignored=0
\end{verbatim}

% ---------------------------------------------------------------------
% SUBSECTION B.4.3: ANSIBLE POST-RUN ARTIFACT DIRECTORY LISTING
% ---------------------------------------------------------------------
% This subsection preserves the authoritative post-run Ansible artifact
% directory listing.
% ---------------------------------------------------------------------

\subsubsection{B.4.3 Ansible Post-Run Artifact Directory Listing}
\label{sec:appendix-b-ansible-post-run-artifact-directory-listing}

\begin{verbatim}
radius-splunk: ls -lh ~/ansible-nac/artifacts
total 44K
-rw-r--r-- 1 nac-lab nac-lab 3.5K Mar 24 12:23 access-node_hostapd_logs.txt
-rw-r--r-- 1 nac-lab nac-lab   25 Mar 24 12:43 access-node_hostapd_restart_result.txt
-rw-r--r-- 1 nac-lab nac-lab 1.4K Mar 24 12:43 access-node_hostapd_status_after.txt
-rw-r--r-- 1 nac-lab nac-lab 1.4K Mar 24 12:43 access-node_hostapd_status_before.txt
-rw-r--r-- 1 root    root    1.4K Mar 23 19:30 access-node_hostapd_status.txt
-rw-r--r-- 1 nac-lab nac-lab 4.0K Mar 24 12:23 radius-splunk_freeradius_journal.txt
-rw-r--r-- 1 nac-lab nac-lab 3.1K Mar 24 12:23 radius-splunk_freeradius_radiuslog.txt
-rw-r--r-- 1 nac-lab nac-lab   25 Mar 24 12:43 radius-splunk_freeradius_restart_result.txt
-rw-r--r-- 1 nac-lab nac-lab 1.8K Mar 24 12:43 radius-splunk_freeradius_status_after.txt
-rw-r--r-- 1 nac-lab nac-lab 1.8K Mar 24 12:43 radius-splunk_freeradius_status_before.txt
-rw-r--r-- 1 root    root    1.8K Mar 23 19:30 radius-splunk_freeradius_status.txt
\end{verbatim}

% ---------------------------------------------------------------------
% SUBSECTION B.4.4: ANSIBLE GENERATED ARTIFACT PREVIEW
% ---------------------------------------------------------------------
% This subsection preserves the first-lines preview of the generated
% Ansible artifact files.
% ---------------------------------------------------------------------

\subsubsection{B.4.4 Ansible Generated Artifact Preview}
\label{sec:appendix-b-ansible-generated-artifact-preview}

\begin{verbatim}
===== /home/nac-lab/ansible-nac/artifacts/access-node_hostapd_logs.txt =====
Mar 24 12:23:12 access-node hostapd[1927]: hostapd_interface_deinit_free(0x61845b677590)
Mar 24 12:23:12 access-node hostapd[1927]: hostapd_interface_deinit_free: num_bss=1 
        conf->num_bss=1
Mar 24 12:23:12 access-node hostapd[1927]: hostapd_interface_deinit(0x61845b677590)
Mar 24 12:23:12 access-node hostapd[1927]: ens38: interface state ENABLED->DISABLED
Mar 24 12:23:12 access-node hostapd[1927]: hostapd_bss_deinit: deinit bss ens38
Mar 24 12:23:12 access-node hostapd[1927]: ens38: Flushing old station entries
Mar 24 12:23:12 access-node hostapd[1927]: ens38: Deauthenticate all stations
Mar 24 12:23:12 access-node hostapd[1927]: ens38: AP-DISABLED

===== /home/nac-lab/ansible-nac/artifacts/access-node_hostapd_restart_result.txt =====
remediation not requested

===== /home/nac-lab/ansible-nac/artifacts/access-node_hostapd_status_after.txt =====
- hostapd-wired.service - hostapd wired 802.1X authenticator (lab)
     Loaded: loaded (/etc/systemd/system/hostapd-wired.service; enabled; preset: enabled)
     Active: active (running) since Tue 2026-03-24 12:23:12 UTC; 20min ago
   Main PID: 2054 (hostapd)
      Tasks: 1 (limit: 2206)
     Memory: 844.0K (peak: 1.0M)
        CPU: 31ms
     CGroup: /system.slice/hostapd-wired.service

===== /home/nac-lab/ansible-nac/artifacts/access-node_hostapd_status_before.txt =====
- hostapd-wired.service - hostapd wired 802.1X authenticator (lab)
     Loaded: loaded (/etc/systemd/system/hostapd-wired.service; enabled; preset: enabled)
     Active: active (running) since Tue 2026-03-24 12:23:12 UTC; 20min ago
   Main PID: 2054 (hostapd)
      Tasks: 1 (limit: 2206)
     Memory: 844.0K (peak: 1.0M)
        CPU: 31ms
     CGroup: /system.slice/hostapd-wired.service

===== /home/nac-lab/ansible-nac/artifacts/access-node_hostapd_status.txt =====
- hostapd-wired.service - hostapd wired 802.1X authenticator (lab)
     Loaded: loaded (/etc/systemd/system/hostapd-wired.service; enabled; preset: enabled)
     Active: active (running) since Mon 2026-03-23 18:51:07 UTC; 39min ago
   Main PID: 1017 (hostapd)
      Tasks: 1 (limit: 2206)
     Memory: 2.9M (peak: 3.1M)
        CPU: 51ms
     CGroup: /system.slice/hostapd-wired.service

===== /home/nac-lab/ansible-nac/artifacts/radius-splunk_freeradius_journal.txt =====
Mar 24 12:22:39 radius-splunk freeradius[3253]: Compiling Post-Auth-Type Challenge for 
        attr Post-Auth-Type
Mar 24 12:22:39 radius-splunk freeradius[3253]: Compiling Post-Auth-Type Client-Lost for 
        attr Post-Auth-Type
Mar 24 12:22:39 radius-splunk freeradius[3253]: radiusd: #### Skipping IP addresses and 
        Ports ####
Mar 24 12:22:39 radius-splunk freeradius[3253]: Configuration appears to be OK
Mar 24 12:22:39 radius-splunk systemd[1]: Started freeradius.service - FreeRADIUS 
        multi-protocol policy server.
Mar 24 12:23:12 radius-splunk systemd[1]: Stopping freeradius.service - FreeRADIUS 
        multi-protocol policy server...
Mar 24 12:23:12 radius-splunk systemd[1]: freeradius.service: Deactivated successfully.
Mar 24 12:23:12 radius-splunk systemd[1]: Stopped freeradius.service - FreeRADIUS 
        multi-protocol policy server.

===== /home/nac-lab/ansible-nac/artifacts/radius-splunk_freeradius_radiuslog.txt =====
Tue Mar 24 11:34:18 2026 : Info: Loaded virtual server <default>
Tue Mar 24 11:34:18 2026 : Warning: Ignoring "sql" (see raddb/mods-available/README.rst)
Tue Mar 24 11:34:18 2026 : Warning: Ignoring "ldap" (see raddb/mods-available/README.rst)
Tue Mar 24 11:34:18 2026 : Info:  # Skipping contents of 'if' as it is always 'false' -- 
        /etc/freeradius/3.0/sites-enabled/inner-tunnel:366
Tue Mar 24 11:34:18 2026 : Info: Loaded virtual server inner-tunnel
Tue Mar 24 11:34:18 2026 : Info: Loaded virtual server default
Tue Mar 24 11:34:18 2026 : Info: Ready to process requests
Tue Mar 24 11:57:46 2026 : Info: Signalled to terminate

===== /home/nac-lab/ansible-nac/artifacts/radius-splunk_freeradius_restart_result.txt =====
remediation not requested

===== /home/nac-lab/ansible-nac/artifacts/radius-splunk_freeradius_status_after.txt =====
- freeradius.service - FreeRADIUS multi-protocol policy server
     Loaded: loaded (/usr/lib/systemd/system/freeradius.service; enabled; preset: enabled)
     Active: active (running) since Tue 2026-03-24 12:23:12 UTC; 20min ago
       Docs: man:radiusd(8)
             man:radiusd.conf(5)
             http://wiki.freeradius.org/
             http://networkradius.com/doc/

===== /home/nac-lab/ansible-nac/artifacts/radius-splunk_freeradius_status_before.txt =====
- freeradius.service - FreeRADIUS multi-protocol policy server
     Loaded: loaded (/usr/lib/systemd/system/freeradius.service; enabled; preset: enabled)
     Active: active (running) since Tue 2026-03-24 12:23:12 UTC; 20min ago
       Docs: man:radiusd(8)
             man:radiusd.conf(5)
             http://wiki.freeradius.org/
             http://networkradius.com/doc/

===== /home/nac-lab/ansible-nac/artifacts/radius-splunk_freeradius_status.txt =====
- freeradius.service - FreeRADIUS multi-protocol policy server
     Loaded: loaded (/usr/lib/systemd/system/freeradius.service; enabled; preset: enabled)
     Active: active (running) since Mon 2026-03-23 18:52:42 UTC; 37min ago
       Docs: man:radiusd(8)
             man:radiusd.conf(5)
             http://wiki.freeradius.org/
             http://networkradius.com/doc/
\end{verbatim}

% ---------------------------------------------------------------------
% SUBSECTION B.5: SELECTED OBSERVABILITY EXCERPTS
% ---------------------------------------------------------------------
% This subsection preserves the key raw observability outputs used in
% the logging and monitoring section.
% ---------------------------------------------------------------------

\subsection{B.5 Selected Observability Excerpts}
\label{sec:appendix-b-selected-observability-excerpts}

% ---------------------------------------------------------------------
% SUBSECTION B.5.1: SPLUNK UF RUNTIME STATUS
% ---------------------------------------------------------------------
% This subsection preserves the authoritative Splunk Universal
% Forwarder runtime status output.
% ---------------------------------------------------------------------

\subsubsection{B.5.1 Splunk UF Runtime Status}
\label{sec:appendix-b-splunk-uf-runtime-status}

\begin{verbatim}
radius-splunk: sudo /opt/splunkforwarder/bin/splunk status
Warning: Attempting to revert the SPLUNK_HOME ownership
Warning: Executing "chown -R root /opt/splunkforwarder"
splunkd is running (PID: 1489)
splunk helpers are running (PIDs: 1490)
\end{verbatim}

% ---------------------------------------------------------------------
% SUBSECTION B.5.2: SYSLOG EXCERPT
% ---------------------------------------------------------------------
% This subsection preserves the selected syslog excerpt used in the
% observability section.
% ---------------------------------------------------------------------

\subsubsection{B.5.2 syslog Excerpt}
\label{sec:appendix-b-syslog-excerpt}

\begin{verbatim}
radius-splunk: sudo tail -n 15 /var/log/syslog
2026-03-24T13:14:41.348438+00:00 radius-splunk multipathd[467]: sda: failed to get sgio 
        uid: No such file or directory
2026-03-24T13:14:41.348468+00:00 radius-splunk multipathd[467]: sda: no WWID in state 
        "undef
2026-03-24T13:14:41.348639+00:00 radius-splunk multipathd[467]: ", giving up
2026-03-24T13:14:41.348690+00:00 radius-splunk multipathd[467]: sda: check_path() failed, 
        removing
2026-03-24T13:15:01.629554+00:00 radius-splunk CRON[1579]: (root) CMD (command -v 
        debian-sa1 > /dev/null && debian-sa1 1 1)
2026-03-24T13:17:01.652202+00:00 radius-splunk CRON[1587]: (root) CMD (cd / && run-parts 
        --report /etc/cron.hourly)
2026-03-24T13:19:05.468094+00:00 radius-splunk systemd[1]: Starting update-notifier-
        download.service - Download data for packages that failed at package install 
        time...
2026-03-24T13:19:05.678561+00:00 radius-splunk systemd[1]: update-notifier-download.
        service: Deactivated successfully.
2026-03-24T13:19:05.678792+00:00 radius-splunk systemd[1]: Finished update-notifier-
        download.service - Download data for packages that failed at package install time.
2026-03-24T13:19:31.934534+00:00 radius-splunk systemd[1376]: launchpadlib-cache-clean.
        service - Clean up old files in the Launchpadlib cache was skipped because of an
        unmet condition check (ConditionPathExists=/home/nac-lab/.launchpadlib/
        api.launchpad.net/cache).
\end{verbatim}

% ---------------------------------------------------------------------
% SUBSECTION B.5.3: AUTH.LOG EXCERPT
% ---------------------------------------------------------------------
% This subsection preserves the selected auth.log excerpt used in the
% observability section.
% ---------------------------------------------------------------------

\subsubsection{B.5.3 auth.log Excerpt}
\label{sec:appendix-b-auth-log-excerpt}

\begin{verbatim}
radius-splunk: sudo tail -n 15 /var/log/auth.log
2026-03-24T13:20:04.756710+00:00 radius-splunk sshd[1628]: pam_unix(sshd:session): session
         opened for user nac-lab(uid=1000) by nac-lab(uid=0)
2026-03-24T13:20:04.762714+00:00 radius-splunk systemd-logind[816]: New session 5 of user 
        nac-lab.
2026-03-24T13:20:05.503276+00:00 radius-splunk sshd[1686]: Received disconnect from 192.16
        8.195.1 port 41932:11: disconnected by user
2026-03-24T13:20:05.503525+00:00 radius-splunk sshd[1686]: Disconnected from user nac-lab 
        192.168.195.1 port 41932
2026-03-24T13:20:05.504875+00:00 radius-splunk sshd[1628]: pam_unix(sshd:session): session
         closed for user nac-lab
2026-03-24T13:20:05.511847+00:00 radius-splunk systemd-logind[816]: Session 5 logged out. 
        Waiting for processes to exit.
2026-03-24T13:20:05.514296+00:00 radius-splunk systemd-logind[816]: Removed session 5.
2026-03-24T13:20:37.703552+00:00 radius-splunk sudo:  nac-lab : TTY=tty1 ; PWD=/home/nac-l
        ab ; USER=root ; COMMAND=/usr/bin/tail -n 15 /var/log/syslog
2026-03-24T13:20:37.704672+00:00 radius-splunk sudo: pam_unix(sudo:session): session opene
        d for user root(uid=0) by nac-lab(uid=1000)
2026-03-24T13:20:37.731117+00:00 radius-splunk sudo: pam_unix(sudo:session): session close
        d for user root
2026-03-24T13:21:39.041208+00:00 radius-splunk sudo:  nac-lab : TTY=tty1 ; PWD=/home/nac-l
        ab ; USER=root ; COMMAND=/usr/bin/tail -n 15 /var/log/auth.log
2026-03-24T13:21:39.042808+00:00 radius-splunk sudo: pam_unix(sudo:session): session opene
        d for user root(uid=0) by nac-lab(uid=1000)
2026-03-24T13:21:39.068864+00:00 radius-splunk sudo: pam_unix(sudo:session): session close
        d for user root
2026-03-24T13:22:03.513481+00:00 radius-splunk sudo:  nac-lab : TTY=tty1 ; PWD=/home/nac-l
        ab ; USER=root ; COMMAND=/usr/bin/tail -n 15 /var/log/auth.log
2026-03-24T13:22:03.514636+00:00 radius-splunk sudo: pam_unix(sudo:session): session opene
        d for user root(uid=0) by nac-lab(uid=1000)
\end{verbatim}

% ---------------------------------------------------------------------
% SUBSECTION B.5.4: FREERADIUS RADIUS.LOG EXCERPT
% ---------------------------------------------------------------------
% This subsection preserves the selected FreeRADIUS application log
% excerpt used in the observability section.
% ---------------------------------------------------------------------

\subsubsection{B.5.4 FreeRADIUS radius.log Excerpt}
\label{sec:appendix-b-freeradius-radius-log-excerpt}

\begin{verbatim}
radius-splunk: sudo tail -n 20 /var/log/freeradius/radius.log
Tue Mar 24 12:23:12 2026 : Info: Debug state unknown (cap_sys_ptrace capability not set)
Tue Mar 24 12:23:12 2026 : Info: systemd watchdog interval is 30.00 secs
Tue Mar 24 12:23:12 2026 : Info: Loaded virtual server <default>
Tue Mar 24 12:23:12 2026 : Warning: Ignoring "sql" (see raddb/mods-available/README.rst)
Tue Mar 24 12:23:12 2026 : Warning: Ignoring "ldap" (see raddb/mods-available/README.rst)
Tue Mar 24 12:23:12 2026 : Info:  # Skipping contents of 'if' as it is always 'false' --
        /etc/freeradius/3.0/sites-enabled/inner-tunnel:366
Tue Mar 24 12:23:12 2026 : Info: Loaded virtual server inner-tunnel
Tue Mar 24 12:23:12 2026 : Info: Loaded virtual server default
Tue Mar 24 12:23:12 2026 : Info: Ready to process requests
Tue Mar 24 12:53:42 2026 : Info: Signalled to terminate
Tue Mar 24 12:53:42 2026 : Info: Exiting normally
Tue Mar 24 13:14:05 2026 : Info: Debug state unknown (cap_sys_ptrace capability not set)
Tue Mar 24 13:14:05 2026 : Info: systemd watchdog interval is 30.00 secs
Tue Mar 24 13:14:05 2026 : Info: Loaded virtual server <default>
Tue Mar 24 13:14:05 2026 : Warning: Ignoring "sql" (see raddb/mods-available/README.rst)
Tue Mar 24 13:14:05 2026 : Warning: Ignoring "ldap" (see raddb/mods-available/README.rst)
Tue Mar 24 13:14:05 2026 : Info:  # Skipping contents of 'if' as it is always 'false' -- 
        /etc/freeradius/3.0/sites-enabled/inner-tunnel:366
Tue Mar 24 13:14:05 2026 : Info: Loaded virtual server inner-tunnel
Tue Mar 24 13:14:05 2026 : Info: Loaded virtual server default
Tue Mar 24 13:14:05 2026 : Info: Ready to process requests
\end{verbatim}

% ---------------------------------------------------------------------
% SUBSECTION B.5.5: HOSTAPD SERVICE STATUS AND JOURNAL EXCERPT
% ---------------------------------------------------------------------
% This subsection preserves the selected hostapd service status and
% journal outputs used in the observability section.
% ---------------------------------------------------------------------

\subsubsection{B.5.5 hostapd Service Status and Journal Excerpt}
\label{sec:appendix-b-hostapd-service-status-and-journal-excerpt}

\begin{verbatim}
access-node: sudo systemctl status hostapd-wired.service --no-pager
- hostapd-wired.service - hostapd wired 802.1X authenticator (lab)
     Loaded: loaded (/etc/systemd/system/hostapd-wired.service; enabled; preset: enabled)
     Active: active (running) since Tue 2026-03-24 04:20:41 UTC; 16min ago
   Main PID: 975 (hostapd)
      Tasks: 1 (limit: 2206)
     Memory: 2.9M (peak: 3.1M)
        CPU: 42ms
     CGroup: /system.slice/hostapd-wired.service
             975 /usr/sbin/hostapd -dd /etc/hostapd/lab/hostapd-wired.conf

Mar 24 04:20:41 access-node hostapd[975]: Using interface ens38 with hwaddr <REDACTED> and
         ssid ""
Mar 24 04:20:41 access-node hostapd[975]: ens38: RADIUS Authentication server 10.0.30.10:1
        812
Mar 24 04:20:41 access-node hostapd[975]: RADIUS local address: 10.0.23.1:45833
Mar 24 04:20:41 access-node hostapd[975]: ens38: IEEE 802.11 Fetching hardware channel/rat
        e support not supported.
Mar 24 04:20:41 access-node hostapd[975]: ens38: RADIUS Authentication server 10.0.30.10:1
        812
Mar 24 04:20:41 access-node hostapd[975]: ens38: Deauthenticate all stations at BSS start
Mar 24 04:20:41 access-node hostapd[975]: ens38: interface state UNINITIALIZED->ENABLED
Mar 24 04:20:41 access-node hostapd[975]: ens38: AP-ENABLED
Mar 24 04:20:41 access-node hostapd[975]: ens38: Setup of interface done.

access-node: sudo journalctl -u hostapd-wired.service -n 20 --no-pager
Mar 24 02:13:55 access-node hostapd[2398]: ens38: CTRL-EVENT-TERMINATING
Mar 24 02:13:55 access-node hostapd[2398]: hostapd_free_hapd_data(ens38)
Mar 24 02:13:55 access-node hostapd[2398]: hostapd_interface_deinit_free: driver=0x565a2d7
        47060 drv_priv=0x565a38444cd0 -> hapd_deinit
Mar 24 02:13:55 access-node hostapd[2398]: hostapd_interface_free(0x565a38442590)
Mar 24 02:13:55 access-node hostapd[2398]: hostapd_interface_free: free hapd 0x565a38443d4
        0
Mar 24 02:13:55 access-node hostapd[2398]: hostapd_cleanup_iface(0x565a38442590)
Mar 24 02:13:55 access-node hostapd[2398]: hostapd_cleanup_iface_partial(0x565a38442590)
Mar 24 02:13:55 access-node hostapd[2398]: hostapd_cleanup_iface: free iface=0x565a3844259
        0
Mar 24 02:13:55 access-node systemd[1]: hostapd-wired.service: Deactivated successfully.
Mar 24 02:13:55 access-node systemd[1]: Stopped hostapd-wired.service - hostapd wired 802.
        1X authenticator (lab).
-- Boot 4efbb5bbbad2458a94bc281c688f0c9b --
Mar 24 04:20:41 access-node systemd[1]: Started hostapd-wired.service - hostapd wired 802.
        1X authenticator (lab).
Mar 24 04:20:41 access-node hostapd[975]: random: getrandom() support available
Mar 24 04:20:41 access-node hostapd[975]: Configuration file: /etc/hostapd/lab/hostapd-wir
        ed.conf
Mar 24 04:20:41 access-node hostapd[975]: eapol_version=2
Mar 24 04:20:41 access-node hostapd[975]: ctrl_interface_group=0
Mar 24 04:20:41 access-node hostapd[975]: Opening raw packet socket for ifindex 3
Mar 24 04:20:41 access-node hostapd[975]: BSS count 1, BSSID mask 00:00:00:00:00:00 (0 bit
        s)
Mar 24 04:20:41 access-node hostapd[975]: ens38: IEEE 802.11 Fetching hardware channel/rat
        e support not supported.
Mar 24 04:20:41 access-node hostapd[975]: Completing interface initialization
Mar 24 04:20:41 access-node hostapd[975]: hostapd_setup_bss(hapd=0x648c62cc3d40 (ens38), f
        irst=1)
Mar 24 04:20:41 access-node hostapd[975]: Using interface ens38 with hwaddr <REDACTED> and
         ssid ""
Mar 24 04:20:41 access-node hostapd[975]: ens38: RADIUS Authentication server 10.0.30.10:1
        812
Mar 24 04:20:41 access-node hostapd[975]: RADIUS local address: 10.0.23.1:45833
Mar 24 04:20:41 access-node hostapd[975]: ens38: IEEE 802.11 Fetching hardware channel/rat
        e support not supported.
Mar 24 04:20:41 access-node hostapd[975]: ens38: RADIUS Authentication server 10.0.30.10:1
        812
Mar 24 04:20:41 access-node hostapd[975]: ens38: Deauthenticate all stations at BSS start
Mar 24 04:20:41 access-node hostapd[975]: ens38: interface state UNINITIALIZED->ENABLED
Mar 24 04:20:41 access-node hostapd[975]: ens38: AP-ENABLED
Mar 24 04:20:41 access-node hostapd[975]: ens38: Setup of interface done.
Mar 24 04:20:41 access-node hostapd[975]: ctrl_iface not configured!
\end{verbatim}

% ---------------------------------------------------------------------
% SUBSECTION B.6: PACKET CAPTURE METADATA AND FIGURE MAPPING
% ---------------------------------------------------------------------
% This subsection preserves the pcap inventory, capture commands,
% Wireshark figure mapping, and the data redaction policy.
% ---------------------------------------------------------------------

\subsection{B.6 Packet Capture Metadata and Figure Mapping}
\label{sec:appendix-b-packet-capture-metadata-and-figure-mapping}

% ---------------------------------------------------------------------
% SUBSECTION B.6.1: PCAP FILE INVENTORY
% ---------------------------------------------------------------------
% This subsection preserves the authoritative pcap inventory listing
% referenced in the packet-analysis section.
% ---------------------------------------------------------------------

\subsubsection{B.6.1 pcap File Inventory}
\label{sec:appendix-b-pcap-file-inventory}

\begin{verbatim}
access-node: ls -lh ~/pcaps
-rw-r--r-- 1 nac-lab nac-lab 385 Mar 23 20:23 failed_8021x_ens38.pcap
-rw-r--r-- 1 nac-lab nac-lab 738 Mar 23 20:23 failed_8021x_ens39.pcap
-rw-r--r-- 1 nac-lab nac-lab 385 Mar 23 20:29 success_8021x_ens38.pcap
-rw-r--r-- 1 nac-lab nac-lab 747 Mar 23 20:29 success_8021x_ens39.pcap
\end{verbatim}

% ---------------------------------------------------------------------
% SUBSECTION B.6.2: PACKET CAPTURE METHODOLOGY AND TCPDUMP COMMANDS
% ---------------------------------------------------------------------
% This subsection preserves the authoritative packet-capture commands
% used to produce the failed and successful authentication pcaps.
% ---------------------------------------------------------------------

\subsubsection{B.6.2 Packet Capture Methodology and tcpdump Commands}
\label{sec:appendix-b-packet-capture-methodology-and-tcpdump-commands}

\begin{verbatim}
access-node:
sudo timeout 20s tcpdump -i ens38 -w /tmp/failed_8021x_ens38.pcap ether proto 0x888e &
sudo timeout 20s tcpdump -i ens39 -w /tmp/failed_8021x_ens39.pcap host 10.0.30.10 and 
        udp port 1812 &
sleep 2

client-1:
sudo systemctl stop wpa_supplicant
sudo timeout 15s wpa_supplicant -i ens33 -D wired -c 
        /etc/wpa_supplicant/lab/wired-8021x.conf -dd

access-node:
sudo timeout 20s tcpdump -i ens38 -w ~/pcaps/success_8021x_ens38.pcap ether proto 0x888e &
sudo timeout 20s tcpdump -i ens39 -w ~/pcaps/success_8021x_ens39.pcap host 10.0.30.10 and
        udp port 1812 &
sleep 2

client-1:
sudo systemctl stop wpa_supplicant
sudo timeout 15s wpa_supplicant -i ens33 -D wired -c 
        /etc/wpa_supplicant/lab/wired-8021x.conf -dd
\end{verbatim}

% ---------------------------------------------------------------------
% SUBSECTION B.6.3: PACKET CAPTURE REDACTION POLICY
% ---------------------------------------------------------------------
% This subsection preserves the publication-safe redaction
% policy for Wireshark packet screenshots.
% ---------------------------------------------------------------------

\subsubsection{B.6.4 Packet Capture Redaction Policy}
\label{sec:appendix-b-packet-capture-redaction-policy}

\begin{verbatim}
Publication policy:
- private IP addresses: keep visible
- MAC addresses: redact
- passwords/shared secrets/auth blobs: redact
- packet payload details that expose secrets: redact

Accordingly, the publication-safe Wireshark images should redact:
- MAC addresses in the packet list pane
- MAC addresses in the packet details pane
- MAC addresses in the hex pane
- RADIUS Authenticator values where shown
- Message-Authenticator values where shown
\end{verbatim}

\bigskip

\noindent\textit{End of document.}

\end{document}
